%% 溯源.tex —— 「题目溯源」报告（47 份口径）
%% **本文件由 工具/README转tex.py 从 ../README.md 的「题目溯源」一节自动生成，不要手改**：
%% 改内容请改 README，改排版请改那个脚本，然后 `make trace` 重生成。
%% 编译：xelatex 溯源.tex（跑两遍，第二遍才有目录与交叉引用）
\documentclass[a4paper,11pt]{article}
\usepackage{common}
\usepackage{booktabs}
\usepackage{xltabular}          % longtable + 固定列宽（宽表见 README 的「题目溯源」）
\usepackage{colortbl}           % 宽表交替底纹（\rowcolors）：行高 3—4 行的表不串行
\usepackage{fancyhdr}
\usepackage{lastpage}           % 页脚「第 N 页 / 共 M 页」
\usepackage{hyperref}
\hypersetup{colorlinks=true,linkcolor=solblue,urlcolor=solblue,
            bookmarksopen=true,bookmarksopenlevel=1,pdfstartview=FitH,
            pdftitle={题目溯源},pdfsubject={上海高一上数学 · 跨卷同题与真题出处考证}}
\setlength{\parskip}{0.28em}
\setlength{\parindent}{0pt}
\renewcommand{\arraystretch}{1.12}
\definecolor{rowgray}{gray}{0.955}   % 交替底纹：比白纸略深一点，不抢正文

%% pandoc 的列表里会插 \tightlist，它只在 pandoc 自带模板里定义，独立编译要自己补，
%% 否则报 `Undefined control sequence \tightlist`（踩过）。定义与 pandoc 模板同文。
\providecommand{\tightlist}{%
  \setlength{\itemsep}{0pt}\setlength{\parskip}{0pt}}

\pagestyle{fancy}
\fancyhf{}
%% 右眉显示**当前节名**：翻到中间页时不至于不知道在哪一节（以前右眉是死的口径标签）。
\fancyhead[L]{\small\color{solblue} 题目溯源}
\fancyhead[R]{\small\nouppercase{\rightmark}}
\fancyfoot[L]{\small 47 份口径（2026-10-02）}
\fancyfoot[C]{\small 第 \thepage 页 / 共 \pageref{LastPage} 页}
\renewcommand{\headrulewidth}{0.4pt}
\renewcommand{\footrulewidth}{0.4pt}

\begin{document}

\begin{titlepage}
\centering
\vspace*{2.6cm}
{\Huge\bfseries 题目溯源}\\[0.9em]
{\Large 上海高一上数学 · 跨卷同题与真题出处考证}\\[2.6em]
{\large 47 份试卷 824 道一级题干 $\times$ 参考库 16{,}273 道高考真题}\\[0.35em]
{\large 跨卷两两 339{,}076 对 \quad|\quad 逐字同题 48 组、同模板变体 20 组}\\[3.2em]
\begin{minipage}{0.8\textwidth}\small\raggedright
本文件由 \texttt{工具/README转tex.py} 从 \texttt{README.md} 的“题目溯源”一节自动生成，
内容与该节逐字一致。判定口径、阈值与踩过的坑见末节“方法与局限”；
两档同题表的每一行都注明了该对“变在哪”，跨卷同题的成员已逐题回 \texttt{原图/} 核过。
\end{minipage}
\end{titlepage}

%% 目录本身也进书签（以前 9 个书签只有 9 节，目录不在里面）
\pdfbookmark[1]{\contentsname}{toc}
\tableofcontents
\newpage
\clearpage
\section{核对摘要}

\newcommand{\tabhdrA}{\rowcolor{white} 项 & 结果 \\}
\rowcolors{1}{white}{rowgray}
\begin{xltabular}{\textwidth}{@{}>{\raggedright\arraybackslash}p{6.55cm}>{\raggedright\arraybackslash}p{8.61cm}@{}}
\toprule
\tabhdrA
\midrule
\endfirsthead
\multicolumn{2}{r}{\footnotesize（续上表）} \\
\midrule
\tabhdrA
\midrule
\endhead
\bottomrule
\endlastfoot
核对范围 & \textbf{47 份试卷的 824 道}一级题干（其中最早入库的 12 份 211 道），逐题比对参考库的 \textbf{16,273 道}高考真题；跨卷两两 \textbf{339,076 对}。\textbf{高一16／高一18 不计入}------它们分别与 高一b／高一c 同卷，重复计入只会造出成对的假同题 \\ \addlinespace[1pt]
定到具体高考试卷 & \textbf{43 道真题}（\textbf{50 条对应}：原题照录 22、改编 25、借定义／借集合另出题 3），对应本库 \textbf{58 个题位}------其中 \textbf{3 道}真题被\textbf{三份}试卷用到（2015 上海卷（春）：高一13／17／25；2003 上海卷（秋理）：高一17／20／36；2010 湖南卷（文）填空 7：高一08／高一a／高一39），\textbf{9 道}各被两份用到（含 2012 江苏 \(f(n)\) 题：高一02／高一42；2019 上海卷（春）填空：高一30／高一45），\textbf{31 道}单独一份使用。三个数与「真题出处」表\textbf{逐行数出来的完全相同}，拆解与该表的出处一并见该节；\textbf{2026-09-28 起这几处数与那张表已由 \texttt{工具/\allowbreak{}计数核对.py} 逐行复算守住} \\ \addlinespace[1pt]
跨卷同题（校际重复） & \textbf{逐字同题 48 组 + 同模板变体 20 组}（另有 8 对分数够不上变体线、但读过原文确属真同题，手工补记），两档成员共 \textbf{130 道}（逐字 96 ＋ 只出现在变体档的 34）------逐字档成员均经逐题对照 \texttt{原图/\allowbreak{}} 核实；高一47 新增 2 组逐字同题与 2 组同模板变体（其中 Q10 与高一07 第12题按完整题目由机判的逐字候选重判为变体） \\ \addlinespace[1pt]
与真题同题型但非同一题 & 约 40 道（上一轮逐条判读；本轮只跑机判，未逐条重判），未计入 \\ \addlinespace[1pt]
其余 & 真题库里既无 $\geqslant$12 字的逐字串、也无「\(r\ge0.71\) 且最长串 $\geqslant$ 12」的候选 $\to$ 非高考真题 \\ \addlinespace[1pt]
\end{xltabular}

\begin{quote}
下面几段是\textbf{上表数字的来历}（溯源一共跑过十五轮：12 份 $\to$ 17 份 $\to$ 18 份 $\to$ 21 份 $\to$ 22 份 $\to$
24 份 $\to$ 28 份 $\to$ 29 份 $\to$ 30 份 $\to$ 36 份 $\to$ 40 份 $\to$ 41 份 $\to$ 43 份 $\to$ 46 份 $\to$ \textbf{47 份}）；前十一轮的数字保留原文以便对照，\textbf{以末一段的 47 份口径为准}
第三轮是 \textbf{2026-09-15 全量重跑}：17 份一起重算（此前只跑过前 12 份，即
高一02---高一09 与 高一a---高一c），新增的五份（高一12／13／14／15／17）由此带出
\textbf{8 个题位}的真题出处与 \textbf{5 组}跨卷同题，另补出旧卷漏检
\textbf{3 条}（高一09 第 5 题、高一05 第 19 题、高一a 附加题第 2 题），合计 11 个题位。
\textbf{2026-09-15 已补跑高一10}（原 17 份 $\to$ 18 份、301 $\to$ 320 道、45,150 $\to$ 51,040 对）：本表、下面
的同题表与真题出处表都已含它。\textbf{高一10 的 19 题在参考库里没有一道达同题线}------最好的匹配只有
11--15 字的通用片段（如 \texttt{集合M=1,2,3,4,5,6}、\texttt{已知集合A=x\textbar{}x2-\allowbreak{}2x-\allowbreak{}3}），整段相似度都 $\leqslant$ 0.60；
但它与库内两份卷比出两组\textbf{逐字同题}（见同题表末两行）。\textbf{高一16／高一18 不计入}：它们分别与
高一b／高一c 同卷，重复计入只会造出成对的假同题。
四份卷的 \textbf{10 道附加题}另见「附加题专查」一节（其中一道可确指校际试卷出处、
一道连网络检索也落空）；另有若干\textbf{非高考但可考据到来源}的题，单列一节。
2026-09-15 的第二轮复核又用「\textbf{数字脱敏}」与「\textbf{自定义名词}」两路重扫，补出
高一02 第 15 题、高一09 第 15 题、高一17 第 11 题三条（都是判据捞不到的形态）。
参考库题干数 16,273（此前记 16,176／16,180）是同一套 \texttt{content/\allowbreak{}**/\allowbreak{}*.tex} 在不同去噪阈值下的读数。
\textbf{2026-09-15 末又把高一19---高一21 三份并入，做了 21 份的全量重跑}：原 18 份 320 道 $\to$ \textbf{21 份
369 道}、跨卷 51,040 $\to$ \textbf{67,896 对}。这一轮的变动分三块------
① \textbf{三份新卷直接带出}的：逐字同题 4 对（02\#6$\times$21\#10「伙伴关系集合」、06\#5$\times$20\#4
「\(ax^2+x+1=0\) 只有一个元素」、07\#9$\times$19\#11「四个数字元素的集合」、09\#13$\times$19\#8「六个关系式」）、
同模板变体 2 组（17\#11$\times$20\#16「数域」、17\#7$\times$20\#15「条件判断」）、新出处 1 条
（高一20 第 15 题＝2003 上海卷（秋理））；
② \textbf{改判据后翻出来的旧卷同题}：同模板变体 2 组（13\#18$\times$15 附加 19「\(x=m^2-n^2\)」、
06\#7$\times$17\#9「填包含关系」）、逐字同题 2 对（07\#10$\times$a\#13「特征数」、15\#14$\times$a\#14，
它们原先都在「手工补记」里）、把 09 第 19 题并入 07／08／11 那一组（3 份 $\to$ \textbf{4 份}）；
③ \textbf{改判与剔除}：05 第 16 题从逐字档移到变体档（05 版漏「非空」、答案 16 与 07／11 的 15
不同），08 第 15 题 $\times$ 19 第 7 题\textbf{回原图发现图形不同}、从两档剔除。
\textbf{判据本身也修了两处}（见「方法与局限」），故两档的组数、成员数与「真题出处」表都随本轮重排。
\textbf{2026-09-16 又把高一22（七宝中学 抱孩子练习9.15）并入，做了 22 份的全量重跑}：原 21 份 369 道 $\to$
\textbf{22 份 387 道}、跨卷 67,896 $\to$ \textbf{74,691 对}。这一轮\textbf{只增不改}，变动全部来自高一22 本身------
① 新出处 \textbf{3 条}：高一22 第 11 题＝\textbf{2011 年大纲卷（文/理）}、第 13 题＝\textbf{2010 年辽宁卷（理）}、
附加 19 题＝\textbf{2006 年湖北卷（理）}（三条都是\textbf{原题照录}，详见「真题出处」表末三行）；
② 该份在\textbf{参考库里再无别的达线命中}（其余 15 题全部「无命中」）；
③ \textbf{跨卷同题与变体一档零新增}------逐字同题仍是 24 组、同模板变体仍是 12 组、两档成员仍是 73 道，
高一22 与库内 21 份\textbf{一组都没比出来}（同校的 高一15 抱孩子练习9.8 与它也不例外）。
三条新出处使「真题」由 24 道增到 \textbf{27 道}、对应关系由 26 条增到 \textbf{29 条}、题位由 32 个增到 \textbf{35 个}，
其中\textbf{原题照录}由 9 条增到 12 条（改编 15 条、借定义／借集合 2 条均不变）。
\textbf{2026-09-20 又把高一23（格致 周末作业3）与高一24（复旦附中 周末作业3）并入，做了 24 份的全量重跑}：
原 22 份 387 道 $\to$ \textbf{24 份 419 道}、跨卷 74,691 $\to$ \textbf{87,571 对}。变动------
① 新出处 \textbf{1 条}：高一23 第 12 题＝\textbf{2004 年北京卷（秋文/秋理）}（\textbf{改编}：机判先落在变体档，
回原文复核确认同源------借真题题干与 A、B 两个选项，把问法由「一定成立的是」反转为「不成立的是」，
C、D 两项各改一处，答案由 A 变 B，详见「真题出处」表末行）；
② 高一24 的 15 道在真题库里\textbf{全部无命中}，高一23 其余 16 道亦无达线命中；
③ \textbf{跨卷同题与变体一档零新增}------逐字同题仍是 24 组、同模板变体仍是 12 组、两档成员仍是 73 道，
两份新卷与库内 22 份\textbf{一组都没比出来}。
新出处使「真题」由 27 道增到 \textbf{28 道}、对应由 29 条增到 \textbf{30 条}、题位由 35 个增到 \textbf{36 个}
（\textbf{改编} 15 $\to$ 16，照录 12、借定义 2 均不变）。
\textbf{2026-09-22 又把高一25---高一28 四份（华二 · 周末作业3、交附闵行 · 周练2、交附闵行 · 周末作业9.18、
七宝 · 周末作业3）并入，做了 28 份的全量重跑}：原 24 份 419 道 $\to$ \textbf{28 份 488 道}、跨卷
87,571 $\to$ \textbf{118,828 对}。变动------
① \textbf{新真题 1 道}：高一28 第 16 题＝\textbf{2013 年浙江卷（文）}（\(\wedge\)／\(\vee\) 两个自定义运算，
\textbf{原题照录}：定义式、四个选项与答案 C 逐字相同，\(r=0.966\)）；
② \textbf{旧真题多出 1 个题位}：高一25 第 16 题＝\textbf{2015 年上海卷（春）}（\(P_1\)／\(P_2\)／\(Q_1\)／\(Q_2\)
四集合版，\textbf{原题照录}------与高一17 第 16 题逐字同题，但高一17 那版 C、D 两项把真题的「存在 \(a\)」
印成了「对任意 \(a\)」（原卷笔误），高一25 这版与真题一致）；
③ 四份新卷在真题库里\textbf{其余各题全部无命中}------机判另给出 4 条候选（高一26 第 7 题 vs 2019 江苏卷、
高一26 第 8 题 vs 1999 上海卷（文）、高一28 第 5 题 vs 2007 北京卷（理）、高一28 第 18 题 vs
2007 上海卷（春）），\textbf{逐条回原文复核后都判为假命中}：共享的只是
「\(B=\{x\mid x>0\}\)」「若 \(A\subseteq B\) 求 \(a\) 的取值范围」「\(x^2-5x+4\geqslant0\)」
「\(>0\) 的解集为 \((-\infty,\)\ldots」这类通用片段，题干的集合与问法都不相同；
④ \textbf{跨卷新增 2 组逐字同题}（高一12 第 9 题 $\times$ 高一26 第 11 题------\textbf{同一所学校}前后两周作业
用了同一道题；高一17 第 16 题 $\times$ 高一25 第 16 题，即上面的 2015 上海春考）与 \textbf{2 组同模板变体}
（高一04 第 15 题 $\times$ 高一26 第 8 题、高一15 第 6 题 $\times$ 高一25 第 6 题）；
另有一对（高一13 第 14 题 $\times$ 高一25 第 16 题）落在已有的上海春考那一组里，只把该组由 2 份扩成 3 份。
以上使「真题」由 28 道增到 \textbf{29 道}、对应由 30 条增到 \textbf{32 条}、题位由 36 个增到 \textbf{38 个}
（\textbf{原题照录} 12 $\to$ 14，改编 16、借定义 2 均不变）；逐字同题 24 $\to$ \textbf{26 组}、同模板变体 12 $\to$ \textbf{14 组}。

\textbf{2026-09-23 又把高一29（延安中学 · 周末作业3）并入，做了 29 份的全量重跑}：原 28 份 488 道 $\to$
\textbf{29 份 506 道}、跨卷 118,828 $\to$ \textbf{127,765 对}。这一轮\textbf{只增不改}，且新增的这一份
\textbf{一处达线的命中都没有}------
① \textbf{真题零新增}：高一29 的 18 题在真题库里全部无命中（最好的匹配只有十几字的通用片段），
故「真题」仍是 29 道／32 条对应／38 个题位；
② \textbf{跨卷同题零新增}：逐字同题仍是 26 组、同模板变体仍是 14 组、两档成员仍是 79 道，
高一29 与库内各卷\textbf{一组都没比出来}；
③ 唯一值得记一笔的是\textbf{高一29 第 19 题与高一15 第 14 题、高一a 第 14 题共用一个含参集合}
（\(A=\{x\mid x^2-ax+a^2-19=0\}\)）------但另两卷的 \(B\)、\(C\) 与三小问都与它不同，
题干相似度未到变体线（\(r<0.71\)），故\textbf{不列入同题表}，只在此备查；
这也再次说明各校备课组用的是\textbf{同一批公共题源}（同一个含参方程在四份卷里出现三次）。
\end{quote}

\begin{quote}
\textbf{2026-09-25 又把高一30（南洋模范中学 · 开学考）并入，做了 30 份的全量重跑}：原 29 份 506 道 $\to$
\textbf{30 份 525 道}、跨卷 127,765 $\to$ \textbf{137,550 对}。变动------
① \textbf{新真题 2 道，两条都是原题照录}：高一30 第 4 题＝\textbf{2010 年安徽卷（秋文）填空}
（命题「存在 \(x\in\R\)，使得 \(x^2+2x+5=0\)」的否定，\(r=0.909\)）、
高一30 第 12 题＝\textbf{2019 年上海卷（春）填空}（\(A=[t,t+1]\cup[t+4,t+9]\)、\(0\notin A\)、
存在正数 \(\lambda\) 使 \(\dfrac\lambda a\in A\)；\textbf{整题一字不差}，\(r=1.000\)、最长公共子串 47 字）；
② 机判另给出 1 条候选（高一30 第 19 题 vs 2017 北京卷（文）），回原文复核后判为\textbf{假命中}------
真题是 \(A=\{x\mid x<-2\text{ 或 }x>2\}\) 求 \(\complement_U A\)，本卷是 \(x\geqslant3\) 且三小问全不同，
共享的只有「集合 \(A=\{x\mid x<-2\) 或 \(x\)\ldots」这 12 个字；
③ \textbf{跨卷新增 4 组逐字同题}（高一11 第 20 题 $\times$ 高一30 第 19 题、高一15 第 6 题 $\times$ 高一30 第 8 题、
高一22 第 5 题 $\times$ 高一30 第 9 题、高一25 第 12 题 $\times$ 高一30 第 15 题），
另有一对（高一25 第 6 题 $\times$ 高一30 第 8 题）落在已有的变体组里、把该组由 2 份扩成 3 份。
以上使「真题」由 29 道增到 \textbf{31 道}、对应由 32 条增到 \textbf{34 条}、题位由 38 个增到 \textbf{40 个}
（\textbf{原题照录} 14 $\to$ 16，改编 16、借定义 2 均不变）；逐字同题 26 $\to$ \textbf{30 组}，同模板变体仍是 \textbf{14 组}。

\textbf{2026-09-26 又把高一31---高一36 六份并入，做了 36 份的全量重跑}：原 30 份 525 道 $\to$
\textbf{36 份 626 道}、跨卷 137,550 $\to$ \textbf{195,625 对}。这六份是 9-25 一天内连发的六篇，
产出比前几轮都多------
① \textbf{新真题 5 道}（2009 天津卷（理）、2008 江苏卷、2012 浙江卷（理）、2004 湖北卷（文）、
2010 福建卷（文）），共添 \textbf{7 条对应}：其中 5 条原题照录、2 条改编（高一35 附加 17 题把
2009 天津卷那道由选择改填空；高一36 第 3 题把 2004 湖北卷的 \(x\leqslant6\) 改成 \(5\)）；
② \textbf{另有 1 条落在已有真题上}：高一36 第 13 题是 2003 上海卷（秋理）的\textbf{第三个题位}
（前两个是 高一17 第 7 题、高一20 第 15 题）；
③ \textbf{机判另给出 2 条候选，复核后判假命中}------高一36 第 9 题 vs 2014 浙江卷（文）
（真题讲的是平方和与平均值，与 \(\dfrac ca\) 的范围无关）；高一31 第 14 题与高一35 附加 17 题
虽同出 2009 天津卷（理），但一为选择、一为填空，判据把它们算作两条；
④ \textbf{跨卷新增 7 组逐字同题}，另有一对（高一08 第 15 题 $\times$ 高一32 第 13 题）把既有的一组
由 2 份扩成 3 份------这一组正是「题干一字不差但图形可能不同」的那类，已单独看图确认
（高一19 第 7 题图形不同，不计）。
以上使「真题」由 31 道增到 \textbf{36 道}、对应由 34 条增到 \textbf{41 条}、题位由 40 个增到 \textbf{45 个}
（\textbf{原题照录} 16 $\to$ 21、改编 16 $\to$ 18，借定义 2 不变）；逐字同题 30 $\to$ \textbf{37 组}，同模板变体仍是 \textbf{14 组}。

\textbf{2026-09-27 又并入高一37---高一40 四份（连载 37---40）}：36 份 626 道 $\to$ \textbf{40 份 701 道}、
跨卷 195,625 $\to$ \textbf{245,350 对}。这一轮------
① \textbf{新真题 3 道}（2009 重庆卷（理）、2010 四川卷（文）、2005 北京卷（春季·文/理）），
共添 \textbf{3 条对应}：其中 \textbf{1 条原题照录}（高一37 第 11 题＝2010 四川卷（文），题干
\(a>b>0\) 与 \(a^2+\frac1{ab}+\frac1{a(a-b)}\) 一字不差、答案同为 \(4\)，只去掉四个选项改填空）、
\textbf{2 条改编}（高一37 第 8 题把 2009 重庆卷（理）的右侧 \(a^2-3a\) 换成 \(\lvert a-3\rvert\)、选择改填空，
答案随之由「\(a\leqslant-1\) 或 \(a\geqslant4\)」变为「\(a\leqslant-1\) 或 \(a\geqslant7\)」；
高一37 第 19 题把 2005 北京卷（春季）第 (1) 问的「精确到 \(0.1\) 千辆/小时」改成
「保留分数形式」，其余分子 \(920\)、分母 \(v^2+3v+1600\) 与第 (2) 问全同）；
② \textbf{机判另给出 1 条候选，复核后判假命中}：高一40 第 9 题 vs 2012 福建卷（文）------
真题是 \(x^2-ax+2a>0\) 在 \(\R\) 上恒成立（\(0<a<8\)），本卷是 \(x^2-ax-2<0\) 在 \((1,2)\) 内
恒成立（\([1,+\infty)\)），除「恒成立，则实数 \(a\) 的取值范围是」这类套话外全不同。
另：高一37 第 8 题机判报出的公共串也是同类套话，但\textbf{同源证据在左侧的 \(\lvert x+3\rvert-\lvert x-1\rvert\)}
（与重庆卷一字未动），故该条仍判\textbf{改编}而非假命中；\textbf{高一38 全卷 17 题零命中}；
③ \textbf{跨卷新增 6 组逐字同题}，另有 2 对把既有组各由 2 份扩成 3 份------新增的 6 组里有 4 组是
\textbf{高一39 与家长拍照卷 高一a 的重题}（第 7／10／附加 2 题各一，加上充要条件与区间那组），
另 2 组是 高一38 与 高一03／07 的（见同题表新增各行）；扩为 3 份的两组是「和谐集」
（＋高一38 第 13 题，38 版另起名「平衡集」）与「特征数」（＋高一39 第 16 题）；
\textbf{高一37 与任何既有卷都不构成同题}（19 题一条都没比出来）。
以上使「真题」由 36 道增到 \textbf{39 道}、对应由 41 条增到 \textbf{44 条}、题位由 45 个增到 \textbf{48 个}
（\textbf{原题照录} 21 $\to$ 22、改编 18 $\to$ 20，借定义 2 不变）；逐字同题 37 $\to$ \textbf{43 组}，同模板变体仍是 \textbf{14 组}。

\textbf{2026-09-28 又并入高一41（上海中学 · 周练2，连载 41）}：40 份 701 道 $\to$ \textbf{41 份 718 道}、
跨卷 245,350 $\to$ \textbf{257,403 对}。这一轮\textbf{一处达线的命中都没有}，是「只增不改」的又一例
（同 2026-09-20 的高一23／24、2026-09-23 的高一29）------
① \textbf{真题零新增}：高一41 的 17 题在真题库里全部无命中（最接近的只有「恒成立，则实数 \(a\) 的
取值范围是」这类套话），故「真题」仍是 \textbf{41 道／46 条对应／54 个题位}（另有两条是同日下一段补登的）
------见下面「2026-09-28 又\ldots」那一段；
② \textbf{跨卷同题零新增}：逐字同题仍是 43 组、同模板变体仍是 14 组、两档成员仍是 110 道，
高一41 与库内各卷\textbf{一组都没比出来}（同校的 高一31 周练1 与它也不例外）；
③ \textbf{补登两个此前漏登的题位}------这一轮顺核「逐字同题组的成员是否都进了真题出处表」才发现：
\textbf{高一39 附加 21 题}（与 高一a 附加 2 题逐字同题）＝ 2010 湖南卷（文）填空 7、
\textbf{高一32 第 13 题}（与 高一08 第 15 题逐字同题、\textbf{图形也一样}）＝ 1999 广东卷／全国卷选择 1。
两者都不改道数／条数，只把题位由 50 补到 \textbf{52}、把 2010 湖南卷（文）从「两份用到」提为
\textbf{三份用到}、1999 广东卷／全国卷从「一份」提为\textbf{两份}。这两个题位当时机判都报过
（高一39 附加 21 题那条更是明写 \texttt{出处=2010\ 年湖南卷（文）｜填空题}），只是没人登记进表。
\end{quote}

\begin{quote}
\textbf{2026-09-28 同日又把①机判清单接进闸}（起因是核「高一35／高一39 两份附加题在专记里是否齐全」）：
先把 \texttt{题目溯源.py} ①「与高考库比对」的命中\textbf{落盘}成 \texttt{工具/\allowbreak{}真题命中清单.txt}（718 个题位一行），
再把\textbf{每个报过的题位}都摊开回原文判一遍------结果\textbf{新查出 2 道真题／2 条对应}：
① \textbf{高一20 第 8 题 ＝ 2023 年全国乙卷（理）选择}（共用 \(U=\R\)、\(M=\{x\mid x<1\}\)、
\(N=\{x\mid-1<x<2\}\) 这一对集合与「用交、并、补表示某个集合」的问法，只换所求集合，
\(r=0.727\)、LCS 22），判\textbf{借同一对集合另问}；
② \textbf{高一21 第 4 题 ＝ 2009 年上海卷（秋文）填空}（真题 \(A=\{x\mid x\leqslant1\}\)、
\(B=\{x\mid x\geqslant a\}\)、\(A\cup B=\R\) 求 \(a\)，本卷把端点与字母对调，答案由 \(a\leqslant1\)
变 \(a\geqslant1\)；\(r=0.912\)、LCS 16），判\textbf{改编}。
其余 \textbf{17 条判为不算}（只共用通用句式或某一个集合的取法），逐条记进新建的
「\textbf{① 机判报了、但未计入本表的题位}」表------以后每轮跑完 ① 若多出一个报过的题位，
\texttt{计数核对.py} 就会报错，逼人在这张表或出处表里 disposition，不再有「报了没人看」的余地。
以上使「真题」由 39 道增到 \textbf{41 道}、对应由 44 条增到 \textbf{46 条}、题位由 52 个增到 \textbf{54 个}
（\textbf{改编} 20 $\to$ 21、\textbf{借定义／借集合} 2 $\to$ 3，原题照录 22 不变）；逐字同题 43 组与变体 14 组不变。
\end{quote}

本轮核对\textbf{发现并订正了 11 处}（都是核对过程本身查出来的，不是新发现而是纠错）：

\begin{enumerate}
\def\labelenumi{\arabic{enumi}.}
\tightlist
\item
  \textbf{转录错误}：高一08 第 15 题（1999 广东卷 Venn 图题）选项 A 误录为 \((M\cap P)\setminus S\)，
  原卷是 \((M\cap P)\cap S\) ------ 已改源码并重编两版。
\item
  \textbf{判定写错}：高一a 第 11 题原判``改编（有序对$\to$差集）''，但真题本身就是差集
  \(B=\{x-y\mid x,y\in A\}\)，应判\textbf{原题照录}（差别只在四个选项的次序被倒过来）。
\item
  \textbf{描述不实}：高一a 第 2 题原写「\(U\) 少一个元素，问法改变」，实际问法未变，只是
  \(U\) 由 \(\{0,\dots,9\}\) 改成 \(\{0,\dots,8\}\)（答案随之由 \(\{7,9\}\) 变为 \(\{7\}\)）。
\item
  \textbf{题号与计数}：三处题号（高一03、高一09、高一c）先前按 \texttt{\textbackslash{}item} 出现次序编号，
  与卷面的 \texttt{\textbackslash{}setcounter\{enumi\}} 口径不符，已订正；「一级题干 230 道」「真题 15,932 道」
  两个数字同属口径错误，已改为 \textbf{211} / \textbf{16,180}。
\item
  \textbf{跨卷同题的判据与计数}：变体档阈值原写 0.74，却与表内最低那对（「全食／偏食」
  \(r=0.7129\)）自相矛盾，已改为 \textbf{0.71}；变体表原多出一行份数 1 的成员，使「7 组」
  与 8 行对不上，已删行、改由正文交代；又补上 4 对分数低于变体线、但逐条读过确属
  真同题的配对（见「判据之外手工补记」）。
\item
  \textbf{变体表里有 6 行的描述与成员原文对不上}（复核成员时逐条比对原始 \texttt{.tex} 才发现）：
  ① 08 第 17 题的 \(P\) 实为\textbf{区间} \((-2,3)\)，本行先写 \(\{2,3\}\)、后照源码的错改成
  \(\{-2,3\}\)（源码也错，见第 8 条）；② 09 第 17 题的 \(P\) 是
  「不等式 \(x^2-x-6<0\) 的解集」，原写成「方程 \(x^2-x-6=0\) 的解集」；③ 03 第 16 题的
  \(A\) 是 \(\{x\mid a\leqslant x\leqslant -a+3\}\)，原写成 \(\{x\mid a<x<a+3\}\)；④ 05 第 7 题的
  \(A\) 是 \(\{x\mid x\geqslant-1\}\)，原行标题却按 c 第 5 题的 \(\{x\mid 1<x<2\}\) 写（且
  c 第 5 题实为 \(1\leqslant x\leqslant 2\)）；⑤ 02／05 第 17 题原写「小问不同」，实则只
  (1) 问不同，(2) 问两版\textbf{完全相同}（同为 \(A\cup B=A\)、答案同为 \((-\infty,-3]\)）；
  ⑥ 三行的「变在哪」漏了条件，已补出实际的 \(\cup\)、\(\cap\)、\(\varnothing\)、\(\subset\)。
  六行的「题目」「变在哪」两栏已按原文重写。
\item
  \textbf{高一05 第 17 题两处转录错误}（对着原卷图片逐字比对时才发现）：① (1) 问原卷作
  「\(B\) 为\textbf{单元素}集合」，本稿误抄成「\(B\) 为单元集集合」；② 答案 (2) 原卷作
  \((-\infty,-3]\)，本稿误抄成 \((-\infty,-3)\)------\textbf{本稿自己的解析写的就是 \((-\infty,-3]\)
  （\(a=-3\) 时 \(B=\{2\}\subseteq A\) 也成立），即解析对、\texttt{\textbackslash{}ans} 错}。已改源码并重编两版。
\item
  \textbf{高一08 第 17 题的 \(P\) 抄成了 \(\{-2,3\}\)}：原卷印的是\textbf{开区间} \((-2,3)\)。源码自己的
  解析就是按开区间做的（\(a+1<3\)、答案 \(\left[-\tfrac{2}{3},+\infty\right)\)），是题干抄错。
  已改源码并重编两版。
\item
  \textbf{高一06 填空题漏了起始偏移}：原卷填空从\textbf{第 3 题}编到第 11 题（1、2 未印在本卷），
  本稿漏写 \texttt{\textbackslash{}setcounter\{enumi\}\{2\}}，于是印成 1--9，而解答题又按原卷印成 12--14，卷面
  凭空多出「10、11 缺号」。同批 8 份（02／03／04／05／07／08／09／11）都有这一行，只
  高一06 漏了。已补上并重编两版------受此影响，本 README 里 高一06 的题号整体 \textbf{+2}
  （孤立元 7$\to$9、全食／偏食 9$\to$11、\(B=\{x\mid ax-1=0\}\) 那题的 1$\to$3）。
\item
  \textbf{高一05 第 10 题的答案端点抄错}：原卷解析得 \(a<-4\) 或 \(2<a\)，本稿的 \texttt{\textbackslash{}ans} 与
  「综上」都写成 \(a<-4\) 或 \(a\geqslant 2\)------\(a=2\) 时 \(B=[4,5]\not\subseteq A=(-\infty,-1)\cup(4,\infty)\)，
  故原卷为是。已改源码并重编两版。
\item
  \textbf{高一07 第 10 题的「特征数」抄错}：原卷作「最大值与最小值\textbf{之和}」，本稿抄成
  「\textbf{之差}」------原卷解析「\(A_1=\{1,2,3,12\}\)，则 \(A_1=13\)」（\(=12+1\)，原卷此处把 \(X_1\) 误印
  成 \(A_1\)，本稿按 \(X_1\) 规范并加注）可证为「之和」。
  这个抄错一度让本 README 误判该对「差／和不同」（见第 6 条附近的手工补记）。已改重编。
\end{enumerate}

以上 11 条是\textbf{溯源核对}那一轮（12 份）的记录。2026-09-15 又对新入库的五份
（高一12／13／14／15／17）做了一轮\textbf{整卷回原图}核对（41 张图逐张放大比对），
另查出 \textbf{5 处转录错误}，清单见文末「9 月中旬五份专记」里的回原图复核表。

\begin{quote}
\textbf{2026-10-02 并入高一42、43并做全库重跑}：41 份 718 道 $\to$ \textbf{43 份 756 道}，跨卷 257,403 $\to$ \textbf{285,390 对}。
① 高一42 第 20 题与 \textbf{2012 江苏卷 \(f(n)\) 计数题}同源（不另增真题道数），但卷面条件③与源题的相对补集表述不同，且本卷按字面会与自身 \(f(4)=4\) 的答案冲突；按\textbf{改编}登记，并在源码与「真题出处」表留下注记。
② 高一43 第 18 题由 \textbf{1997 年全国卷（文/理）} 的高速运输成本模型改编，加入每小时费用限制问，按一题一条对应登记（文、理两版是同一道题）。
由此真题出处由 41 道／46 条／54 个题位增至 \textbf{42 道／48 条／56 个题位}（原题照录 22、改编 23、借定义／借集合 3）；2012 江苏卷成为第 8 道被两份本库卷采用的真题。
③ 跨卷复核确认高一26 第 8 题与高一43 第 7 题逐字同题，新增一组；高一39 附加22 与高一42 第11题、高一37 第20题与高一43 第20题、高一c 第2题与高一43 第3题新增三组同模板变体，并将高一43 第7题补入既有高一04／高一26 变体组。两条只共用参数不等式句式的候选（高一12 第14题$\times$高一42 第7题、高一33 第17题$\times$高一42 第19题）判为非同源，不入表。该轮更新到 \textbf{45 组逐字同题、17 组同模板变体}。

\textbf{2026-10-02 再并入高一44---46并全量重跑}：43 份 756 道 $\to$ \textbf{46 份 806 道}，跨卷 285,390 $\to$ \textbf{324,415 对}。
① 高一45 第 10 题＝\textbf{2019 年上海卷（春）填空}的改编（LCS 22、\(r=0.832\)）：保留``正区间集合在 \(a\mapsto\lambda/a\) 下不变''的设定，更换区间端点并把 \(0\notin A\) 改为 \(t>0\)。该题源已由高一30 第12题登记，故只增加一条对应与一个题位、不增加真题道数。
② 高一44、45、46 在真题库均无新增可登记题；高一45 第10题是本批唯一需要对高考机判命中作出处登记的题位。
③ 跨卷复核：高一44 第7题并入高一12 第9题／高一26 第11题逐字同题组；高一45 第14题与高一08 第14题逐字同题，新增一组；高一45 第10题与高一30 第12题属同模板变体。高一46 第7题仅共享含参集合问法，不入表。该轮更新到 \textbf{46 组逐字同题、18 组同模板变体}，两档共122个成员位置（逐字92、仅变体30），另有8对手工补记。真题出处更新为 \textbf{42 道／49 条／57 个题位}（原题照录22、改编24、借定义／借集合3），2019上海春考现有高一30与高一45两个对应题位。

\textbf{2026-10-02 并入高一47并全库重跑}：46 份 806 道 $\to$ \textbf{47 份 824 道}，跨卷 324,415 $\to$ \textbf{339,076 对}。
① 高一47 第 6 题是 \textbf{2012 年江苏卷填空}的一道新改编：两题都给二次函数值域 \([0,+\infty)\)，再由 \(f(x)<c\) 的解集为区间求 \(c\)；本卷将区间长度由 6 改为 8，答案相应为 16。它不同于表中高一02／高一42 所据的同卷 \(f(n)\) 计数题，故新增一道真题及一条改编对应。
② 高一47 第 11 题的两条机判候选（2012 浙江卷（理）、2007 安徽卷（文））只共享集合 \(B\) 的二次式；\(A\) 的定义及所求集合不同，判为通用集合片段，不登记真题出处，列入下方「未计入」台账。
③ 跨卷候选逐题复核后，高一47 第 5 题与高一27 第 6 题、附加 1 题与高一13 第 10 题为两组逐字同题；第 4 题与高一c 第 4 题、第 10 题与高一07 第 12 题为两组同模板变体（后一对完整题目不同，虽落在机判逐字候选中仍归变体）。全库更新为 \textbf{48 组逐字同题、20 组同模板变体}，两档共130个成员位置（逐字96、仅变体34），另有8对手工补记。真题出处更新为 \textbf{43 道／50 条／58 个题位}（原题照录22、改编25、借定义／借集合3）。
\end{quote}

\clearpage
\section{回原图逐题核对（2026-09-15 起，按批推进）}

\textbf{页级核对（已全部完成）}：把公众号原文的图片逐张下载，与 \texttt{原图/\allowbreak{}} 里归档的文件比对------
有原文链接的现为 \textbf{44 份 / 376 张}（高一02---高一09 与 高一10---高一15、高一17，加高一19---高一47），其中最早一轮
逐张核过的是 \textbf{13 份共 95 张}，\textbf{张数、顺序、逐张 MD5 全部一致}；其余三十一份不在那一轮内
（高一12 交附闵行 11 张、高一10 复附浦东 9 张、高一19 延安 6 张、高一20 进才 9 张、
高一21 上师大宝山 6 张、高一22 七宝 7 张、高一23 格致 5 张、高一24 复旦附中 7 张、
高一25 华二 10 张、高一26 交附闵行 10 张、高一27 交附闵行 7 张、高一28 七宝 11 张、
高一29 延安 6 张、高一30 南模 11 张、高一31 上海中学 10 张、高一32 华二 10 张、
高一33 交附闵行 13 张、高一34 交附闵行 7 张、高一35 七宝 12 张、高一36 格致 6 张、
高一37 七宝 10 张、高一38 大同 12 张、高一39 大同 11 张、高一40 大同 4 张、
高一41 上海中学 15 张、高一42 复旦附中 8 张、高一43 交附闵行 12 张、高一44 交附闵行 8 张、
高一45 交附闵行 10 张、高一46 交附闵行 9 张、高一47 七宝 9 张------这些均为按页序保存的微信公众号原图;
\textbf{高一37---高一47 十一份中，除高一40 四张都为 4762$\times$6735 外，其余十份均有「同一篇里各页分辨率不一致」的情形}
（高一37：第 3、4、6 页 4762$\times$6735，其余 7 页 2560$\times$3620；高一38：第 1、10、11、12 页
2560$\times$3620，第 2---9 页 4762$\times$6735；高一39 亦混排；\textbf{高一41：第 2、3、4、6、8、11、12、13、14 页
4762$\times$6735，第 1、5、7、9、10、15 页 2560$\times$3620}；高一42 第 02、08 页 2560$\times$3620、其余六页 4762$\times$6735；
高一43 第 01、05、09、12 页 2560$\times$3620、其余八页 4762$\times$6735；高一44 第 08 页 2560$\times$3620、其余七页 4762$\times$6735；
高一45 第 07、08、10 页 2560$\times$3620、其余七页 4762$\times$6735；高一46 第 03---06、08---09 页 2560$\times$3620、其余三页 4762$\times$6735；高一47 第 08 页 2560$\times$3620、其余八页 4762$\times$6735）；
高一22 的 7 张另有两点与前几份不同：一是\textbf{同一篇里各页分辨率不一致}------第 1、2、4、6 页
4762$\times$6735，第 3、5、7 页 2560$\times$3620（已逐页打开确认，非下载缺档）；二是这一份取的是
\texttt{mmbiz.qpic.cn/\allowbreak{}\textless{}hash\textgreater{}/\allowbreak{}0} 的\textbf{原图}（4762$\times$6735 那一档，与 高一20／高一21 同档），
比早期入库的几份清晰------高一02、高一15 那一批的 \texttt{原图/\allowbreak{}} 是 1080$\times$1527 的窄幅档）。
\textbf{高一25---高一28 这四份又回到 1080$\times$1527 那一档}（四篇各页同尺寸，与高一23／24 同档）。
高一a---高一c、高一16 与 高一18 是家长拍摄的作业照片（后两者的原图是 \texttt{高一b}／\texttt{高一c} 同一组照片的副本），
不适用这一路。这一步只能证明「纸面都搬进来了」，逐题逐字
是否转录完整另见下表。

\textbf{逐题核对}（每卷都逐页读原图，比对题干、选项、\texttt{\textbackslash{}ans} 与 \texttt{\textbackslash{}ifanswer} 解析全文）。\textbf{截至
2026-10-02，除 \texttt{高一16}／\texttt{高一18} 两份「对号副本」外，47 套卷子都过了一遍独立于转写的
回原图逐题核对}------工作分几批推进：下表前 12 行（高一02---高一09、高一a---高一c 与 高一11，分四批）；
另五份（高一12／13／14／15／17）的\textbf{整卷回原图}核对（41 张图逐张放大比对）在时间上更早，
清单见文末「9 月中旬五份专记」里的回原图复核表；高一25---高一28 于 2026-09-22 独立做过一轮、
高一29 于 2026-09-23、高一30---高一36 于 2026-09-26、高一37---高一40 于 2026-09-26（另一名
worker 只读原图）、高一41 于 2026-09-28、高一42／高一43 与高一44---高一47 于 2026-10-02 各做过一轮（后几轮见下表相应各行）。

\begin{quote}
\textbf{2026-09-26 这一轮补的是最后一批「只有转写、没人复核过」的卷子}：高一10 与 高一19---高一24
七份入库时都是「转录时逐页回对」，此后一直没有独立核对（下表里原先写着「暂无独立于转写的
第二轮回原图复核记录」）。这一轮给七份各派一名 worker、\textbf{只读原图、不给它看既有结论}，
并额外要求\textbf{逐题独立验算答案}。结果：\textbf{七份全部查出问题}------
\end{quote}

\newcommand{\tabhdrB}{\rowcolor{white} 卷 & 这一轮的发现 & 处理 \\}
\rowcolors{1}{white}{rowgray}
\begin{xltabular}{\textwidth}{@{}>{\raggedright\arraybackslash}p{1.98cm}>{\raggedright\arraybackslash}p{7.78cm}>{\centering\arraybackslash}p{4.98cm}@{}}
\toprule
\tabhdrB
\midrule
\endfirsthead
\multicolumn{3}{r}{\footnotesize（续上表）} \\
\midrule
\tabhdrB
\midrule
\endhead
\bottomrule
\endlastfoot
高一10 & 第 10 题【解析】原卷「这个\textbf{元}不是\ldots」漏「素」字（本稿当初已补字，但文件头误称「未作订正」） & 保留补字、登记进文件头（注释） \\ \addlinespace[1pt]
高一19 & \textbf{1 处真转录错误}：第 12 题【解析】末句原卷是 \(0\in\overline A\)，本稿漏抄上横线（且文件头误把它记成原卷问题） & 订正并重编 \\ \addlinespace[1pt]
高一20 & 第 21(3)【解析】原卷作 \(\overline{B}\cap\allowbreak  A\)、题干作 \(\overline{A}\cap\allowbreak  B\)，本稿当初统一成题干的写法 & 恢复原卷次序并重编 \\ \addlinespace[1pt]
高一21 & 第 8 题②分支原卷作「当 \(B\neq\varnothing\) 时，\textbf{则} \ldots」，本稿漏「则」 & 补字并重编 \\ \addlinespace[1pt]
高一22 & 无转录错误；四处「原卷存疑」逐处核实吻合，另\textbf{新查出第五处}原卷问题（第 14(2) 分析与答案的区间开闭矛盾） & 加注、存疑改五处（注释） \\ \addlinespace[1pt]
高一23 & 无转录错误；\textbf{新查出三处原卷自身问题}（取等条件、第 11 题答案、第 15(2) 答案漏负号） & 加注（注释） \\ \addlinespace[1pt]
高一24 & 第 14 题题干原卷把「\(=\allowbreak 0\)」印在花括号外、其解析却印在括号内 & 保留修正并登记（注释） \\ \addlinespace[1pt]
\end{xltabular}

\begin{quote}
七份的\textbf{答案}（共 123 道：19＋18＋20＋16＋18＋17＋15）都逐一验算过、无一处与解析矛盾；
\texttt{高一19} 的 3 张韦恩图与 \texttt{高一24} 的三圆图也逐张与原卷比过（后者用二值化 mask 比出 IoU$\approx$0.91），
另四份（高一10／20／21／22）经核对确认\textbf{全卷无插图}。
\end{quote}

\newcommand{\tabhdrC}{\rowcolor{white} 卷 & 原图/题数 & 核对结果 \\}
\rowcolors{1}{white}{rowgray}
\begin{xltabular}{\textwidth}{@{}>{\raggedright\arraybackslash}p{3.39cm}>{\centering\arraybackslash}p{3.39cm}>{\raggedright\arraybackslash}p{7.96cm}@{}}
\toprule
\tabhdrC
\midrule
\endfirsthead
\multicolumn{3}{r}{\footnotesize（续上表）} \\
\midrule
\tabhdrC
\midrule
\endhead
\bottomrule
\endlastfoot
高一02 & 8 张 / 17 题 & \textbf{无转录错误}。唯一差异是第 18 题【解析】原卷把第 (2) 问印成「（3）」（原卷跳号），本稿按题干口径规范为 (2)，已加注 \\ \addlinespace[1pt]
高一03 & 6 张 / 16 题 & \textbf{3 处解析转录错误}，均已订正重编：① 第 10 题【解析】多写一对「\(S_3\) 和 \(S_5\) 不能同时取数」（原卷只列 \(S_1S_6\)／\(S_2S_5\)／\(S_3S_4\) 三对，且 \(3+\allowbreak 5=\allowbreak 8\equiv1\pmod 7\)，本可同时取）；② 同题末句原卷作「和 \(S_0\) 中的一个」，本稿误作「\(S_6\)」；③ 第 18(3)【解析】等号左端首项原卷作 \(\dfrac{1}{x-1}\)，本稿误作 \(\dfrac{1}{x}\) \\ \addlinespace[1pt]
高一04 & 3 张 / 14 题 & \textbf{2 处实质错误 + 2 处记号差异}，均已订正重编：① 第 3 题题干原卷作「\(y=\allowbreak \dfrac{12}{3-x}\in\mathbb N\)，\(x\in\mathbb Z\)」，本稿把 \(\in\mathbb N\) 挂到了 \(x\) 上（那样只能得 \(\{4,\allowbreak 6,\allowbreak 12\}\)，与本稿自己的答案冲突）；② 第 7 题【答案】原卷作「\(A\subset\allowbreak  B\)」，本稿抄成「\(A\supsetneq B\)」------\textbf{方向反了}；③ 第 6 题【答案】、④ 第 16 题(2) 的 \(\subset\allowbreak \) 本稿原写作 \(\subsetneq\allowbreak \)／\(\subseteq\allowbreak \)，已按原卷记号改回 \\ \addlinespace[1pt]
高一05 & 9 张 / 18 题 & \textbf{16 处}（题干/选项/答案 3、解析 13），均已订正重编。题干/选项/答案：① 第 19 题题干性质 \(M\) 中 \(x_ix_j\) 与 \(\dfrac{x_j}{x_i}\) 的分式被抄成 \(\dfrac{x_i}{x_j}\)------\textbf{方向反了}；② 第 20 题记号 \(N_+\allowbreak \) 写成 \(\mathbb N^*\)；③ \textbf{第 15 题的选项 B、选项 D 与【解析】末句三处把「开左闭右」的 \(\left(\frac43,\allowbreak \frac53\right]\) 抄成 \(\left[\frac43,\allowbreak \frac53\right)\)（开闭抄反）}------与它自己第 ① 步解出的「\(\frac43<m\leqslant\allowbreak \frac53\)」自相矛盾（2026-09-15 复核题号时补出）。解析（集中在第 4、5、9、11、19、20 题）：第 4 题末项丢负号（\(-\frac52\) 抄成 \(\frac52\)）；第 5 题把原卷「\textbf{两个方程组用『或』并列}」抄成一个三层嵌套的 \texttt{cases}；第 9 题把原卷「集合 \(\overline B\) 中含 1、5」改写成「集合 \(B\) 不含 1、5」（两句等价，按原卷录回）；第 11 题两处「\textbf{一元}一次方程」抄成「二元一次方程」；\textbf{第 19 题解析 5 处}------下标 \(x_1\)／\(x_2\) 抄成 \(i\)／\(j\)、\(x_1\geqslant\allowbreak  1\) 抄成 \(x_2\geqslant\allowbreak  1\)、\(x_i\cdot x_5>x_5\) 抄成 \(>x_i\)、\(\dfrac{x_5}{x_i}\) 抄成 \(\dfrac{x_5}{x_1}\)、\(x_3x_4>x_2x_4=x_5\) 不等号写反且漏一句（末两行不等式链整个走样）；\textbf{第 20 题解析 5 处}------\(0\pm 0\) 抄成 \(0+\allowbreak 0\)、不等链末项 \(a_3-a_1\) 抄成 \(a_2-a_1\) 且 \(a_3-a_2=a_2\) 抄成 \(a_3-a_3=a_2\)（后者会使 \(a_2=a_3=\allowbreak 2026\)，与本卷答案 \(a_2=\allowbreak 1013\) 冲突）、三处下标 \(n-i\) 抄成 \(n-1\)、展开式首项 \((a_n-a_n)\) 抄成 \((a_n-a_1)\) \\ \addlinespace[1pt]
高一06 & 4 张 / 12 题 & \textbf{10 处}（题干/答案 3、记号 1、题号口径 1、解析 5），均已订正重编：① 第 3 题【答案】漏 \(a=\allowbreak 1\)（原卷作「0, 1, -1」，\(a=\allowbreak 1\) 时 \(B=\allowbreak \{1\}\subseteq\allowbreak  A\)）；② 第 4 题题干把 \(\in\mathbb Z\) 挂到了 \(x\) 上（原卷作 \(A=\allowbreak \{x\in N\mid\allowbreak  y=\allowbreak \dfrac{12}{x+\allowbreak 3}\in Z\}\)，即「\(y\) 是整数」）；③ 第 7 题题干与【答案】的 \(\subset\allowbreak \)／\(\supset\) 照原卷改回（本稿原写作 \(\subsetneq\allowbreak \)／\(\supsetneq\)），其解析方向也是反的（本稿作「\(T\subseteq\allowbreak  S\)，故 \(S\supsetneq T\)」，与本题答案 \(\subset\allowbreak \) 自相矛盾）；④ 文末「参考答案」块原从 1 起编，与卷面第 3---14 题错位，已补 \texttt{\textbackslash{}setcounter\{enumi\}\{2\}}；⑤ 解析：第 11 题 \(B=\allowbreak \{-\frac12,\allowbreak \frac12\}\) 丢负号（原卷此句还把 \(A\cap\allowbreak  B\neq\varnothing\) 印成 \(=\allowbreak \varnothing\)，已加注）、第 13 题题干漏「其和」且解析首段被压缩成一句（已按原卷补回）、第 14 题两小问漏了原卷的「（给出结论即可）」「（需说明理由）」、其 (2) 解析末句「相异正整数」抄成「正非零整数」 \\ \addlinespace[1pt]
高一07 & 11 张 / 19 题 & \textbf{13 处}（按源码里的 \texttt{\%} 注计：漏排整图 1、原卷笔误加注 2、转录错误 10；其中第 16 题一处含 4 个错点，合并为 3 条注），均已订正重编：① 第 14 题解析原卷印有一张韦恩图（\(A\) 独 5、\(B\) 独 5、三集合公共 10、\(A\cap\allowbreak  C\) 独 4、\(B\cap\allowbreak  C\) 独 4、\(C\) 独 2、三集合之外 10），本稿此前\textbf{漏排此图而正文仍作「如图」}，已按原卷截图补入（\texttt{figures/\allowbreak{}venn\_50\_ABC.png}）；② 第 16 题：差集记号 \(A\setminus\allowbreak  B\) 抄成 \(A\vee B\)、\(C\) 描述里的待填横线抄成 \(\vee\vee\)，选项 D 的 \(\complement_{\mathbb N}\) 抄成 \(\varnothing\)、\(\complement_{\mathbb N}C\) 抄成 \(C\setminus\allowbreak  C\)，同题解析里两处 \(\complement_{\mathbb N}\) 抄成 \(C\setminus\allowbreak \)；③ 第 12 题解析「得 \(x_2=\allowbreak -\dfrac{3a}{2}\)」漏掉负号（本稿写成 \(\dfrac{3a}{2}\)）；④ 第 20 题【答案】漏 \(\frac34\)（与本稿自己的解析结论不一致）；⑤ 第 11 题② \(c\notin M\) 抄成 \(c\in M\)（那样②为假，与本卷答案 ①②③ 冲突）、第 19 题定义 \(2-x^2\notin D\) 抄成 \(\in D\)（那样 \(E=\allowbreak \varnothing\)，与答案 \(-5\) 冲突）、第 20 题解析 \(xy=\allowbreak 0\in D\) 抄成 \(\notin D\)、第 21 题解析上界 \(x_4-x_3<x_4\) 抄成 \(<x_3\)，其第 (3) 问不等式链抄乱且首项 \(a_1-a_1\) 抄成 \(a_i-a_1\)；⑥ 第 10 题原卷自身把 \(X_1=\allowbreak 13\) 印成「\(A_1=\allowbreak 13\)」、第 18 题原卷子集列举印成 \(\{0-4\}\)（两处均为原卷笔误，按文意规范并加注） \\ \addlinespace[1pt]
高一08 & 9 张 / 19 题 & \textbf{1 处漏排整图 + 5 处转录错误}（源码内 5 条 \texttt{\%} 注，其中第 11 题一条含 3 个错点、另有 1 条是分页说明），均已订正重编：① 第 15 题原卷印有一个韦恩图（全集 \(I\)、\(M\) 与 \(P\) 两圆相交、\(S\) 为含于 \(P\) 的小圆、阴影为 \((M\cap\allowbreak  P)\setminus\allowbreak  S\)），本稿\textbf{漏排此图而题干仍作「如图」}，已按原卷截图补入（\texttt{figures/\allowbreak{}venn\_MP\_minus\_S.png}）；② 第 6 题解析原卷作「所以 \(3\in C\)，\(2\notin C\)，\(-4\notin C\)」，本稿把 \(2\notin C\) 抄成 \(2\in C\)------\textbf{与题干 \(A\cap\allowbreak  C=\allowbreak \varnothing\) 及 \(2\in A\) 直接矛盾}；③ 第 11 题题干原卷作「集合 \(E=\allowbreak \{0,\allowbreak 1,\allowbreak 2,\allowbreak \cdots,n\}\) 的子集 \(\{a_1,a_2,\allowbreak \dots,a_m\}\)」（原卷上标用 \(n\)、下标用 \(m\)，本身不统一），本稿把上标也写成 \(m\)；④ 同题解析原卷作「\(211=\allowbreak 2^7+\allowbreak 2^6+\allowbreak 2^4+\allowbreak 2^1+\allowbreak 2^0\)」（本稿写成 \(2^2\)，那样和为 213）；⑤ 同题【答案】原卷作 \(\{0,\allowbreak 1,\allowbreak 4,\allowbreak 6,\allowbreak 7\}\)，本稿漏掉 7（作 \(\{0,\allowbreak 1,\allowbreak 4,\allowbreak 6\}\)，仅 83）；⑥ 第 17(1) 题解析原卷的方程组含三行「\(3a\geqslant\allowbreak  -2\)／\(a+\allowbreak 1<3\)／\(3a<a+\allowbreak 1\)」，本稿略去了末行 \\ \addlinespace[1pt]
高一09 & 8 张 / 20 题 & \textbf{1 处漏排整图 + 2 处转录错误}，均已订正重编：① 第 16 题原卷印有一个韦恩图（全集 \(U\)、相交的椭圆 \(A\) 与 \(B\)、阴影为 \(A\cup\allowbreak  B\) 之外），本稿\textbf{漏排此图而题干仍作「如图」}，已按原卷截图补入（\texttt{figures/\allowbreak{}venn\_U\_minus\_AB.png}）；② 第 9 题解析 ② 原卷作「但 \(-1+\allowbreak \sqrt3\notin\{x\mid\allowbreak  x=t^2,\allowbreak \,t\in M\}\)」，本稿把 \(\notin\) 抄成 \(\in\)------那样②反而与 \(M\) 相等，与同句「即②与集合 \(M\) 不相等」自相矛盾；③ 第 12 题题干原卷作「\textbf{高一的}同学发现\ldots\ldots」（逐字放大后「高一」与「的」之间只有常规字距，无班号），本稿写成「高一（13）班的同学发现」，已按原卷删去「（13）班」 \\ \addlinespace[1pt]
高一a & 4 张 / 17 题 & \textbf{无转录错误}。这一份是家长拍摄的学生作业照片，逐页比对后题干、选项、分值行、附加题编号与卷面笔色分工（黑笔学生作答／蓝笔教师订正／红笔打钩$\times$）全部与 \texttt{\%} 注记录一致；卷面 10 处学生错答与教师的蓝笔订正也逐处核对无误（清单见文末「高一a 专记」） \\ \addlinespace[1pt]
高一b & 4 张 / 21 题 & \textbf{无转录错误}。三处已知差异（第 11 题的答案取 \(5\) 而非学生所写的 \(10\)、第 13 题原卷漏排「\(A=\allowbreak \)」、第 16---18 题的作答区）与卷首两处排版差异（原卷未印分节标题与分值、本稿补「2026-\/-2027 学年度」前缀与校名「大同中学」）均已在 \texttt{\%} 注与「高一b 专记」写明，本轮逐页放大复核无误；第 4 题的韦恩图与 \texttt{figures/\allowbreak{}venn\_AB\_complement.png} 一致（阴影在 \(A\cup\allowbreak  B\) 之外，与【答案】\(\overline{A\cup\allowbreak  B}\) 相符） \\ \addlinespace[1pt]
高一c & 4 张 / 19 题 & \textbf{无转录错误}；只有 1 处非题面标签：第 17 题题号后原卷印有一个虚线框小标签「手批」（全卷仅此一处，疑为标注该题由教师手批），本稿不排，已在卷首 \texttt{\%} 注与源码内写明。第 5 题的 \(\subset\allowbreak \)（无下划线，即真包含，答案 \((2,\allowbreak +\allowbreak \infty)\)）与同卷第 9、13 题的 \(\subseteq\allowbreak \) 本轮再逐处放大比对，确认两者在卷面上确有区别；第 1、2、6、9、10、11、12 题与选择 13---16 的题干、选项、答案亦与卷面逐字一致 \\ \addlinespace[1pt]
高一11 & 7 张 / 19 题 & \textbf{无转录错误}。四处原卷笔误（第 10 题【解析】①③ 的「为真命题」、第 19(1) 的 \(\{0-4\}\)、第 19(2) 的 \(a^2-1=\allowbreak -4\times\allowbreak (-4)\)）与三处排版口径（答案直接印在题下、\textbf{第 14 题的答案「D」更是印在题干括注里}、实数集 \(R\) 与 \(\mathbf{R}\) 混用）均已在卷首 \texttt{\%} 注与「高一11 专记」写明，本轮逐页放大复核无误；第 21(3) 的 \(\subset\allowbreak \) 与第 10②、12、17、19、20 题的 \(\subseteq\allowbreak \) 也逐处放大确认；第 6 题的韦恩图与 \texttt{figures/\allowbreak{}venn\_A\_minus\_B.png} 一致（阴影为 \(A\setminus\allowbreak  B\)，与【解析】\(A\cap\allowbreak \overline B\) 相符） \\ \addlinespace[1pt]
高一10 & 9 张 / 19 题 & \textbf{新入库（本次转写）}。填空 3---12、选择 13---16、解答 17---21，题干、选项、\texttt{\textbackslash{}ans} 与 \texttt{\textbackslash{}ifanswer} 解析逐页照原卷录入（沿用同一学校的 \texttt{高一11} 的体例）；转写时逐页放大比对过一遍，\textbf{未发现原卷笔误}------题 16 的 \(1995=\allowbreak 15\times\allowbreak 133\) 与「\(15^2\) 的倍数有 8 个」两处最易被误判为笔误，实测 \(133\)、\(8\) 均正确。\textbf{第二轮}（2026-09-26）由另一名 worker 独立核过一遍：19 题答案逐一验算一致；查出 \textbf{1 处}------第 10 题【解析】首句原卷作「这个\textbf{元}不是 \(1\) 就是 \(m\)」（「元素」漏了「素」字），本稿当初补了「素」字却在文件头写着「未发现笔误、未作任何按文意的订正」，自相矛盾；已按 高一b 第 13 题「原卷漏排、本稿补上」的先例保留补字，并把该处登记进文件头的笔误清单、加行内注（注释改动，未重编 PDF）。全卷 9 页无图 \\ \addlinespace[1pt]
高一16 & 4 张 / 21 题 & \textbf{与 \texttt{高一b} 同卷}（源码逐字相同），故不另做一轮核对------上表 \texttt{高一b} 那一行的结论直接适用于本卷。两卷只差卷面标题的「（补）」二字（已用 \texttt{pdftotext} 逐行比对确认） \\ \addlinespace[1pt]
高一18 & 4 张 / 19 题 & \textbf{与 \texttt{高一c} 同卷}（源码逐字相同），同上不作独立核对。两卷只差卷面标题的「（补）」二字（已用 \texttt{pdftotext} 逐行比对确认） \\ \addlinespace[1pt]
高一19 & 6 张 / 18 题 & \textbf{新入库（本次转写）}。填空 3---13、解答 14---20，题干、\texttt{\textbackslash{}ans} 与 \texttt{\textbackslash{}ifanswer} 解析逐页照原卷录入；转写时逐页放大比对了题干、答案与解析全文，并单独核过三张韦恩图（第 4、6、7 题）的裁剪范围与阴影（第 4 题的图原卷印在【解析】右侧，故入答案版；第 6、7 题是题干图，两版都排------三张都取自原卷截图、4$\times$ LANCZOS 放大）。\textbf{原卷存疑三处}已照录加注：第 11 题「已知一个有四个数字元素的集合 \(S\)」、第 16(1) 结论「\(a\) 的值为 \(5\backslash-3\)」（「\textbackslash」应是顿号）、第 12 题【解析】末句「同理当 \(0\in A\) 时」（按文意应为 \(0\notin A\)，但卷面字形确为 \(\in\)、无斜杠）。\textbf{第二轮}（2026-09-26）由另一名 worker 独立核过一遍（含 3 张插图逐张比对、18 题答案逐一验算）：查出 \textbf{1 处真转录错误}------第 12 题【解析】末句原卷是 \(0\in\overline A\)（\(A\) 上\textbf{有横杠}），本稿漏抄了 \(\overline{}\)；更麻烦的是文件头那条「原卷存疑③」把\textbf{本稿的漏抄}误记成了原卷问题（还写着「卷面字形确为 \(\in\)、无斜杠」）。已按原卷补回上横线、删去该条存疑（清单改回两处）并重编 \\ \addlinespace[1pt]
高一20 & 9 张 / 20 题 & \textbf{新入库（本次转写）}。填空 3---12、选择 13---16、解答 17---21，逐页照原卷录入。\textbf{订正原卷笔误两处}：① 第 13 题【解析】印作「所以 AD 错误，B 正确」「故选 B」，与同句「\(a\) 为奇数、\(b\) 为偶数 \(\Rightarrow a+b\) 为奇数」自相矛盾（\(a+b\) 为奇数就该落在奇数集 \(A\) 里），按文意订正为「BD 错误，A 正确」「故选 A」；② 第 20(5) 的【解析】原卷误标为「(2)」（同卷 (2) 的解答另有其文），按小问编号规范为 (5)。另第 14 题题干印「则下列四个命题」而其后列了 ①---⑤ 五条，照录加注。本卷无插图。\textbf{第二轮}（2026-09-26）由另一名 worker 独立核过一遍：20 题答案逐一验算一致；\textbf{两处原卷笔误的订正}（第 13 题答案 AD$\to$BD、第 20(5) 误标 (2)）经它放大复核，确认「改得对、没改过头」；查出 \textbf{1 处}------第 21(3)【解析】原卷印 \(\overline{B}\cap\allowbreak  A=\allowbreak \varnothing\)，而题干的 (3) 问印 \(\overline{A}\cap\allowbreak  B=\allowbreak \varnothing\)（原卷自己次序不一致，两者因交集可交换而等价），本稿当初统一成了题干的写法；已恢复原卷解析的次序并加注、重编 \\ \addlinespace[1pt]
高一21 & 6 张 / 16 题 & \textbf{新入库（本次转写）}。填空 3---11、解答 12---14（末题原卷标「（选做题）」），逐页照原卷录入。第 11 题的【答案】原卷印在【解析】末句（「综上，以上正确的命题是②④」），本稿按全稿惯例另提为 \texttt{\textbackslash{}ans\{\}}；第 14(1)【解析】首句原卷连用两个「因为」，照录加注。第 9 题的 \(\subset\allowbreak \)（真包含，与解析「真子集个数 \(2^8-1=\allowbreak 255\)」自洽）与第 11④ 的 \(\subseteq\allowbreak \) 逐处放大确认确有区别。本卷无插图。\textbf{第二轮}（2026-09-26）由另一名 worker 独立核过一遍：16 题答案逐一验算一致；查出 \textbf{1 处漏字}------第 8 题②分支原卷作「②当 \(B\neq\varnothing\) 时，\textbf{则} \ldots」，本稿漏了「则」（同卷第 12 题同位置的「有」是保留的，故确属漏字、不是全稿统一省略），已补并重编。它另把 9 处 $\subset$／$\subseteq$ 与 cases 逐行核对过，并逐段数了原图行数 vs \texttt{\textbackslash{}step} 条数、未发现去重 \\ \addlinespace[1pt]
高一22 & 7 张 / 18 题 & \textbf{新入库（本次转写）}。填空 3---10、选择 11---13、解答 14---16、附加题 17---20，逐页照原卷录入。原图是 4762$\times$6735／2560$\times$3620 两档高清档，逐段裁开放大后比对了题干、选项、\texttt{\textbackslash{}ans} 与 \texttt{\textbackslash{}ifanswer} 解析全文。\textbf{原卷存疑四处}已照录加注：① 第 8 题【解析】方程组第三行印成 \(x_1+x_2=a>0\)（前两行已用过 \(x_1+x_2\)，按两根为正数的判据此处应为两根之积 \(x_1x_2\)）；② \textbf{第 15 题题干的 \(B\) 与原卷自己的答案、解析对不上}（题干 / 答案 / 解析三处互相矛盾，已把三处都 4$\times$ 放大逐字核对）------题干印 \(B=\allowbreak \{(x,y)\mid\allowbreak  x=n,y=n^2-n+a+\allowbreak 1\}\)，按此则 \(A\cap\allowbreak  B\) 要在 \(m=n\) 时满足 \(3m+\allowbreak 1=m^2-m+a+\allowbreak 1\)、即 \(m^2-4m+a=\allowbreak 0\)：(1) 取 \(a=\allowbreak 1\) 时\textbf{无整数解}、\(A\cap\allowbreak  B\) 应为 \(\varnothing\)，(2) 的正整数解只有 \(a=\allowbreak 3\) 或 \(4\)；而原卷 (1) 的答案作 \(\{(0,\allowbreak 1),\allowbreak (4,\allowbreak 13)\}\)、(2) 的解析把 \(B\) 改写成 \(y=a(x^2-x+\allowbreak 1)\) 后得 \(a=\allowbreak 1\) 或 \(4\)。\textbf{只有把题干那处读作 \(y=a(n^2-n+\allowbreak 1)\) 才能同时解释答案与解析}（此时 \(3k+\allowbreak 1=k^2-k+\allowbreak 1\Rightarrow k=\allowbreak 0,\allowbreak 4\)，且 \(a=\allowbreak \frac{3x+\allowbreak 1}{x^2-x+\allowbreak 1}\) 在 \(x=\allowbreak 0,\allowbreak 1,\allowbreak 4\) 上取 \(1,\allowbreak 4,\allowbreak 1\)）；若读作 \(y=n^2-n+a\)，则 (1) 合、(2) 变成 \(a=\allowbreak 3,\allowbreak 4\)；③ 第 16(3) 的【解析】误标为「（2）」；④ 第 18 题题干印「已知\textbf{实数} \(M=\allowbreak \dots\)」（应为「集合」）。本卷无插图（已全卷检索确认无「如图」）。\textbf{第二轮}（2026-09-26）由另一名 worker 独立核过一遍：18 题答案逐一验算一致；\textbf{四处「原卷存疑」逐处放大核实，全部与原稿记载相符}（第 15 题那条「唯一能同时解释答案与解析的读法是 \(y=a(n^2-n+\allowbreak 1)\)」也被它独立复算证实），源码是忠实照抄；另\textbf{新查出第五处原卷问题}------第 14(2)【解析】的分支条件是\textbf{严格}的（\(3m-1>-m\) 即 \(m>\dfrac14\) 时 \(B\sub A\) 成立，\(m\leqslant\allowbreak \dfrac14\) 时无解），末行答案却印成\textbf{闭}区间 \(\left[\dfrac14,\allowbreak +\allowbreak \infty\right)\)，按它自己的分析应为开区间（已 3$\times$ 放大确认末行左侧确是方括号）。已加注、文件头的存疑清单由四处改为五处（注释改动，未重编 PDF） \\ \addlinespace[1pt]
高一23 & 5 张 / 17 题 & \textbf{第一轮（转写时逐页回对）之后一直没有独立核对，2026-09-26 才补上}：17 题的题干、选项、\texttt{\textbackslash{}ans}、\texttt{\textbackslash{}ifanswer} 逐句比对一字不差，各题答案也逐一验算通过（含第 12 题那道 2004 北京卷改编的答案 B）------\textbf{未发现转录错误}。但它查出的\textbf{三处原卷自身问题}推翻了该卷文件头原来那句「本卷未发现原卷笔误或存疑处」，已按「照录＋加注」补记：① 第 10 题【解析】的取等条件印作 \(x=\allowbreak 3y=\allowbreak 3z=\allowbreak \dfrac37\)，按 \(7M\geqslant\allowbreak  x+\allowbreak 3y+\allowbreak 3z\) 取等需 \(M=x=y=z\)（即 \(x=y=z=\allowbreak \dfrac37\)），终值 \(\dfrac37\) 无误、是书写自相矛盾；② 第 11 题原卷答案标 C，而取 \(0\)、\(1\)、\(-1\) 时三数之积为 \(0\)、\(0\)、\(-1\)、并不存在「两数之积为正数」，不正确的应是 \textbf{D}；③ 第 15(2) 答案印成 \(\left(\dfrac14,\allowbreak +\allowbreak \infty\right)\)，由 (1) 的 \(b=\allowbreak \dfrac78a\)、\(a>0\) 代入应为 \(\left(-\dfrac14,\allowbreak +\allowbreak \infty\right)\)（原卷漏了负号）。三处均已加注（注释改动，未重编 PDF） \\ \addlinespace[1pt]
高一24 & 7 张 / 15 题 & \textbf{第一轮（转写时逐页回对）之后一直没有独立核对，2026-09-26 才补上}：填空 3---10、选择 11---13、解答 14---17 逐句比对，\textbf{只查出 1 处}------第 14 题题干原卷把「\(=\allowbreak 0\)」印在右花括号\textbf{之外}（\(A=\allowbreak \{x\mid\allowbreak  x^2-mx+\allowbreak 4m-16\}=\allowbreak 0\)），而它\textbf{自己的【解析】却印在括号之内}（原卷前后不一致）；本稿当初跟了解析，按 高一b 第 13 题「原卷漏排 \(A=\allowbreak \)、本稿补上」的先例\textbf{保留修正}，并补了行内注与文件头「排版差异⑧」。唯一插图（第 11 题三圆图）用二值化 mask 与原卷比对，IoU$\approx$0.91、未镜像、\(A\) 左上 \(B\) 右上 \(C\) 正下，与答案 \(A\) 相符；全卷 7 页仅此一图。第 15(2) 的原卷笔误（\(+\allowbreak \frac ba\) 印成 \(-\frac ba\)）经放大核实确系原卷如此（注释改动，未重编 PDF） \\ \addlinespace[1pt]
高一25 & 10 张 / 19 题 & \textbf{无转录错误}（2026-09-22 独立于转录的第二轮逐题核对）。填空 3---12、选择 13---16、解答 17---21，题干、选项、\texttt{\textbackslash{}ans} 与 \texttt{\textbackslash{}ifanswer} 解析逐页照原卷录入。\textbf{原卷存疑两处}已照录加注：① 第 8 题【解析】「当 \(\Delta=a^2-16>0\) 时\ldots 有两个\textbf{相同}的根」（\(\Delta>0\) 应为「不相同」，其自己的后文就把 \(x_1\)、\(x_2\) 当不同的两根在讨论）；② 第 11 题【解析】先写 \(m\)、\(n\in\N^*\)，下一行又写「\(0\leqslant\allowbreak  m\leqslant\allowbreak 10\)，\(0\leqslant\allowbreak  n\leqslant\allowbreak 10\)，共有 121 种」（按 \(\N^*\) 应从 1 起，至多 100 种；121＝11$\times$11 含 0）。第 17 题的 Venn 图取自原卷截图 4$\times$ 放大（\texttt{figures/\allowbreak{}venn\_Abar\_cap\_B.png}，阴影为 \(\overline{A}\cap\allowbreak  B\)），原卷印在题干右侧、本稿移至题干之后 \\ \addlinespace[1pt]
高一26 & 10 张 / 19 题 & \textbf{3 处转录错误}（2026-09-22 第二轮核对查出），均已订正重编：① 第 13 题选项 C 的补集下标原卷作 \(\complement_U\)，本稿误抄成 \(\complement_I\)（放大 16$\times$ 与图中全集标号比对后确认是斜体 \(U\)）；② 第 20(2) 题【解析】原卷作真包含的 \(B\subset\allowbreak \complement_R A\)，本稿抄成带下划线的 \(\sub\)（同页 21(2) 的 \(\subseteq\allowbreak \) 多一条横杠，对比清晰）；③ 第 12 题【解析】三处原卷作 \(\dfrac85\leqslant\allowbreak  t\leqslant\allowbreak 2\)，本稿统一改成 \(t<2\)------改完与同题末行答案 \(\left(\dfrac95,\allowbreak 2\right]\) 自相矛盾（该答案含 \(2\)）。\textbf{原卷自身矛盾一处}已照录加注：第 12 题【解析】首步的方程组里原卷印 \(t<2\)，与它后面三处及末行答案的 \(t\leqslant\allowbreak 2\) 对不上（按题意 \(t=\allowbreak 2\) 时 \(N=\allowbreak \left[\dfrac75,\allowbreak 2\right]\) 确含于 \(\{x\mid\allowbreak 1\leqslant\allowbreak  x\leqslant\allowbreak 2\}\)，故应为 \(t\leqslant\allowbreak 2\)）。第 13 题的 Venn 图取自原卷截图 4$\times$ 放大（\texttt{figures/\allowbreak{}venn\_ABC\_lenses.png}，阴影为 \((A\cap\allowbreak  B)\cup\allowbreak (B\cap\allowbreak  C)\)），原卷印在题干右侧、本稿移至题干之后 \\ \addlinespace[1pt]
高一27 & 7 张 / 12 题 & \textbf{无转录错误}（2026-09-22 第二轮核对）。本卷只有「一、填空题」「二、解答题」两节、无选择题：填空 3---10、解答 11---14，逐页照原卷录入；\(\subset\allowbreak \)（第 7、10 题解析）与 \(\sub\)（第 11(2)、12(2) 题）逐处放大确认、\(\overline{A}\)／\(\complement_R\) 两种补集记号照录无误。核对时另发现\textbf{文件头两处注释列得不齐}，已订正（真包含的 \(\subset\allowbreak \) 除第 7 题解析外第 10 题解析也有两处；\(\complement_R\) 只在第 11 题出现、第 12 题全题无补集符号）------注释订正不影响正文，故未重编 PDF \\ \addlinespace[1pt]
高一28 & 11 张 / 19 题 & \textbf{无转录错误}（2026-09-22 第二轮核对）。填空 3---12、选择 13---16、解答 17---21。原卷第 3---6 题只印【答案】行、第 7 题起印【解析】，本稿统一改为试卷版隐去、答案版以 \texttt{\textbackslash{}ans\{\}}／\texttt{\textbackslash{}ifanswer} 给出。\textbf{原卷存疑一处}已照录加注：第 17(2) 题【解析】由 \(A\cup\allowbreak  B=\allowbreak \{x\mid\allowbreak  x>-2\}\)、\(A\cap\allowbreak  B=\allowbreak \{x\mid\allowbreak 1<x\leqslant\allowbreak 3\}\) 推出 \(B=\allowbreak [1,\allowbreak 3]\)、\(a=\allowbreak -4\)、\(b=\allowbreak 3\)------但 \(B=\allowbreak [1,\allowbreak 3]\) 时 \(A\cup\allowbreak  B=\allowbreak (-2,\allowbreak -1)\cup\allowbreak [1,\allowbreak +\allowbreak \infty)\) 并非 \(\{x\mid\allowbreak  x>-2\}\)；按题设应有 \(B=\allowbreak [-1,\allowbreak 3]\)（即 \(a=\allowbreak -2\)、\(b=\allowbreak -3\)），原卷答案显系错误。本卷无插图（全卷检索确认无「如图」） \\ \addlinespace[1pt]
高一29 & 6 张 / 18 题 & \textbf{无转录错误}（2026-09-23 独立于转写的第二轮回原图核对）。填空 3---12、选择 13---14、解答 15---20（末题是证明题），逐页照原卷录入；转写与复核两轮都逐页放大比对了题干、选项、\texttt{\textbackslash{}ans} 与 \texttt{\textbackslash{}ifanswer} 解析全文（原图第 2、3 页 4762$\times$6735，第 1、4、5、6 页 2560$\times$3620）。第二轮复核按全稿查出过错的五类逐处放大核对：全卷 \textbf{10 处多行 \texttt{cases}}（第 8、10、12、13 题解析、第 14 题的定义 2 处与解析 4 处、第 18(2) 题）、\textbf{补集记号}（第 14① 与第 16 题的 \(\overline{B}\)、第 14 题解析的 \(C_UB\)------下标放大到 16$\times$ 与题首全集 \(U\) 比对，两竖并列为斜体 \(U\)）、\textbf{$\subset$ 1 处与 $\subseteq$ 10 处}（第 4 题题干的 \(\subset\allowbreak \) 与 \(\sub\) 并列、第 13 题解析、第 14② 及其解析、第 17 题题干与解析）、以及\textbf{答案与解析是否自相矛盾}（第 6、13、14、16、17、19 题的 \texttt{\textbackslash{}ans} 与解析结论逐一对照），均与卷面一致。\textbf{原卷存疑三处}已照录加注：① 第 6 题【解析】\textbf{自相矛盾}------推出 \(m\leqslant\allowbreak 2\)，结论却写 \([2,\allowbreak +\allowbreak \infty)\)；② 第 14 题解析的补集记号\textbf{前后不一}（题干与解析首行作 \(\overline{B}\)，后文改作 \(C_UB\)）；③ 第 18 题题干作「必要\textbf{非}充分条件」、其解析作「必要\textbf{不}充分条件」（同义异写），另第 15 题解析「方程的\textbf{解}为 \(R\)」与同段后文「方程的\textbf{解集}为 \(R\)」措辞不一。本卷无插图（全卷检索确认无「如图」） \\ \addlinespace[1pt]
高一30 & 11 张 / 19 题 & \textbf{2 处标点级差异}（2026-09-25 独立于转录的第二轮核对），均已订正重编：① 第 12 题【解析】原卷作「\(\dfrac\lambda a\in[t,t+\allowbreak 1]\)．」（句末是实心句点，6$\times$ 放大确认），本稿误作逗号；② 第 21(3) 小题原卷作「\(\left\{n\mid\allowbreak \dfrac{495}{n+\allowbreak 2}\in\N,n\in\Z\right\}\) 求其所有子集\ldots」（集括号后\textbf{无}逗号），本稿多写了一个逗号。其余题干、选项、\texttt{\textbackslash{}ans} 与 \texttt{\textbackslash{}ifanswer} 解析全文逐题照原卷：第 7、19(2) 题带下划线的 \(\sub\) 与第 17(2) 题真包含的 \(\subset\allowbreak \) 逐处放大确认；第 13 题四个选项的 \(\in\)／\(\notin\) 逐个核对；第 17(2)、19(3) 的 \texttt{cases} 与第 18(2) 的三行分段函数、各处区间开闭均一致；\(k_1\)／\(k_2\)、\(A_i\)／\(A_j\)、\((-1)^m\)、\(2^9\)／\(2^n\)／\(2^{12}\) 等上下标无误；\texttt{\textbackslash{}ans} 与解析结论无矛盾、无漏项。本卷 \textbf{6 张插图}（第 13、14 题的韦恩图、第 16 题【解析】的三幅图、第 18 题的速度---时间图、第 20 题的题干图与【解析】图）逐一与原卷比对，画的都是原卷那个图形，原卷有图处源码都已排图。\textbf{原卷存疑三处}已照录加注：① 第 10 题【解析】把该集合记作 \(A\)（题设里叫 \(P\)）；② 第 21 题【解析】(3) 的「把这把这」重复；③ 同题 (2) 的结论式分子用花括号 \\ \addlinespace[1pt]
高一31 & 10 张 / 16 题 & \textbf{无转录错误}（2026-09-26 独立于转录的第二轮回原图核对）。填空 3---10、选择 11---14、解答 15---18，全卷无插图；\textbf{全卷没有单独的【答案】行}，每题下方直接印【解析】，答案写在解析末句（选择题作「故选 D／B／C」）。逐题核对：$\subseteq$（第 10 题条件①、第 11 题题干与选项 A、第 15(2) 题）与真包含（第 11 题选项 B、C）四处均放大确认；第 18 题题干两套 cases 的变化量分支顺序与「\(i=\allowbreak 1,\allowbreak 2,\allowbreak \cdots,n-1\)」的位置都与原卷一致；\(A_1\)／\(A_2\)、\(x_i\)／\(x_{i+\allowbreak 1}\)、\(2^{n-1}\)、\((x-2)^{999}\) 等指数下标无误；答案（第 11 题 D、第 12 题 B、第 13 题 B、第 14 题 C）与解析末句一致。\textbf{原卷存疑两处}已照录加注：① 第 11 题【解析】首句「因为 \(A=\allowbreak \{0,\allowbreak 1\}\) 是 \(A\) 的一个子集」（表述绕，结论 \(A\in B\) 成立）；② 第 18(3) 的「\(x\in\{7,\allowbreak 2,\allowbreak 3,\allowbreak 4\}\)」（按上下文应为 \(\{1,\allowbreak 2,\allowbreak 3,\allowbreak 4\}\)）。另：原卷并列各项一律用逗号或不加标点，本稿按全稿惯例改排顿号，已列入该卷「排版差异」⑨ \\ \addlinespace[1pt]
高一32 & 10 张 / 19 题 & 转录后已逐页回图核对（2026-09-26）。填空 3---12、选择 13---16、解答 17---21，只有第 13 题一图（韦恩图，取自原卷截图 4$\times$ 放大）。第 6、8、9 题只印【答案】、第 18 题的【答案】印作「略」、第 15 题的答案印在题干括号内，其余印【解析】。\textbf{原卷存疑四处}照录加注：① 第 19(2)【解析】「即 \(\dfrac12<m<\dfrac32\)」（按方程组与最终答案应为 \(\dfrac12<m<\dfrac23\)）；② 第 21(2)①「显然 \(\lvert X\rvert\geqslant\allowbreak 0\)」（疑为 \(\lvert S\rvert\)）；③ 第 20(3) 答案的 \(\{0,\allowbreak 4,\allowbreak 5,\allowbreak 6\}\) 是 4 元集而题问「3 元差集」；④ 第 21(2)② 的解析\textbf{原卷就未写完}（第 10 页止于 \(\lvert S-A\rvert=n-s\)、\(\lvert B-S\rvert=n-t\)）。\textbf{卷首年份}：卷面印作「\textbf{2025-2026 年}」，与连载「27学年高一32」及公众号原文标题的 2026-2027 不符；经三人次分别放大独立判读，三次均读作 2025-2026，本稿按「标题行照原卷」排作 2025-\/-2026------\textbf{全稿唯一一份年份不是 2026-\/-2027 的卷子}，见「高一32 专记」。\textbf{第二轮}（2026-09-26）又独立核过一遍，查出 \textbf{1 处漏抄}：第 12 题【解析】法二的行首等号链（\(\dfrac4{a^2+\allowbreak 2a}+\allowbreak \dfrac1{2b^2+b}=\allowbreak \dfrac{4x^2}{2x+\allowbreak 1}+\allowbreak \dfrac{y^2}{y+\allowbreak 2}\geqslant\allowbreak \cdots\)）被省成只剩右边，已按原卷补回并重编 \\ \addlinespace[1pt]
高一33 & 13 张 / 19 题 & \textbf{未发现转录错误}（2026-09-26 第二轮核对）。填空 3---12、选择 13---16、解答 17---21，无插图。逐题核对：第 3 题四个不等式的 \(>\)／\(\geqslant\allowbreak \) 与 \(\dfrac1{\sqrt{ab}}\) 一致；第 6、7、8、11、12 题的最小值／最大值与解析结论一致；第 9 题的 cases 与两条不等式链、第 19 题的 400／150／7200／14400 等报价数字、第 21 题的 \(\sqrt[n]{}\) 与 \(n\in\N^*\) 均无误；选择题答案 B、B、B、D 与解析末句一致。\textbf{原卷存疑三处}已照录加注（第 9 题解析的自相矛盾、第 19 题「一\textbf{面}」、第 5 题解析末句），并据核对补记两处排版差异（第 9 题「（\(a\in R\)）」移出 cases 花括号；文首注释④里原写的「第 18 题」其实没有 \(R\)）。\textbf{第二轮}（2026-09-26）又独立核过一遍，查出 \textbf{第 9 题【解析】漏排原卷重复的 cases 框}：原卷在解析里三次重新排出题干那个两行不等式组（第 3 页两处、第 4 页一处，其中一处是整行重复），源码只留「关于 \(x\) 的不等式组仅有一个整数解」、掉了框（有一处整行略去）------已按「照录」补回三处框与那一整行，并重编（答案版因此由 9 页增到 10 页）；另补记一条并列标点的排版差异⑨ \\ \addlinespace[1pt]
高一34 & 7 张 / 12 题 & \textbf{未发现转录错误}（2026-09-26 第二轮核对）。填空 3---10、解答 11---14（\textbf{无选择题}），全卷\textbf{无一处印「【答案】」}------12 题全部只有【解析】，答案写在解析末句。第 11(2) 按 \(a\) 分的五种情况与第 12(2) 的联立方程组逐行比对无串行漏行；各最小值（\(4/3\)、\(\dfrac{\sqrt5+\allowbreak 1}4\)、\(2\)、\(2\)、\(\dfrac{\sqrt2}2\)、\(3\)、\(\dfrac94\) 与 \(\dfrac32\)、\(\dfrac{4\sqrt6}3\)、\(4\) 与 \(\dfrac92\)、\(9\) 与 \(m\leqslant\allowbreak 8\)、\(\dfrac{\sqrt{10}}2\)）全部验算通过；标点逐句一致。唯一配图（第 13 题解析的坐标系）经比对确认画的正是原卷第 6 页那幅，无漏排。\textbf{第二轮}（2026-09-26）又独立核过一遍，查出 \textbf{1 处漏抄}：第 11(2)【解析】\(a<0\) 的三种情况原卷作「解 \((x+\allowbreak 3)\left(x-\dfrac4a\right)\leqslant\allowbreak 0\)，\textbf{得} \ldots\ldots」（三行同构），源码简写成「\textbf{解得} \ldots\ldots」（同题 \(a>0\) 那一支原卷本就作「解得」，故只这三行是漏抄）------已按原卷补回并重编。\textbf{原卷存疑五处}照录加注 \\ \addlinespace[1pt]
高一35 & 12 张 / 18 题 & \textbf{无符号／下标／标点类转录错误}（2026-09-26 第二轮核对）。填空 3---10、选择 11---13、解答 14---16、\textbf{附加题 17---20}（四节齐全），无插图。第 7 题三行不等式组、第 13 题大量三行组、第 18 题各两行方程组逐行对齐；指数与上下标（\((x-2)^{999}\)、\((x-1)^{2015}\)、\((a\pm1)^2\)、\(t_1\)---\(t_4\)）无误；区间开闭与正负号、第 16 题三小问的分段答案、第 18 题四组三元数对、第 19 题 \(a=\allowbreak -3\) 与结果 4、第 20 题的 \(c<x\leqslant\allowbreak  e\) 或 \(b\leqslant\allowbreak  x<d\) 全部一致。核对另指出原卷第 9、10 页在解析中有 \textbf{3 行重复陈述 \(S\)／\(T\) 定义}被源码省略，已按「照录」补回并重编。\textbf{第二轮}（2026-09-26）又独立核过一遍：\textbf{无实质转录错误}，只报出 \textbf{1 处压缩照录}------第 13 题【解析】把原卷照列的三个二次函数式（\(y=ax^2+bx+c\)、\(y=bx^2+cx+a\)、\(y=cx^2+ax+b\)）压成了「三个二次函数」，而同题 ① 处却是照列的，前后不一致，已改回照列并重编。\textbf{原卷存疑三处}照录加注 \\ \addlinespace[1pt]
高一36 & 6 张 / 17 题 & \textbf{2 处标点不一致，已订正}（2026-09-26 第二轮核对）：① 第 7 题题干原卷「\(<0\) 　\(x\in R\) 恒成立」中间只有空格、无逗号；② 第 10 题【解析】① 两处原卷用逗号（「\(k\in\Z\)，若 \(k=\allowbreak 2n\)」「故 \(m\in[1]_4\)，若 \(k=\allowbreak 2n+\allowbreak 1\)」），源码写成了分号。其余全部一致：第 13 题的 \(\varnothing\subset\allowbreak  M\subset\allowbreak \R\) 放大确认是真包含；cases／不等式链（第 4、5、7、8、10、16、19 题）逐行无误；\([k]_t\) 系列下标、\(a_1\)---\(c_2\)、\(\sqrt{5k+\allowbreak 1}\) 无误；答案（第 11 题 C、第 12 题 B、第 13 题 D、第 14 题 D）与解析一致；\(\dfrac{13}6\)、\(-\dfrac12<\dfrac ca\)、\(\dfrac79\)、\(1\leqslant\allowbreak  m<\dfrac75\) 或 \(25<m\leqslant\allowbreak 49\) 等数字均与原卷一致。无插图。\textbf{原卷存疑两处}照录加注（第 11 题「至少一个\textbf{数}为 1」、第 14 题称「四个命题」实列三条------后者可由真题原文确证是笔误） \\ \addlinespace[1pt]
高一37 & 10 张 / 19 题 & \textbf{正文与原图零差异}（2026-09-26 独立于转录的第二轮核对------另一名 worker、只读原图、不参考既有结论，并逐题独立验算答案）。填空 3---12、选择 13---16、解答 17---21，全卷无插图。19 题答案逐一验算一致；题干、选项、\texttt{\textbackslash{}ans} 与 \texttt{\textbackslash{}ifanswer} 全文逐句对照一字不差；\texttt{\textbackslash{}step}/\texttt{\textbackslash{}stepm} \textbf{98 条}与原卷【解析】的印刷行数逐题相等。\textbf{三处卷面异常}由它独立报出、已逐处放大核实并登记进文件头：① 第 20(3) 题干作「若\textbf{实数}\ldots」而解析作「若\textbf{正实数}\ldots，\textbf{且 \(a>b\)}」（解析自行加强了条件）；② 第 21(3) 解析由 \(a_1\geqslant\allowbreak 1\) 直接取严格「\(>1>\)」（对结论 \(n<10\) 无影响）；③ 第 21(2) 设问作「并求出 \(a_3\) 的\textbf{值}」（单数）而结论给出四个值。另第 17(1) 小问末原卷用冒号「：」（同题 (2) 用句点）------ 一并照录加注 \\ \addlinespace[1pt]
高一38 & 12 张 / 17 题 & \textbf{1 处须改 + 文件头漏登 3 条}（2026-09-26 两轮：先只读原图独立核对、再比对转录稿）。17 题答案全部验算一致。第 6 题 Venn 图（阴影确为 \(A\setminus\allowbreak  B\)、未镜像）、第 9 题 \(3\times\allowbreak 3\) 方格（tabular 九格逐格与原卷一致）、第 14 题作答括注位置------三处均核实「文件头说法成立」。\textbf{须改 1 处}：第 9 题【解析】原卷 9 行、本稿并成 8 条 \texttt{\textbackslash{}step}（原卷第 1 行「由表格数据得所有数之和为」单独成行），已拆开；现逐块条数 \([5,\allowbreak 5,\allowbreak 1,\allowbreak 9,\allowbreak 9,\allowbreak 9,\allowbreak 4,\allowbreak 9,\allowbreak 9,\allowbreak 12,\allowbreak 6,\allowbreak 14,\allowbreak 21,\allowbreak 13,\allowbreak 65]\)、\textbf{合计 191 = 原卷 191 行}------第 17 题末三处原卷折行（「\ldots 的」／「必要条件.」等）也据此按「一条 \texttt{\textbackslash{}step} ＝ 原卷一个印刷行」补拆。另报出文件头\textbf{漏登的 3 条}原卷问题（第 19 题题干漏符号 \(U\)、\(c_1+c_2=\allowbreak 2\) 与其上一行 \(2(c_1+c_2)=\allowbreak 21\) 不自洽、第 19(3) 的 \(A''=A\cup\allowbreak  A'\) 疑当作 \(2A\cup\allowbreak  A'\)）------已逐页放大复核后补登记并加注；文件头另两处计数（「三分完美数」实为 5 次、「非空数集」实为 1 处）也据它改正 \\ \addlinespace[1pt]
高一39 & 11 张 / 22 题 & \textbf{3 处须改 + 头注 2 处修正}（2026-09-26：先只读原图独立核对 22 题答案，再比对转录稿）。填空 3---12、选择 13---16、解答 17---20、\textbf{附加题 21---24}，共 22 题、3 张插图。\textbf{唯一答案性错误在原卷}：第 17 题【解析】推出「\(a\in\R\)」，与题干「有且仅有一个是真命题」矛盾（正确为 \((-\infty,\allowbreak -6]\cup\allowbreak (-5,\allowbreak +\allowbreak \infty)\)）------本稿\textbf{照录未改}。比对另查出 \textbf{3 处须改}：第 16 题解析的两处原卷作\textbf{数字 \(A_1\)}（同题另一处才是斜体 \(A_i\)），本稿统一成了 \(A_i\)------已改回；第 20(2) 的 11 个集合列表里 \(\{8,\allowbreak 11\}\) 与 \(\{9,\allowbreak 12\}\) 之间原卷误用\textbf{句号}，本稿顺手改成逗号------已照录回来；文件头④举例的题号错（第 9 题全文没有数集记号）------已改成「第 7、11、14、17、18、19 题」。\textbf{\texttt{\textbackslash{}step} 163 条与原卷印刷行逐题全等}（0 处漏拆／多拆）；\textbf{原卷笔误/存疑由 7 条补到 13 条}。另：本卷另一名 worker 独立写出的第二份转录与采纳稿做了 token 级比对，\textbf{相似度 0.9889、残差全为排版差异}（\texttt{\textbackslash{}\textbackslash{}{[}4pt{]}} vs \texttt{\textbackslash{}\textbackslash{}}、\texttt{\textbackslash{}dfrac} vs \texttt{\textbackslash{}frac}），数学内容一致------是对本稿内容的强旁证 \\ \addlinespace[1pt]
高一40 & 4 张 / 17 题 & \textbf{正文零须改}（2026-09-26 独立于转录的第二轮核对 + 转录稿比对）。填空 3---10、选择 11---14、解答 15---19，全卷 17 题、无插图。17 题答案逐一验算一致；题干、选项、\texttt{\textbackslash{}ans} 与 \texttt{\textbackslash{}ifanswer} 全文逐字照录，\textbf{7 处【解析】的 \texttt{\textbackslash{}step} 条数与原卷行数逐一相等（合计 27 = 27）}。三处原卷异常照录加注：第 13 题 A／D 两选项文字完全相同（原卷笔误）、第 5 题「全集 \(U=R\)」与「集合 \(M\)」之间漏排标点（本稿按 高一b 第 13 题「原卷漏排、本稿补上」先例补逗号并加注）、第 10 题「\ldots 的取值范围」漏「是」（文件头原缺此条，已补登记）。\textbf{副标题定为「不等式」}：17 题里约 15 题为不等式，集合仅第 4、5 两题、逻辑 0 题 \\ \addlinespace[1pt]
高一41 & 15 张 / 17 题 & \textbf{正文零须改}（2026-09-28 独立于转录的第二轮核对：另一名 worker \textbf{只读原图}先独立成稿、逐题独立验算答案，之后才打开转录稿比对）。填空 3---12、选择 13---16、解答 17---19，全卷 17 题、\textbf{无插图}。17 题答案独立验算，\textbf{16 题与卷面一致}；唯一不一致是\textbf{第 9 题的原卷之误}（末句答案漏了 \(t=\allowbreak -1\)，正确应为 \((-\infty,\allowbreak -2]\cup\allowbreak [-1,\allowbreak 1]\)），本稿\textbf{照录未改}。题干、选项、\texttt{\textbackslash{}ans} 与 \texttt{\textbackslash{}ifanswer} 全文逐句比对无差异；\texttt{\textbackslash{}step}／\texttt{\textbackslash{}stepm} \textbf{216 条与原卷印刷行逐题全等}（17／7／10／7／4／7／15／18／17／11／4／9／15／9／17／28／21）。历史高发项逐一放大核过：\(\subset\allowbreak \)（第 13 题，无下划线）与 \(\sub\)（第 14／16／17 题）分得清；第 16 题末句右端是\textbf{圆括号} \([-\frac{15}2,\allowbreak -7)\)；第 9 题的 \(-3\)、第 13 题的 \(\Rightarrow\)（无否定斜线）、第 12 题的 \(\{7\}\) 都确认\textbf{原卷如此}。另\textbf{查出文件头两处表述不准}并已订正（②⑤ 漏列第 10 题的斜体 \(R\)；②⑩ 原写「第 3 题两处」实为第 3、4、15 题各一处）。\textbf{副标题定为「集合与常用逻辑用语、不等式」}：17 题里不等式约 9---11 题、集合与常用逻辑用语 6 题 \\ \addlinespace[1pt]
高一42 & 8 张 / 19 题 & \textbf{无题面转录错误}（2026-10-02 独立第二轮逐页核图）。题干、答案与解析、题号及分节均与原图一致；原卷存疑四处照录，见本卷专记。原图第 02、08 页为 2560$\times$3620、其余六页 4762$\times$6735；第 11 题解析图为 TikZ 重排 \\ \addlinespace[1pt]
高一43 & 12 张 / 19 题 & \textbf{无题面转录错误}（2026-10-02 独立第二轮逐页核图）。题干、答案与解析、题号及分节均与原图一致；原卷第 13 题的答案与解析矛盾，照录加注。第 5 题解析图取自第 1 页原图，已放入 \texttt{\textbackslash{}ifanswer}，仅答案版显示 \\ \addlinespace[1pt]
高一44 & 8 张 / 12 题 & \textbf{无题面转录错误}（2026-10-02 独立第二轮逐页核图）；题号填空3--10、解答11--14。原卷错误/论证缺口加注：第3题漏论 \(a=\allowbreak 1\)、第4题条件不唯一、第10题对任意集合假设最大值、第12题端点包含错、第13题收入单位不换算 \\ \addlinespace[1pt]
高一45 & 10 张 / 19 题 & \textbf{无题面转录错误}（2026-10-02 独立第二轮逐页核图）；题号填空3--12、选择13--16、解答17--21。原卷错误加注：第11题补集推出关系不符、第15题解析与条件冲突、第18题真子集端点、第19题并集区间错误、第20题总费用算式与结论不符 \\ \addlinespace[1pt]
高一46 & 9 张 / 19 题 & \textbf{无题面转录错误}（2026-10-02 独立第二轮逐页核图）；题号填空3--12、选择13--16、解答17--21。第12题题干说50个实数解而解析按整数根计数，照录加注；三幅图均按原卷图片裁取 \\ \addlinespace[1pt]
高一47 & 9 张 / 18 题 & 独立逐页核对原图（2026-10-02）发现 \textbf{1 处转录错误}：第12题解析原图第4页以 \(x=\allowbreak 0,y=\allowbreak 3\) 为反例，源码曾误录 \(y=\allowbreak 2\)，已订正为 \(y=\allowbreak 3\)；其余题干、答案、解析、题号、选项与图均无差异。第9题集合 \(C\) 中未定义 \(y\)、附加3题「后西方数学家」均为原卷存疑，照录加注。 \\ \addlinespace[1pt]
\end{xltabular}

（体例沿用全稿：订正处一律在 \texttt{.tex} 里加 \texttt{\%} 注写明原卷原样与本稿错处，并重编试卷版/答案版。）

\textbf{另做过一次「半开半闭区间」专项扫描}（2026-09-15）：把（当时）17 份源码里的半开半闭区间全捞出来，
共 17 处、分布在 6 份；其中 12 处另一侧端点是 \(+\infty\)（如 \([-\frac23,+\infty)\)），写法唯一、
无歧义，剩下 \textbf{5 处两端都是有限数}的逐处回原图核了括号形状------\textbf{高一05 第 15 题的选项 B、D
与【解析】末句三处确有「开闭抄反」}（\((\frac43,\frac53]\) 抄成 \([\frac43,\frac53)\)，已订正），
另外 4 处（高一05 第 18(3) 题的 \((2,\frac52]\)、高一09 第 20 题的 \([\frac14,1)\)、高一03 第 7 题的
\([\frac32,2)\)、高一13 第 3 题的 \((\frac83,6]\)）与卷面逐字一致。\textbf{后入库的高一10 也同法扫过}：
它的半开半闭区间只有 \([6,+\infty)\) 与 \((-\infty,1)\) 两处，\textbf{另一侧端点都是无穷}，
没有「两端都是有限数」的情形。\textbf{2026-09-16 入库的高一22 同样扫过}：它有 \((0,1]\) 一处
两端有限（第 8 题【解析】），已放大核过括号形状------\textbf{左圆括、右方括}，与同句前面的
\(0<a\leqslant1\) 自洽；另有 \((-\infty,-4]\)、\(\left[\frac14,+\infty\right)\) 两处都是单侧无穷。

另把历史上\textbf{改过区间或端点}的题逐条回原图复核了一遍（共 5 处，含上面提到的第 15 题）：
高一05 第 17(2) 题的答案 \((-\infty,-3]\)（原卷【解析】末句即带 \(]\)）、第 10 题的「\(a<-4\) 或
\(a>2\)」（原卷末句作「\(2<a\)」，右端点不取）、高一08 第 17 题的 \(P=(-2,3)\)（\textbf{开}区间）、
高一07 第 21 题【解析】里的 \((t,t+2)\)（随「当 \(t<0<t+2\) 即 \(-2<t<0\) 时」一句照原卷录回）
------四处均与卷面一致，无需改动。

\begin{quote}
\textbf{区间开闭这件事到此已核完，不必重扫}：全库 17 处半开半闭区间（含 5 处两端有限的）与
历史上改过区间/端点的 5 处，都已逐处放大到看清括号形状核对过；只有高一05 第 15 题有错
并已订正。今后只需在\textbf{改动过某题的区间}时，重新核那一处。
\end{quote}

\clearpage
\section{材料流转}

各市重点中学高一数学备课组的校本周末作业/周练 $\to$ 公众号「\textbf{魔都数学加油站}」（作者署名
「破晓」，上海本地数学自媒体）收集整理后图文发布 $\to$ 本仓库 OCR 重排。

\begin{itemize}
\tightlist
\item
  46 篇已收录连载原文\textbf{逐篇核实}都出自该号（标题统一作「27学年高一 N------2026-2027年上海市$\times$$\times$中学
  高一上\ldots」，\texttt{N} 是该号的连载编号；高一46公众号标题有「你那」笔误，详见表注）。
\item
  \textbf{连载号与本地编号的对应关系}见下表。2026-09-15 逐篇抓原文标题核对过一遍：\texttt{N} \textbf{一律以标题
  为准}，不能按发布日期判------9-12 当天一次发了三篇。\textbf{本地数字编号一律与连载号对齐}（\texttt{高一NN}
  的 \texttt{NN} 就是连载号 \texttt{N}）；\textbf{没有连载号的三份家长拍照卷改用字母}编号（\texttt{高一a}／\texttt{高一b}／
  \texttt{高一c}），数字号段因此全部留给连载卷。
\item
  连载 2---47 共 46 篇\textbf{现已全部入库}。其中 \textbf{N=16}（大同中学 周末作业2.1）\textbf{没有从公众号原文
  另转一份}------经逐题核对，那份原文与家长拍照卷 \texttt{高一b} 是同一份作业（公众号版还略去填空
  第 1、2 题），故用 \texttt{高一b} 的内容做了编号对齐的 \texttt{高一16}；同理 \textbf{N=18}（大同中学 周末作业2.2，
  9-15 发布）就是拍照卷 \texttt{高一c}（见下表 N=16／18 两行与「试卷清单」）。
  \textbf{N=23／24}（9-18／9-19 发布，格致 · 周末作业3、复旦附中 · 周末作业3）是 2026-09-20 入库的
  两篇，\textbf{N=25---28}（9-21 一次发了四篇：华二 · 周末作业3、交附闵行 · 周练2、
  交附闵行 · 周末作业9.18、七宝 · 周末作业3）是 2026-09-21 入库的四篇，六篇均为公众号原文直转，
  没有对应的拍照卷；\textbf{N=42／43}（复旦附中 · 周末作业4、交附闵行 · 周练3）于 2026-10-01 发布、
  2026-10-02 入库；\textbf{N=44---46}（交附闵行 · 国庆作业1---3）与 \textbf{N=47}（七宝中学 · 抱孩子练习9.26）均于 2026-10-02 发布、2026-10-02 入库，均为公众号原文直转。
\end{itemize}

\begingroup\small
\newcommand{\tabhdrD}{\rowcolor{white} 连载 N & 原文标题里的学校 · 作业 & 发布 & 本地编号 & 状态 \\}
\rowcolors{1}{white}{rowgray}
\begin{xltabular}{\textwidth}{@{}>{\centering\arraybackslash}p{1.45cm}>{\raggedright\arraybackslash}p{5.55cm}>{\raggedright\arraybackslash}p{1.10cm}>{\centering\arraybackslash}p{1.93cm}>{\raggedright\arraybackslash}p{3.86cm}@{}}
\toprule
\tabhdrD
\midrule
\endfirsthead
\multicolumn{5}{r}{\footnotesize（续上表）} \\
\midrule
\tabhdrD
\midrule
\endhead
\bottomrule
\endlastfoot
2 & 大同中学 · 集合综合练习 & 9-05 & 高一02 & 已入库，编号一致 \\ \addlinespace[1pt]
3 & 格致中学 · 周末作业1 & 9-05 & 高一03 & 已入库，编号一致 \\ \addlinespace[1pt]
4 & 位育中学 · 周末作业1 & 9-05 & 高一04 & 已入库，编号一致 \\ \addlinespace[1pt]
5 & 七宝中学 · 周末作业9.5 & 9-06 & 高一05 & 已入库，编号一致 \\ \addlinespace[1pt]
6 & 延安中学 · 周末作业1 & 9-07 & 高一06 & 已入库，编号一致 \\ \addlinespace[1pt]
7 & 交附闵行 · 周练1 & 9-08 & 高一07 & 已入库，编号一致 \\ \addlinespace[1pt]
8 & 七宝中学 · 周末作业1 & 9-08 & 高一08 & 已入库，编号一致 \\ \addlinespace[1pt]
9 & 进才中学 · 周末作业1 & 9-11 & 高一09 & 已入库，编号一致 \\ \addlinespace[1pt]
\textbf{10} & 复附浦东 · 周末作业1 & 9-12 & 高一10 & 已入库，编号一致 \\ \addlinespace[1pt]
11 & 复附浦东 · 周末作业2 & 9-12 & 高一11 & 已入库，编号一致 \\ \addlinespace[1pt]
12 & 交附闵行 · 周末作业9.11 & 9-12 & 高一12 & 已入库，编号一致 \\ \addlinespace[1pt]
13 & 格致中学 · 周末作业2 & 9-12 & 高一13 & 已入库，编号一致 \\ \addlinespace[1pt]
14 & 华师大二附中 · 集合与逻辑单元练习 & 9-13 & 高一14 & 已入库，编号一致 \\ \addlinespace[1pt]
15 & 七宝中学 · 抱孩子练习9.8 & 9-13 & 高一15 & 已入库，编号一致 \\ \addlinespace[1pt]
16 & 大同中学 · 周末作业2.1 & 9-13 & 高一16（对号副本）＋ \textbf{高一b（补）}（同一份拍照卷） & 已入库，编号一致 \\ \addlinespace[1pt]
17 & 七宝中学 · 周末作业2 & 9-14 & 高一17 & 已入库，编号一致 \\ \addlinespace[1pt]
18 & 大同中学 · 周末作业2.2 & 9-15 & 高一18（对号副本）＋ \textbf{高一c（补）}（同一份拍照卷） & 已入库，编号一致 \\ \addlinespace[1pt]
19 & 延安中学 · 周末作业2 & 9-16 & 高一19 & 已入库，编号一致 \\ \addlinespace[1pt]
20 & 进才中学 · 周末作业2 & 9-16 & 高一20 & 已入库，编号一致 \\ \addlinespace[1pt]
21 & 上师大宝山 · 限时练习1 & 9-16 & 高一21 & 已入库，编号一致 \\ \addlinespace[1pt]
22 & 七宝中学 · 抱孩子练习9.15 & 9-16 & 高一22 & 已入库，编号一致 \\ \addlinespace[1pt]
23 & 格致中学 · 周末作业3 & 9-18 & 高一23 & 已入库，编号一致 \\ \addlinespace[1pt]
24 & 复旦附中 · 周末作业3 & 9-19 & 高一24 & 已入库，编号一致 \\ \addlinespace[1pt]
25 & 华师大二附中 · 周末作业3 & 9-21 & 高一25 & 已入库，编号一致 \\ \addlinespace[1pt]
26 & 交附闵行 · 周练2 & 9-21 & 高一26 & 已入库，编号一致 \\ \addlinespace[1pt]
27 & 交附闵行 · 周末作业9.18 & 9-21 & 高一27 & 已入库，编号一致 \\ \addlinespace[1pt]
28 & 七宝中学 · 周末作业3 & 9-21 & 高一28 & 已入库，编号一致 \\ \addlinespace[1pt]
29 & 延安中学 · 周末作业3 & 9-23 & 高一29 & 已入库，编号一致 \\ \addlinespace[1pt]
30 & 南洋模范中学 · 开学考 & 9-23 & 高一30 & 已入库，编号一致 \\ \addlinespace[1pt]
31 & 上海中学 · 周练1 & 9-25 & 高一31 & 已入库，编号一致 \\ \addlinespace[1pt]
32 & 华师大二附中 · 周末作业4 & 9-25 & 高一32 & 已入库，编号一致（卷首年份照原卷作 \textbf{2025-2026}，与连载的 2026-2027 不同，见「高一32 专记」） \\ \addlinespace[1pt]
33 & 交附闵行 · 中秋作业 & 9-25 & 高一33 & 已入库，编号一致 \\ \addlinespace[1pt]
34 & 交附闵行 · 查漏补缺卷 & 9-25 & 高一34 & 已入库，编号一致 \\ \addlinespace[1pt]
35 & 七宝中学 · 抱孩子练习9.22 & 9-25 & 高一35 & 已入库，编号一致 \\ \addlinespace[1pt]
36 & 格致中学 · 周末作业4 & 9-25 & 高一36 & 已入库，编号一致 \\ \addlinespace[1pt]
37 & 七宝中学 · 周末作业4 & 9-26 & 高一37 & 已入库，编号一致 \\ \addlinespace[1pt]
38 & 大同中学 · 周末作业3 & 9-26 & 高一38 & 已入库，编号一致 \\ \addlinespace[1pt]
39 & 大同中学 · 周末作业4.1 & 9-26 & 高一39 & 已入库，编号一致 \\ \addlinespace[1pt]
40 & 大同中学 · 周末作业4.2 & 9-26 & 高一40 & 已入库，编号一致 \\ \addlinespace[1pt]
41 & 上海中学 · 周练2 & 9-28 & 高一41 & 已入库，编号一致 \\ \addlinespace[1pt]
42 & 复旦附中 · 周末作业4 & 10-01 & 高一42 & 已入库，编号一致 \\ \addlinespace[1pt]
43 & 交附闵行 · 周练3 & 10-01 & 高一43 & 已入库，编号一致 \\ \addlinespace[1pt]
44 & 交附闵行 · 国庆作业1 & 10-02 & 高一44 & 已入库，编号一致 \\ \addlinespace[1pt]
45 & 交附闵行 · 国庆作业2 & 10-02 & 高一45 & 已入库，编号一致 \\ \addlinespace[1pt]
46 & 交附闵行 · 国庆作业3 & 10-02 & 高一46 & 已入库，公众号标题有「你那」二字，卷面标题不含 \\ \addlinespace[1pt]
47 & 七宝中学 · 抱孩子练习9.26 & 10-02 & 高一47 & 已入库，编号与连载一致 \\ \addlinespace[1pt]
--- & 大同中学 · 周中小练习1 & --- & 高一a（补） & 无连载号，家长拍照卷（原卷未印校名，判为大同） \\ \addlinespace[1pt]
--- & 大同中学 · 第二周周末练习2 & --- & 高一c（补） & 无连载号，家长拍照卷 \\ \addlinespace[1pt]
\end{xltabular}
\endgroup

\begin{itemize}
\tightlist
\item
  \textbf{为什么拍照卷用字母}：三份卷的原始材料是家长拍摄的学生作业，\textbf{不经公众号、原卷也没有
  连载号}，若沿用数字就会与连载卷抢号（此前正是「重号 + 加「（补）」」的做法，需靠校名才能
  区分）。现改为字母前缀，\texttt{高一02}---\texttt{高一28} 全部留给连载卷，一一对应、不再有歧义。
  \textbf{例外是 \texttt{高一16} 与 \texttt{高一18}}：它们的内容来自拍照卷（见上表 N=16／18 两行），但为了让数字号段
  与连载号一一对应，各出了一份\textbf{编号对齐的副本}（\texttt{高一16\_大同中学\_第二周周末练习1}、
  \texttt{高一18\_大同中学\_第二周周末练习2}，卷面标题都不带「（补）」），内容分别与 \texttt{高一b（补）}、
  \texttt{高一c（补）} 一字不差------两份都留着：带「（补）」的作拍照原件的转录本、对号的那份作编号副本。
  \textbf{改动其一时须同步改另一份。}
\item
  \textbf{「（补）」怎么用}：只用在\textbf{三份家长拍照卷}上（\texttt{高一a\_大同中学\_周中小练习1（补）}、
  \texttt{高一b\_大同中学\_第二周周末练习1（补）}、\texttt{高一c\_大同中学\_第二周周末练习2（补）}），表示
  本篇不在连载序列内（原卷标题也没有这三个字，是\textbf{本稿所加}，已写进各卷头部注释）。
  其余各卷编号与连载号一致，文件名与卷面标题都\textbf{不带}任何标记。
\item
  高一a---高一c 是家长拍摄的学生作业，不经公众号。\textbf{旁证}：高一a 与 高一b 有两道题
  \textbf{逐字相同}（\(A=\{(x,y)\mid y=x^2\}\) 求交、满足 \(M\cap\{a_1,a_2,a_3\}=\{a_1,a_2\}\)
  的 \(M\) 个数），而 高一b 卷眉印「上海市大同\ldots」、高一a 原卷未印校名------\textbf{据此把 高一a 判为大同
  中学}，文件名记为 \texttt{高一a\_大同中学\_周中小练习1（补）}（卷面标题仍照原卷、不带校名）。
\end{itemize}

\clearpage
\section{上游是「共享题库」：48 组逐字同题 + 20 组同模板变体}

把 \textbf{47 份试卷的 824 道}一级题干两两交叉比对（跨卷 \textbf{339,076 对}），校际重复分两档，两档的
成员都按卷面题号逐条调出、逐字核过题干。判据见「方法与局限」。（本表是 2026-10-02
\textbf{47 份全量重跑并人工复核}后的口径；17 份、18 份、21 份、22 份、24 份、28 份、29 份、30 份、36 份、40 份、41 份、43 份
那几轮的组数见「核对摘要」的注。高一16／高一18 分别与 高一b／高一c 同卷，不重复计入。

\textbf{逐字同题 48 组}（骨架最长公共子串 $\geqslant$ 25 字且含 $\geqslant$ 8 个汉字，或整段相似度 $\geqslant$ 0.95------
几乎全是符号的题另有一条\textbf{符号档}，见「方法与局限」；★ 为 2026-09-15 重跑新增，
$\oplus$ 为补跑高一10 新增，▲ 为 2026-09-15 末并入高一19---21 后的新增或改判，
◇ 为 2026-09-22 并入高一25---28 后新增）：

\newcommand{\tabhdrE}{\rowcolor{white} 题目 & 份数 & 分布 \\}
\rowcolors{1}{white}{rowgray}
\begin{xltabular}{\textwidth}{@{}>{\raggedright\arraybackslash}p{5.83cm}>{\centering\arraybackslash}p{1.23cm}>{\raggedright\arraybackslash}p{7.68cm}@{}}
\toprule
\tabhdrE
\midrule
\endfirsthead
\multicolumn{3}{r}{\footnotesize（续上表）} \\
\midrule
\tabhdrE
\midrule
\endhead
\bottomrule
\endlastfoot
\(A=\allowbreak \{x\mid\allowbreak  2a+\allowbreak 1\leqslant\allowbreak  x\leqslant\allowbreak  3a-5\}\)，\(B=\allowbreak \{x\mid\allowbreak  x<0\text{ 或 }x>19\}\) & 3 & 02 第 10 题、a 第 8 题、c 第 9 题 \\ \addlinespace[1pt]
对加减封闭的数集 \(S\)（\(x\pm y\in S\)，问最小的正数） & 3 & 05 第 12 题、09 第 21 题、b 第 18 题　\textbf{数字有变体}：05 版是「最小正数为 6」的填空，09／b 版是「最小正数为 5」的三小问 \\ \addlinespace[1pt]
集合 \(M\) 的 \(k\) 个不同\textbf{非空}子集 \(N_1\dots N_k\) 且并集 \(\neq M\)，求 \(k\) 最大 & 2 & 07 第 15 题、11 第 16 题　▲\textbf{05 第 16 题已移出本组}：它漏了「非空」二字、\(M\) 也换成 \(\{a,b,c,d,e\}\)，答案 16 与 07／11 的 15 不同，两两只有 \(0.797\)／\(0.813\)，改列变体档 \\ \addlinespace[1pt]
\(A=\allowbreak \{x\mid\allowbreak  x^2+\allowbreak 4x=\allowbreak 0\}\) 与含参 \(B=\allowbreak \{x\mid\allowbreak  x^2+\allowbreak 2(a+\allowbreak 1)x+a^2-1=\allowbreak 0\}\) 的子集讨论 & \textbf{4} & 07 第 18 题、08 第 18 题、11 第 19 题、\textbf{▲09 第 19 题}　▲09 版题干与前三份\textbf{一字不差}（\(r=\allowbreak 0.840\)），只是只问「\(B\subseteq\allowbreak  A\) 求 \(a\)」；它原先是作为「单成员变体」附在本表之外的，本轮并入本组 \\ \addlinespace[1pt]
「和谐集」（去掉任一元素后余下元素可平分） & \textbf{3} & 02 第 19 题、09 \textbf{附加题}（卷面无编号）、\textbf{◇38 第 13 题}　\textbf{题干逐字相同}：09 版是《上海市上海中学东校 2023--2024 学年高一上期中卷》的原题（三小问），02 版是它的改编（只问「元素个数的最小值」）------见「附加题专查」；\textbf{38 版是同一「和谐集」定义的第三份}（正整数组 \(A=\allowbreak \{a_1,a_2,\allowbreak \cdots,a_n\}\)、去掉其中任一元素后余下元素可平分，38 版另起名叫「平衡集」，见「高一38 专记」） \\ \addlinespace[1pt]
\(f(x)=\allowbreak \dfrac1{1-x}\) 迭代（\(x\in A\Rightarrow\dfrac1{1-x}\in A\)） & 2 & 03 第 18 题、c 第 19 题 \\ \addlinespace[1pt]
全集 \(U=\allowbreak \{x\in\mathbb N\mid\allowbreak  x\leqslant\allowbreak  9\}\)，由 \(A\cup\allowbreak  B\)、\(A\cap\allowbreak \overline{B}\)、\(A\cap\allowbreak  B\) 反求 \(A\) & 2 & 05 第 9 题、07 第 8 题 \\ \addlinespace[1pt]
\(\{0,\allowbreak -1,\allowbreak 2a\}=\allowbreak \{a-1,\allowbreak -\lvert a\rvert,a+\allowbreak 1\}\)，求 \(a\) & 2 & 07 第 5 题、08 第 5 题 \\ \addlinespace[1pt]
满足 \(M\cap\allowbreak \{a_1,a_2,a_3\}=\allowbreak \{a_1,a_2\}\) 的 \(M\) 个数 & 2 & a 第 4 题、b 第 7 题　\textbf{这一组同时见于上面的高考出处表}（改编自 2008 山东卷（文）选择 1） \\ \addlinespace[1pt]
\(A=\allowbreak \{(x,y)\mid\allowbreak  y=x^2\}\) 与 \(B=\allowbreak \{(x,y)\mid\allowbreak  (y-1)/(x-1)=\allowbreak \tfrac12\}\) 求 \(A\cap\allowbreak  B\) & 2 & a 第 5 题、b 第 8 题 \\ \addlinespace[1pt]
★「孤立元」（\(S=\allowbreak \{1,\allowbreak \dots,\allowbreak 8\}\) 的 3 元子集不含孤立元） & 2 & 06 第 9 题、15 附加 17 题　\textbf{两版逐字相同，且同为题源 2009 北京卷（文）填空 6 的原题照录} \\ \addlinespace[1pt]
★「好集」（\(0,\allowbreak 1\in A\)；\(x-y\in A\)，\(x\neq0\) 时 \(1/x\in A\)） & 2 & 08 第 21 题、14 第 20 题　\textbf{定义与 (1)(2) 问逐字相同}；14 版第 (3) 问多一个命题 \(q\)（\(y/x\in A\)） \\ \addlinespace[1pt]
★非空数集 \(A\) 的 \(S=\allowbreak \{x\mid\allowbreak  x=a+b\}\)、\(T=\allowbreak \{x\mid\allowbreak  x=\allowbreak \lvert a-b\rvert\}\) & 2 & 07 第 21 题、15 第 16 题　\textbf{定义与 (2) 问逐字相同}，只换 \(A\)（\(\{2,\allowbreak 5\}\)／\(\{1,\allowbreak 3\}\)）与第 (3) 问上界（2025／2021） \\ \addlinespace[1pt]
★\(\alpha:4\leqslant\allowbreak  x\leqslant\allowbreak  8\)，\(\beta:m+\allowbreak 1\leqslant\allowbreak  x\leqslant\allowbreak  2m+\allowbreak 4\)，\(\alpha\) 是 \(\beta\) 的充分不必要条件 & 2 & 11 第 11 题、12 第 4 题　\textbf{逐字相同}（\(r=\allowbreak 0.987\)） \\ \addlinespace[1pt]
★\(A=\allowbreak \{x\mid\allowbreak  x=a^2+\allowbreak 2a+\allowbreak 4,a\in\mathbb{R}\}\) 与 \(B=\allowbreak \{y\mid\allowbreak  y=b^2-4b+\allowbreak 3,b\in\mathbb{R}\}\) 的关系 & 2 & 11 第 8 题、12 第 3 题　\textbf{逐字相同}（\(r=\allowbreak 0.971\)） \\ \addlinespace[1pt]
$\oplus$集合 \(\{(x,y)\mid\allowbreak  x+\allowbreak 2y=\allowbreak 2,\allowbreak \ x,y\in\mathbb{N}\}\) 用列举法表示 & 2 & 10 第 4 题、11 第 7 题　\textbf{逐字相同}（\(r=\allowbreak 1.000\)）------同一学校相邻两次作业用了同一道题 \\ \addlinespace[1pt]
$\oplus$\(A=\allowbreak \{x\mid\allowbreak  x^2-3x-4=\allowbreak 0\}\)、\(B=\allowbreak \{x\mid\allowbreak  ax-1=\allowbreak 0\}\)，\(A\cap\allowbreak  B=B\) 求 \(a\) & 2 & 05 第 6 题、10 第 9 题　\textbf{逐字相同}（\(r=\allowbreak 0.955\)）；同族还有 c 第 3 题（改问「充分非必要条件」，见变体档）与 a 第 9 题（把 \(A\) 换成含参方程、改问「只有一个真子集」） \\ \addlinespace[1pt]
▲\(A=\allowbreak \{x\mid\allowbreak  x^2-3x+\allowbreak 2=\allowbreak 0\}\) 与含参 \(B=\allowbreak \{x\mid\allowbreak  x^2+\allowbreak 2(a+\allowbreak 1)x+a^2-5=\allowbreak 0\}\) & 2 & 02 第 17 题、05 第 17 题　\textbf{两版集合一字不差}，只差「设集合／已知集合」；\textbf{小问有出入}------(1) 问 02 作 \(A\cap\allowbreak  B=\allowbreak \{2\}\)、05 作「\(B\) 为单元素集合」，(2) 问两版相同（\(A\cup\allowbreak  B=A\)，答案同为 \((-\infty,\allowbreak -3]\)）。原列变体档首行，本轮按题干逐字改入本档 \\ \addlinespace[1pt]
▲「伙伴关系集合」（\(x\in A\Rightarrow\frac1x\in A\)），\(M=\allowbreak \{-1,\allowbreak 0,\allowbreak \frac13,\allowbreak \frac12,\allowbreak 1,\allowbreak 2,\allowbreak 3,\allowbreak 4\}\) & 2 & 02 第 6 题、21 第 10 题　定义措辞略有出入（「若 \(x\in A\)，则 \(\frac1x\in A\)，就称 \(A\) 是\ldots」／「若对任意的 \(x\in A\)，都有 \(\frac1x\in A\)，则称 \(A\) 是\ldots」），集合与问法相同，\textbf{答案同为 15} \\ \addlinespace[1pt]
▲\(A=\allowbreak \{x\mid\allowbreak  ax^2+x+\allowbreak 1=\allowbreak 0\}\) 只有一个元素 & 2 & 06 第 5 题、20 第 4 题　\textbf{逐字相同}（\(r=\allowbreak 0.939\)），只差「若集合／已知集合」，答案同为 \(a=\allowbreak 0\) 或 \(\frac14\) \\ \addlinespace[1pt]
▲四个数字元素的集合 \(S\) 的\textbf{全子集元素和}（空集记 0） & 2 & 07 第 9 题、19 第 11 题　\textbf{只换数字}：07 版和 \(16208\)（元素和 \(2026\)）、19 版和 \(16192\)（\(2024\)） \\ \addlinespace[1pt]
▲六个关系式（\(\varnothing\)、\(\{a\}\)、\(a\) 与 \(\{a,b\}\) 间的 \(\in/\subseteq\allowbreak \)） & 2 & 09 第 13 题、19 第 8 题　\textbf{只换记号}：①② 09 印 \(\subsetneq\allowbreak \)、19 印 \(\subseteq\allowbreak \)；因 \(\varnothing\subsetneq\allowbreak \{a\}\) 与 \(\varnothing\subseteq\allowbreak \{a\}\) 同真、\(a\subsetneq\allowbreak \{a\}\) 与 \(a\subseteq\allowbreak \{a\}\) 同假，两版答案都是 ①③⑤ \\ \addlinespace[1pt]
▲\(A=\allowbreak \{x\mid\allowbreak  x^2-ax+a^2-19=\allowbreak 0\}\)、\(B=\allowbreak \{x\mid\allowbreak  x^2-5x+\allowbreak 6=\allowbreak 0\}\)、\(C=\allowbreak \{x\mid\allowbreak  x^2+\allowbreak 2x-8=\allowbreak 0\}\) & 2 & 15 第 14 题、a 第 14 题　\textbf{三式一字不差}（\(r=\allowbreak 0.978\)，公共串 44 字）；15 版两小问、a 版三小问（a 版多「\(A\cap\allowbreak  B=A\cup\allowbreak  B\)」一问），公共的后两问相同。原列「判据之外手工补记」，本轮改入本档 \\ \addlinespace[1pt]
▲「特征数」定义题（把 \(M\) 分成三个集合，\(X_i\) 取「最大＋最小」之和） & \textbf{3} & 07 第 10 题、a 第 13 题、\textbf{◇39 第 16 题}　同一个自定义概念，差别在 \(M\)（\(1\sim12\)／\(0<x\leqslant\allowbreak 15\)）、每组元素个数（4／5）与问法（求最大值与最小值的积 \(=\allowbreak 1485\)／问哪个值不可能）；公共串 27 字全落在「特征数」的定义句上。原列「判据之外手工补记」，本轮改入本档；\textbf{39 版是同一「特征数」定义的第三份}（\(M=\allowbreak \{x\mid\allowbreak  x\in\N,\allowbreak 0<x\leqslant\allowbreak 15\}\)、每组恰 5 个元素、问 \(X_1+X_2+X_3\) 的取值范围） \\ \addlinespace[1pt]
◇\(x^2+ax+a^2-8=\allowbreak 0\) 两根 \(x_1\)、\(x_2\) 且 \(\dfrac1{x_1}+\allowbreak \dfrac1{x_2}=\allowbreak -\dfrac12\)，求 \(a\) & 3 & 12 第 9 题、26 第 11 题、\textbf{44 第 7 题}　\textbf{同一所学校（交附闵行）沿用同一道题}：题干与设问一致，答案均为 \(-2\)；本轮把该组由 2 份扩为 3 份 \\ \addlinespace[1pt]
◇\(P_1\)／\(P_2\)／\(Q_1\)／\(Q_2\) 四集合版「下列说法正确的是」 & 2 & 17 第 16 题、25 第 16 题　\textbf{题干、四个选项与答案 B 逐字相同}，同出 2015 上海卷（春）；差别只在 17 那版的 C、D 两项把真题的「存在 \(a\)」印成「对任意 \(a\)」（原卷笔误，见「真题出处」表），25 这版与真题一致 \\ \addlinespace[1pt]
◇已知集合 \(A=\allowbreak \{x\mid\allowbreak  x<-2\text{ 或 }x\geqslant\allowbreak 3\}\)，\(U=\allowbreak \R\)（三小问） & 2 & 11 第 20 题、30 第 19 题　\textbf{题干、三个小问与三个答案全部一字不差}（复附浦东 $\times$ 南洋模范，跨校用了同一道题，连解析也同文） \\ \addlinespace[1pt]
◇集合 \(A=\allowbreak \{1,\allowbreak 2\}\)、\(B=\allowbreak \{x\mid\allowbreak  ax=\allowbreak 6\}\)，\(A\cup\allowbreak  B=A\) 求 \(a\) 的取值集合 & 2 & 15 第 6 题、30 第 8 题　\textbf{逐字相同}，答案同为 \(\{0,\allowbreak 3,\allowbreak 6\}\)；同族另有 25 第 6 题（另一型，见变体档） \\ \addlinespace[1pt]
◇先定义 \(M-N\) 与 \(M\oplus N=\allowbreak (M-N)\cup\allowbreak (N-M)\)，再求 \(A\oplus B\) & 2 & 22 第 5 题、30 第 9 题　\textbf{定义那两句一字不差}（占题干大半，故判据把它判进逐字档）；两卷给的 \(A\)、\(B\) 不同（22 版 \(A=\allowbreak \left\{x\mid\allowbreak  x\geqslant\allowbreak -\dfrac94\right\}\)、\(B=\allowbreak \{x\mid\allowbreak  x<0\}\)，30 版 \(A=\allowbreak [0,\allowbreak +\allowbreak \infty)\)、\(B=\allowbreak (-\infty,\allowbreak 2]\)），答案随之不同 \\ \addlinespace[1pt]
◇区间 \(M=\allowbreak \left[m,m+\allowbreak \dfrac34\right]\)、\(N=\allowbreak \left[n-\dfrac13,n\right]\) 含于 \([0,\allowbreak 1]\)，求 \(M\cap\allowbreak  N\)``长度''的最小值 & 2 & 25 第 12 题、30 第 15 题　\textbf{逐字相同}，答案同为 \(\dfrac1{12}\) \\ \addlinespace[1pt]
◇集合 \(A=\allowbreak \{x\mid\allowbreak  y=\allowbreak \sqrt{1-x^2},x\in\R\}\)、\(B=\allowbreak \{y\mid\allowbreak  y=x^2+\allowbreak 1,x\in A\}\)，求 \(A\cup\allowbreak  B\) & 2 & 03 第 8 题、31 第 3 题　\textbf{逐字相同}（答案同为 \([-1,\allowbreak 2]\)） \\ \addlinespace[1pt]
◇``\(\dfrac{a_1}{a_2}=\allowbreak \dfrac{b_1}{b_2}=\allowbreak \dfrac{c_1}{c_2}\)''是``\(M=N\)''的什么条件 & 2 & 20 第 15 题、36 第 13 题　\textbf{逐字相同}（同出 2003 上海卷（秋理），见出处表） \\ \addlinespace[1pt]
◇非空集合 \(A\) 具有性质：\(\dfrac xy\in A\)、\(x+y\in A\)，判断哪个说法错误 & 2 & 24 第 13 题、32 第 16 题　\textbf{题干与四个选项逐字相同}，答案同为 D \\ \addlinespace[1pt]
◇\(A=\allowbreak \{x\mid\allowbreak  x^2+\allowbreak 2x-8\geqslant\allowbreak 0\}\)、\(B=\allowbreak \{x\mid\allowbreak  x^2-2ax+\allowbreak 4\leqslant\allowbreak 0\}\)，\(A\cap\allowbreak  B\cap\allowbreak \N\) 恰有 \(2\) 个元素 & 2 & 25 第 9 题、36 第 8 题　\textbf{逐字相同}，答案同为 \(\left[\dfrac{13}6,\allowbreak \dfrac52\right)\) \\ \addlinespace[1pt]
◇\(\lvert mx^2-4x+\allowbreak 3\rvert\geqslant\allowbreak 2\) 的必要非充分条件是``\(x\leqslant\allowbreak 1\) 或 \(x\geqslant\allowbreak 4\)'' & 2 & 27 第 10 题、32 第 10 题　\textbf{逐字相同}，答案同为 \(\left(\dfrac45,\allowbreak \dfrac{15}{16}\right]\)；两卷连解析也同文 \\ \addlinespace[1pt]
◇已知 \(a\)、\(b\in\R^+\allowbreak \) 且 \(ab=\allowbreak 2\)，求 \(\dfrac b{2+a^2}+\allowbreak \dfrac a{2+b^2}\) 的最小值 & 2 & 33 第 6 题、34 第 7 题　\textbf{逐字相同}（同属交附闵行，答案同为 \(\dfrac{\sqrt2}2\)） \\ \addlinespace[1pt]
◇已知实数 \(x\)、\(y\) 满足 \(3x^2+\allowbreak 3y^2-2xy=\allowbreak 4\)，求 \(x+y\) 的最大值 & 2 & 33 第 8 题、34 第 5 题　\textbf{逐字相同}（同上，两卷连解析也同文，答案同为 \(2\)） \\ \addlinespace[1pt]
◇``\(M\)、\(P\)、\(S\) 是全集 \(I\) 的三个子集，则阴影部分所表示的集合是'' & \textbf{2} & 08 第 15 题、32 第 13 题　\textbf{题干与四个选项逐字相同}：两版的图形\textbf{也一样}（\(S\) 内含于 \(P\)、阴影是 \((M\cap\allowbreak  P)\setminus\allowbreak  S\)），答案同为 C。\textbf{注意 19 第 7 题虽被判据拉进同一组，但它的图形是 \(P\) 内含于 \(S\)、答案不同，不是同题}（沿用既有结论，见「回原图逐题核对」） \\ \addlinespace[1pt]
◇\(A=\allowbreak \{x\mid\allowbreak  a\leqslant\allowbreak  x\leqslant\allowbreak -a+\allowbreak 3\}\)、\(B=\allowbreak \{x\mid\allowbreak  x<-1\) 或 \(x>5\}\) & 2 & 03 第 16 题、\textbf{38 第 15 题}　\textbf{题干逐字相同}（公共串 12 字全落在 \(B\) 的定义上），按条件讨论 \(a\) 的范围 \\ \addlinespace[1pt]
◇对于集合 \(M=\allowbreak \{a\mid\allowbreak  a=x^2-y^2,\allowbreak \ x\in\Z,\allowbreak \ y\in\Z\}\)，给出如下三个结论 & 2 & 07 第 11 题、\textbf{38 第 14 题}　\textbf{题干逐字相同}（\(M\) 的定义句一字不差），判断正确结论的个数 \\ \addlinespace[1pt]
◇关于 \(x\) 的不等式组 \(\begin{cases}x^2-x-2>0\\ 2x^2+\allowbreak (2k+\allowbreak 5)x+\allowbreak 5k<0\end{cases}\) 的整数解的集合为 \(\{-2\}\) & 2 & 28 第 7 题、\textbf{◇39 第 9 题}　\textbf{题干逐字相同}（同上------两个不等式、整数解恰为 \(\{-2\}\)），求 \(k\) 的范围 \\ \addlinespace[1pt]
◇``\(x\geqslant\allowbreak  a\)''是``\(0<x<2\)''的必要非充分条件 & 2 & a 第 7 题、\textbf{◇39 第 3 题}　\textbf{题干逐字相同}（充要条件与区间端点的关系），答案同为 \(a\leqslant\allowbreak 0\) \\ \addlinespace[1pt]
◇高斯取整函数 \(y=\allowbreak [x]\)（题面先介绍「数学王子」高斯再给定义） & 2 & a 第 10 题、\textbf{◇39 第 11 题}　\textbf{介绍段与定义句逐字相同}（a 版问取值个数、39 版问 \(x\) 的范围） \\ \addlinespace[1pt]
◇规定集合 \(M=\allowbreak \{a_1,a_2,\allowbreak \cdots,a_n\}\) 的子集 \(\{a_{i_1},a_{i_2},\allowbreak \cdots,a_{i_m}\}\) 为 \(M\) 的第 \(k\) 个子集 & 2 & a \textbf{附加 2 题}、\textbf{◇39 附加 21 题}　\textbf{定义句逐字相同}（自定义「第 \(k\) 个子集」后的枚举题） \\ \addlinespace[1pt]
高一交附闵行同校的 \(A\subseteq\allowbreak  B\) 参数集合题 & 2 & 26 第 8 题、43 第 7 题　\textbf{逐字同题}：两题的 \(A\)、\(B\) 定义、包含关系设问与答案范围均相同；另与 04 第 15 题构成同模板变体，见下表 \\ \addlinespace[1pt]
\(a>\lvert b\rvert\) 与 \(a^2>b^2\) 的充要条件判断题 & 2 & 08 第 14 题、45 第 14 题　\textbf{题干、四个选项及答案 A 完全相同} \\ \addlinespace[1pt]
集合 \(A=\allowbreak \{x\mid\allowbreak (a-1)x^2+\allowbreak 3x-2=\allowbreak 0\}\) 有且仅有两个子集 & 2 & 27 第 6 题、47 第 5 题　\textbf{题干、答案逐字相同}，均得 \(a=\allowbreak 1\) 或 \(a=\allowbreak -\dfrac18\) \\ \addlinespace[1pt]
\(I=\allowbreak \{1,\allowbreak 2,\allowbreak 3,\allowbreak 4,\allowbreak 5\}\) 的「好子集」（\(\lvert A\rvert\leqslant\allowbreak \min(A)\)）计数 & 2 & 13 第 10 题、47 附加 1 题　定义与所求逐字相同，均得 12 \\ \addlinespace[1pt]
\end{xltabular}

\textbf{同模板变体 20 组}（整段相似度 $\geqslant$ 0.71 且最长公共子串 $\geqslant$ 12，且公共串里含汉字／字母／关系符号
这类\textbf{内容字符}------纯数字串不算；成员之间换了数字、问法或设定，\textbf{答案往往随之不同}。
判据本身在本轮修过，见「方法与局限」）：

\begingroup\small
\newcommand{\tabhdrF}{\rowcolor{white} 题目 & 份数 & 分布 & 变在哪 \\}
\rowcolors{1}{white}{rowgray}
\begin{xltabular}{\textwidth}{@{}>{\raggedright\arraybackslash}p{3.63cm}>{\centering\arraybackslash}p{1.36cm}>{\raggedright\arraybackslash}p{1.84cm}>{\raggedright\arraybackslash}p{7.48cm}@{}}
\toprule
\tabhdrF
\midrule
\endfirsthead
\multicolumn{4}{r}{\footnotesize（续上表）} \\
\midrule
\tabhdrF
\midrule
\endhead
\bottomrule
\endlastfoot
含参二次方程与子集讨论\textbf{族内}的交叉（02／05 那一版 $\leftrightarrow$ 07／08／09／11 那一版） & 2 $\leftrightarrow$ 4 & 02 第 17 题、05 第 17 题 $\times$ 07 第 18 题、08 第 18 题、09 第 19 题、11 第 19 题（共 \textbf{7 对}） & 两版\textbf{不是一个组}：\(B=\allowbreak \{x\mid\allowbreak  x^2+\allowbreak 2(a+\allowbreak 1)x+a^2-1=\allowbreak 0\}\) 只在 07／08／09／11 四份一字不差，02／05 那版是 \(a^2-5\)；\(A\) 的方程也是 \(x^2-3x+\allowbreak 2=\allowbreak 0\)／\(x^2+\allowbreak 4x=\allowbreak 0\) 两种。\(r=\allowbreak 0.747\sim0.886\)，都够变体线但够不上逐字线。\textbf{这是全稿最大的一个族}，本稿按相似度拆成逐字档（07／08／09／11）与变体档（02／05）两处，合起来看才是它的全貌 \\ \addlinespace[1pt]
含参区间 \(B\) 与两射线之并 \(A\) 的包含关系 & 2 & 05 第 10 题、08 第 10 题 & 条件相同（都是 \(B\subseteq\allowbreak  A\)），只换端点------\(A\) 左支 \(x<-1\)／\(x<-5\)，右支 \(x>4\)／\(x\geqslant\allowbreak  4\)；\(B\) 左端 \(2a\)／\(a+\allowbreak 1\)，右端同为 \(a+\allowbreak 3\) \\ \addlinespace[1pt]
▲集合 \(M\) 的 \(k\) 个不同子集 \(N_1\dots N_k\) 且并集 \(\neq M\) & 2 & 05 第 16 题 $\times$ 07 第 15 题、11 第 16 题（共 2 对） & 05 版 \(M=\allowbreak \{a,b,c,d,e\}\)、只写「不同子集」（答案 16）；07／11 版 \(M=\allowbreak \{1,\allowbreak 2,\allowbreak 3,\allowbreak 4,\allowbreak 5\}\)、写「不同\textbf{非空}子集」（答案 15）。▲本轮从逐字档移到本档 \\ \addlinespace[1pt]
▲\(A=\allowbreak \{x\mid\allowbreak  x=m^2-n^2,\allowbreak \ m\)、\(n\in\mathbb Z\}\) 的「所有满足 \(A\) 的偶数」 & 2 & 13 第 18 题、15 附加 19 题 & \textbf{本表里重合最强的一对}：13 版是三小问（判断 \(8,\allowbreak 9,\allowbreak 10\) 是否属于 \(A\)／证「\(x\in A\)」是「\(x\in B\)」的必要非充分条件／求所有满足 \(A\) 的偶数），其中第 (3) 问与 15 附加 19 题的\textbf{题干与问法一字不差}；只因判据只看题干（小问被切掉）才落进本档。答案 \(4n,\allowbreak \,n\in\mathbb Z\) \\ \addlinespace[1pt]
▲「数域」定义题（\(a+b\)、\(a-b\)、\(ab\)、\(\frac ab\) 都在集合内） & 2 & 17 第 11 题、20 第 16 题 & 定义句与四个命题的骨架几乎逐字相同，差别在字母（\(F\)／\(G\)）、数字（\(2026\)／\(2023\)）、第三条命题的集合（\(P=\allowbreak \{x\mid\allowbreak  x=\allowbreak 3k\}\)／\(\{x\mid\allowbreak  x=\allowbreak 2k\}\)）与问法（问\textbf{真命题编号}／问\textbf{假命题个数}）；答案（1）（2）（4）／B（1 个） \\ \addlinespace[1pt]
\(M=\allowbreak \{1,\allowbreak \dots,\allowbreak 10\}\) 的非空子集按 \(\sum k(-1)^k\) 求和 & 2 & 09 第 11 题、b 附加 20 题 & 一问填空、一问解答 \\ \addlinespace[1pt]
「全食／偏食」定义题 & 2 & 06 第 11 题、c 第 11 题 & 同一个 \(A=\allowbreak \{-1,\allowbreak \tfrac12,\allowbreak 1\}\)，但 \(B\) 的定义与问法不同（求范围 / 求值） \\ \addlinespace[1pt]
★集合 \(A\) 满足「\(x\)、\(y\in A\Rightarrow x-y\in A\)」（或 \(\tfrac xy\in\mathbb N\)）的元素个数/关系题 & 2 & 03 第 12 题、09 第 7 题 & 两版分别改编自\textbf{同一道} 2012 新课标卷（理）选择 1（见出处表那两条）：09 只把 \(A\) 的元素由 5 个改成 6 个，03 把条件换成 \(\tfrac xy\in\mathbb N\)、\(A\) 换成 \(\{-1,\allowbreak 0,\allowbreak 1\}\)，\(r=\allowbreak 0.734\)。★本对在改判据前是\textbf{漏检}的（公共串全符号、\(r=\allowbreak 0.800\)） \\ \addlinespace[1pt]
▲「\(\tfrac{a_1}{a_2}=\allowbreak \tfrac{b_1}{b_2}=\allowbreak \tfrac{c_1}{c_2}\) 是 \(M=N\) 的什么条件」 & 2 & 17 第 7 题、20 第 15 题 & 20 版\textbf{就是 2003 上海卷（秋理）的那道真题}（解集为两个\textbf{不等式} \(>0\) 的集合，本卷补了 \(\varnothing\subsetneq\allowbreak  M\subsetneq\allowbreak \mathbb R\)、\(\varnothing\subsetneq\allowbreak  N\subsetneq\allowbreak \mathbb R\)），答案「必要非充分」；17 版把不等式改成\textbf{方程} \(=\allowbreak 0\)、去掉两个补条件，答案随之变成「充分非必要」。\textbf{同结构、答案相反} \\ \addlinespace[1pt]
▲\(S=\allowbreak \{s\mid\allowbreak  s=\allowbreak 2n+\allowbreak 1\}\) 与 \(T\) 填包含关系 & 2 & 06 第 7 题、17 第 9 题 & 06 版 \(n\in\mathbb N\)、\(T=\allowbreak \{t\mid\allowbreak  t=\allowbreak 2n-1\}\)、只能填「\(=\allowbreak \)／\(\subset\allowbreak \)／\(\supset\)」，答案 \(\subset\allowbreak \)；17 版 \(n\in\mathbb Z\)、\(T=\allowbreak \{t\mid\allowbreak  t=\allowbreak 4n+\allowbreak 1\}\)、四个符号都能填，解析结论 \(S\supset T\) \\ \addlinespace[1pt]
▲\(P_1=\allowbreak \{x\mid\allowbreak  x^2+ax+\allowbreak 1>0\}\)、\(P_2=\allowbreak \{x\mid\allowbreak  x^2+ax+\allowbreak 2>0\}\) 那道上海春考题 & \textbf{3} & 13 第 14 题、17 第 16 题、\textbf{◇25 第 16 题} & \(r=\allowbreak 0.726\sim0.732\)／公共串 31 字；17／25 两版多出一对 \(Q_1\)、\(Q_2\)（\(x^2+x+b>0\)／\(x^2+\allowbreak 2x+b>0\)），且两版彼此\textbf{逐字相同}（另列逐字档末行），13 版只留 \(P_1\)、\(P_2\) 半题、四个选项改写。原列「判据之外手工补记」，24 份那轮够变体线移入本档，本轮因高一25 入库扩成 3 份 \\ \addlinespace[1pt]
▲「封闭集／单封集」（\(a_i+a_j\) 与 \(a_j-a_i\) 至少有一个属于该集合） & 2 & 05 第 20 题、12 第 21 题 & \(r=\allowbreak 0.726\)／公共串 18 字。原列「判据之外手工补记」，本轮够变体线，移入本档 \\ \addlinespace[1pt]
◇固定 \(B=\allowbreak \{x\mid\allowbreak  x^2-3x+\allowbreak 2=\allowbreak 0\}\) 与含参 \(A\)，\(A\subseteq\allowbreak  B\) 求 \(a\) & 3 & 04 第 15 题、26 第 8 题、43 第 7 题 & 26 与 43 两版的 \(A\)、\(B\) 定义与设问逐字相同（见逐字档）；04 版的含参 \(A\) 不同，但同为讨论 \(A\subseteq\allowbreak  B\) 的范围，故扩为同模板变体组 \\ \addlinespace[1pt]
◇「\(A\) 固定、\(B=\allowbreak \{x\mid\allowbreak  ax+c=\allowbreak 0\}\) 含参，\(B\subseteq\allowbreak  A\) 求 \(a\) 的取值集合」 & \textbf{3} & 15 第 6 题、25 第 6 题、\textbf{◇30 第 8 题} & \(r=\allowbreak 0.719\)／公共串 15 字；\textbf{只共享句型骨架}（\(A\) 由 \(\{1,\allowbreak 2\}\) 换成 \(x^2+x-6=\allowbreak 0\) 的解集 \(\{-3,\allowbreak 2\}\)、\(B\) 由 \(ax=\allowbreak 6\) 换成 \(ax-2=\allowbreak 0\)、条件由 \(A\cup\allowbreak  B=A\) 换成 \(B\subseteq\allowbreak  A\)），与末两组同型，证据强度弱于前面几组，单列备查。\textbf{◇本轮因高一30 入库扩成 3 份}：30 第 8 题与 15 第 6 题是逐字同题（另列逐字档），只与 25 第 6 题够不上逐字线 \\ \addlinespace[1pt]
三集合两两交集各为单元素且三重交集为空 & 2 & 39 附加 22 题、42 第 11 题 & 两题都数三集合交集结构，但底集由 \(\{1,\allowbreak 2,\allowbreak 3\}\) 扩为 \(\{1,\allowbreak 2,\allowbreak 3,\allowbreak 4\}\)，问法与答案分别不同（6／96），属同模板变式，不算逐字同题 \\ \addlinespace[1pt]
正数 \(a,b\) 满足 \(a+b=\allowbreak 1\)，按给定解法最小化倒数线性组合 & 2 & 37 第 20 题、43 第 20 题 & 给定解法相同，目标式分别为 \(\dfrac1a+\allowbreak \dfrac2b\) 与 \(\dfrac1a+\allowbreak \dfrac1b\)，系数不同，属同模板变式 \\ \addlinespace[1pt]
\(p\) 与 \(q\) 对应区间的充分不必要条件 & 2 & c 第 2 题、43 第 3 题 & 同以区间真包含转化为充要关系求参数，区间端点与答案范围不同 \\ \addlinespace[1pt]
正参数下两个闭区间的倒数映射不变性 \(\lambda/a\in A\) & 2 & 30 第 12 题、45 第 10 题 & 都求正数 \(\lambda\) 使 \(a\in A\Rightarrow\lambda/a\in A\)，但两段区间端点与参数条件由 \([t,t+\allowbreak 1]\cup\allowbreak [t+\allowbreak 4,t+\allowbreak 9],\allowbreak \ 0\notin A\) 改为 \([t,t+\allowbreak 1]\cup\allowbreak [t+\allowbreak 3,t+\allowbreak 6],\allowbreak \ t>0\)；判为同模板变式，不把新卷题当作该高考题的另一个对应 \\ \addlinespace[1pt]
参数化二次式的存在命题为假，求参数范围 & 2 & c 第 4 题、47 第 4 题 & 都以二次函数／方程的判别式判断参数化的存在命题为假；表达式、真假条件与答案范围不同（机判 LCS 12、\(r=\allowbreak 0.714\)） \\ \addlinespace[1pt]
用 \(C(A)\) 定义集合基数差 \(A*B\)，求使 \(A*B=\allowbreak 1\) 的参数个数 & 2 & 07 第 12 题、47 第 10 题 & 都由 \(C(A)\) 与 \(A*B=\allowbreak \lvert C(A)-C(B)\rvert\) 引入基数差，求满足条件的参数取值个数；多项式构造不同，均答 5，按完整题目判同模板变体（机判将其列入逐字候选） \\ \addlinespace[1pt]
\end{xltabular}
\endgroup

\textbf{有 3 行从本表移出、改列「判据之外手工补记」}------它们共享的只是句型骨架（「含参集合 + 包含
关系 + 求 \(a\) 范围」），集合与条件都不相同，两版的\textbf{公共串短于 12 字}，任何文本判据都捞不到：
\texttt{08\ 第\ 17\ 题\ ×\ 09\ 第\ 17\ 题}（公共串 5 字）、\texttt{02\ 第\ 18\ 题\ ×\ 03\ 第\ 16\ 题}（4 字）、
\texttt{05\ 第\ 7\ 题\ ×\ c\ 第\ 5\ 题}（9 字）。它们不是被新判据筛掉的，而是\textbf{从来就没进过候选池}。

\textbf{另有一个族跨在「逐字档」与「判据之外」之间}------\textbf{「性质 \(P\)／封闭集／单封集」}：三者是
同一副定义的三种说法（对任意 \(i\leqslant j\)，\(a_ia_j\) 与 \(a_j/a_i\)／\(a_i+a_j\) 与 \(a_j-a_i\)
两数中至少有一个属于该集合），源头是 \textbf{2009 北京卷（理）的「性质 \(P\)」}（乘除版）。本库里
高一05 第 19 题是它的乘除版（换名「性质 \(M\)」），高一05 第 20 题「封闭集」与高一12 第 21 题
「单封集」是它的\textbf{加减版}，两版定义逐字程度很高（见「判据之外手工补记」与出处表）。

\textbf{全稿最大的一个族藏在上面两档之间}：「含参二次方程 + 子集讨论」共出现在 \textbf{6 份}里------
高一02 第 17 题、05 第 17 题、07 第 18 题、08 第 18 题、09 第 19 题、11 第 19 题，其中
\(B=\{x\mid x^2+2(a+1)x+a^2-1=0\}\) 这一表达式\textbf{四份一字不差}（02／05 那份是 \(a^2-5\)），
差别只在 \(A\) 的方程（\(x^2+4x=0\) 或 \(x^2-3x+2=0\)）与小问。本稿按相似度把它拆在了逐字档
（07／08／11）、变体档（02／05）和单成员变体（09）三处，合起来看才是它的全貌。另有一份
（高一04 第 15 题）是它的\textbf{反向形式}（含参的换成了 \(A\)、固定的换成 \(B\)）。

后一档里最后两组\textbf{只共享句型骨架}（「含参集合 + 包含关系 + 求 \(a\) 范围」），集合与条件都不
相同，证据强度明显弱于前五组，单列备查；前五组的题设骨架则一致，只是参数、问法或小问有出入。

\textbf{判据之外手工补记 8 对（7 行）}------这一批的分数都在变体线以下，但逐条读过原文、确实是真同题。它们正好
说明 \textbf{「骨架文本相似度」量的不是「数学内容」}：下文「方法与局限」里 \texttt{高一07\ 第\ 4\ 题\ ×\ 高一b\ 第\ 8\ 题} 的 \(r=0.783\) 是虚高（两句都短），下面七对则都是虚低（题面文字改了、数学没改）。
（★ 为 2026-09-15 重跑新增；▲ 为 2026-09-15 末并入高一19---21 那一轮新增；分数是本轮口径，
与上一轮有 $\pm$0.02 的出入。）

\begingroup\small
\newcommand{\tabhdrG}{\rowcolor{white} 题目 & 份数 & 分布 & \(r\)／LCS \\}
\rowcolors{1}{white}{rowgray}
\begin{xltabular}{\textwidth}{@{}>{\raggedright\arraybackslash}p{6.78cm}>{\centering\arraybackslash}p{1.10cm}>{\raggedright\arraybackslash}p{4.68cm}>{\centering\arraybackslash}p{1.75cm}@{}}
\toprule
\tabhdrG
\midrule
\endfirsthead
\multicolumn{4}{r}{\footnotesize（续上表）} \\
\midrule
\tabhdrG
\midrule
\endhead
\bottomrule
\endlastfoot
\(A=\allowbreak \{x\mid\allowbreak  x^2-3x-4=\allowbreak 0\}\)、\(B=\allowbreak \{x\mid\allowbreak  ax-1=\allowbreak 0\}\) & 2 & 05 第 6 题、c 第 3 题 & 0.582／14 \\ \addlinespace[1pt]
\(A=\allowbreak \{x\mid\allowbreak  x<-1\text{ 或 }x>4\}\) 与含参 \(B\) & 2 & 05 第 10 题、07 第 17 题 & 0.575／13 \\ \addlinespace[1pt]
固定方程 \(x^2+\allowbreak 2x-8=\allowbreak 0\)、\(x^2-5x+\allowbreak 6=\allowbreak 0\) ＋ 含参 \(x^2-tx+t^2-c=\allowbreak 0\) & 3 & 08 第 6 题、a 第 14 题、15 第 14 题 & 0.448／13（a 与 15 两版已升入逐字档，\textbf{本行指 08 版与这两份的关系}，共 2 对） \\ \addlinespace[1pt]
★「第 \(k\) 个子集」（2 的幂编码） & 2 & 08 第 11 题、a 附加题第 2 题 & 0.614／10　\textbf{同源}：都改编自 2010 湖南卷（文）填空 7，改法不同（见表下说明） \\ \addlinespace[1pt]
▲同一个 \(P=\allowbreak (-2,\allowbreak 3)\) 与含参 \(Q\) 的包含关系 & 2 & 08 第 17 题、09 第 17 题 & 0.640／\textbf{5} \\ \addlinespace[1pt]
▲含参区间与包含关系（含参的那个一在 \(B\)、一在 \(A\)） & 2 & 02 第 18 题、03 第 16 题 & 0.448／\textbf{4} \\ \addlinespace[1pt]
▲\(A\) 为固定区间、\(B\) 含参时的包含关系 & 2 & 05 第 7 题、c 第 5 题 & 0.635／\textbf{9} \\ \addlinespace[1pt]
\end{xltabular}
\endgroup

后 3 对（▲）是\textbf{从变体档移下来的}，缘由见变体表下的那段说明：它们的公共串只有 4---9 字，
短于候选取线（LCS $\geqslant$ 12），\textbf{根本进不了候选池}------不是判据放宽或收紧能救回来的。
上一轮把它们列在变体表里，靠的是「只共享句型骨架」这句话，与判据口径不符，本轮据实归位。
另外 3 对（「特征数」、上海春考的 \(P_1/P_2\)、「封闭集／单封集」）本轮\textbf{分数够了}，
已分别移入逐字档与变体档。

\begin{itemize}
\tightlist
\item
  \textbf{「特征数」}（原列本表首行，本轮\textbf{移入逐字档}）：同一个自定义概念------把 \(M\) 分成三个集合、
  \(X_i\) 都取「最大元素与最小元素\textbf{之和}」。差别只在 \(M\)（\(1\sim12\)／\(0<x\leqslant15\)）、
  每组元素个数（4／5）与问法（求最大值与最小值的积 \(=1485\)／问哪个值不可能）。判据这次能捞到它，
  是因为\textbf{公共串 27 字全落在「特征数」那句定义上}，越过了逐字档的「$\geqslant$ 25 字」线。这个名词与
  高考里的「特征数」\textbf{同名不同义}（见「考据不出来的」一节）。
\item
  \textbf{\(x^2-3x-4=0\) 那一对}：\textbf{两版集合一字不差}（与变体档首行那对一样强），只是 05 问
  \(A\cap B=B\)、12 问「\(x\in B\)」是「\(x\in A\)」的充分非必要条件------问法占了句子一大截，
  把分数拉了下去。
\item
  \textbf{\(A=\{x\mid x<-1\text{ 或 }x>4\}\) 那一对}：\(A\) 一字不差，\(B\) 都是含参区间、都求参数范围。
  \textbf{据此上表「含参区间 \(B\) 与两射线之并 \(A\)」那一族实为 3 份}（05／07／08），不只是 2 份。
\item
  \textbf{最后一对是这一轮新翻出来的}，也是四对里唯一\textbf{连 \(r\ge 0.50\) 的候选池都没进}的：08 第 6 题
  给 \(A=\{x\mid x^2+2x-8=0\}\)、\(B=\{x\mid x^2-5x+6=0\}\)、\(C=\{x\mid x^2-mx+m^2-13=0\}\)，
  问「\(B\cap C\neq\varnothing\) 且 \(A\cap C=\varnothing\) 求 \(m\)」；a 第 14 题把三条式子摆成
  \(A=\{x\mid x^2-ax+a^2-19=0\}\)、\(B=\{x\mid x^2-5x+6=0\}\)、\(C=\{x\mid x^2+2x-8=0\}\)，
  三小问全用 \(\cap\)、\(\subseteq\) 与 \(\varnothing\) 串起来。\textbf{两条固定方程一字不差}，第三条同为
  \(x^2-tx+t^2-c=0\) 型，只是含参的那条在两题里换了位置。
  \textbf{前两档按「骨架文本」找，它按「共享方程」才找得出来}------题面措辞改得多，\(r\) 就掉到 0.448，
  但数学骨架是同一副。
\end{itemize}

（上面的 \(r\)／LCS 都是\textbf{本轮 29 份口径}；下面这段阈值分析里的数字是 \textbf{12 份}那一轮跑的，
22 份口径下未重算，结论不变。）

这 4 对\textbf{不能靠降线收进来}：后两对里最松的是 0.597，要收它得把变体线从 0.71 压到 0.597，而
\(r\) 落在 0.597--0.71、LCS $\geqslant$ 12 的共 15 对，真同题只有 2 对，其余 13 对全是套话；最后一对
（0.448）更在 0.50 之外，再降就把 646 对翻成上千对。故按「分数不够、内容够」单列在此。

同一轮扫描里\textbf{更弱的几对没有收}------它们只共用一段常见表达式，不足以断定同源，故留在此处备查
不列成组：\texttt{05\ 第\ 6\ 题／06\ 第\ 3\ 题／c\ 第\ 3\ 题} 三者共用 \(B=\{x\mid ax-1=0\}\)（其中 05 与 c
另共用 \(A\) 一字不差，已收在上表；\texttt{06} 只共用 \(B\)。2026-09-15 补跑又发现 \textbf{05 与 10 的那一版
逐字相同}，一并收在上表末两行）；\texttt{05\ 第\ 7\ 题\ ×\ c\ 第\ 15\ 题} 共用
\(B=\{x\mid a-1\leqslant x\leqslant a+2\}\)；\texttt{02／05\ 第\ 17\ 题\ ×\ 08\ 第\ 3\ 题} 共用
\(A=\{x\mid x^2-3x+2=0\}=\{1,2\}\)------这个集合太常见，几乎不构成证据。

不同学校、不同周次出现这么多同题，说明各校备课组\textbf{大量取材于同一批公共题源}（教辅、网络
题库、区统考、教材习题），而非各校独立原创。旁证是这些题在菁优网、1010jiajiao、百度题库、
作业帮上都能搜到近乎一致的版本。

\clearpage
\section{能考据到的高考出处：43 道真题（50 条对应／58 个题位）}

拿 1952--2026 年的\textbf{全部高考真题}逐字比对（参考库 \texttt{\textasciitilde{}/\allowbreak{}Sources/\allowbreak{}AI/\allowbreak{}Gaokao-\allowbreak{}Math-\allowbreak{}Problems-\allowbreak{}Compilation}，
含上海春/秋/文/理及各地方卷；去噪后有效题干 \textbf{16,273 道}），能定到具体试卷的\textbf{真题有 43 道}，
对应本库 \textbf{58 个题位}、\textbf{50 条对应关系}：

\begin{itemize}
\tightlist
\item
  \textbf{按关系分}：原题照录 \textbf{22} 条、改编 \textbf{25} 条、借定义／借集合另出题 \textbf{3} 条 ＝ 50 条;
\item
  \textbf{真题侧}：\textbf{3 道}真题被\textbf{三份}试卷用到（2015 上海卷（春）：高一13／17／25；
  2003 上海卷（秋理）：高一17／20／36；2010 湖南卷（文）填空 7：高一08／高一a／高一39），
  \textbf{9 道}各被两份用到（2009 北京卷（文）填空 6「孤立元」、2017 江苏卷、2012 新课标卷（理）、
  2012 江苏卷 \(f(n)\) 计数题、1999 广东卷／全国卷、2008 山东卷（文/理）、2009 北京卷（理）「性质 \(P\)」、2009 天津卷（理）、2019 上海卷（春））
  $\to$ 3 $\times$ 3 ＋ 9 $\times$ 2 ＋ 31 $\times$ 1 ＝ \textbf{58 个题位}（另有 31 道各被一份用到）；\textbf{7 条}对应关系是同一道真题的第二种／第三种改法
  （2012 新课标卷（理）选择 1 两种、2012 江苏卷两种、2015 上海卷（春）三种、2003 上海卷（秋理）两种、
  2009 天津卷（理）两种、2019 上海卷（春）两种）$\to$ 6 道真题共多出 1 ＋ 1 ＋ 2 ＋ 1 ＋ 1 ＋ 1 ＝ \textbf{7 条}，故 43 ＋ 7 ＝ \textbf{50 条对应}。
\end{itemize}

判定口径是「脱去 LaTeX 与标点后的骨架文本最长公共子串」+ 整段相似度，下表每一条都逐字看过原文：
（★ 为 2026-09-15 全量重跑新增，含旧卷漏检 2 条；▲ 为并入高一19---21 后新增；● 为 2026-09-16
并入高一22 后新增------那一轮三条\textbf{全是原题照录}；◆ 为 2026-09-20 并入高一23／24 后新增------
那一轮一条，系\textbf{改编}；◇ 为 2026-09-22 并入高一25---28 后新增------那一轮两条，\textbf{全是原题照录}。
\textbf{■} 为 2026-09-28 顺核①机判命中清单时补登的两条------清单里另外 17 条判为不算，逐条见下。）

\begin{quote}
\textbf{本节原与「核对摘要」那一行的口径不同}（摘要写「26 道／31 个题位」，本节写「24 道／27 个题位」）。
根子是\textbf{两处数的不是一个东西}：摘要数的是「题目 $\leftrightarrow$ 真题」的\textbf{对应条数}，本节数的是\textbf{真题的道数}，
而两处都没把这层区分写出来；本节那个「27 个题位」还漏算了 5 个。\textbf{2026-09-15 已统一}------
两个地方改用同一套数（本节标题、正文与摘要那一行逐处对齐）。
\textbf{2026-09-28 又发现}上一轮（并入高一37---高一40）\textbf{只加了表行、没改本节标题}，于是本节停在
「36 道／41 条／45 个题位」而摘要已走到「39／44／48」；而且\textbf{两处的题位数都对不上那张表}
------把表逐行数出来是 \textbf{50 个题位}（摘要写 48、本节写 45，两处自己的算式也各差 1），
摘要的「6 道各被两份用到」也漏了 1 道（2003 上海卷（秋理）其实是\textbf{三份}在用）。这一轮一并改正，
并按那张表给这几个数立了断言（见 \texttt{工具/\allowbreak{}计数核对.py}）。数「道」时要把\textbf{同一份卷里的两道不同真题
分开算}（2009 北京卷（文）的「填空 6 孤立元」与「选择 8」是两道），否则会少算 1。

另：本节早年那几轮的\textbf{题位增量也记少过}------并入高一31---高一36 那一轮实际新增 \textbf{7 个题位}
（高一31 第 14、高一34 第 8、高一35 第 12、高一35 附加 17、高一36 第 3、高一36 第 13、
高一36 第 14），当时却记成「40 $\to$ 45」（+5），此后一路少 2（它写的 45／48 应为 \textbf{47／50}）；
并入高一37---高一40 那一轮的 +3 是对的。上表与本节标题现按\textbf{逐行数出来的真值}写。

另需说明：本节对\textbf{后入库各卷}的并入并不齐------\textbf{高一19---21 只并入了下面这一条}：
\textbf{高一20 第 15 题＝2003 上海卷（秋理）}（骨架 LCS 55 字、\(r=0.935\)，本卷补了
\(\varnothing\subsetneq M\subsetneq\mathbb R\)、\(\varnothing\subsetneq N\subsetneq\mathbb R\) 两个条件），
它和 17 第 7 题那一行是\textbf{同一道真题的两种改法}；\textbf{高一22 的三条则全部入了表}（表末三行，
2026-09-16 补）；\textbf{高一23 的一条也已入表}（◆那一行，2026-09-20 补）；\textbf{高一25／高一28 的两条
也已入表}（◇表末两行，2026-09-22 补，两条都是原题照录）。这几份新卷的其余题在真题库里
\textbf{都没有达逐字线的
命中}（19 第 3 题、20 第 3／8 题只有 12---22 字的通用片段，属「同题型但非同一题」；高一25---28
机判另给出 4 条候选，回原文复核后判为假命中，见「核对摘要」的注）。

其后各轮的新真题与新题位\textbf{都已入表}：并入高一30 的两条（◇ 两行）、并入高一31---高一36 的
\textbf{5 道新真题 ＋ 高一36 第 13 题那个新题位}、并入高一37---高一40 的 \textbf{3 条}（表末三行）；
并入高一29 的那一轮整批零命中，\textbf{并入高一41 的 2026-09-28 这一轮同样零命中}。
\end{quote}

\newcommand{\tabhdrH}{\rowcolor{white} 本库题目 & 出处 & 关系 \\}
\rowcolors{1}{white}{rowgray}
\begin{xltabular}{\textwidth}{@{}>{\raggedright\arraybackslash}p{2.69cm}>{\raggedright\arraybackslash}p{4.20cm}>{\raggedright\arraybackslash}p{7.85cm}@{}}
\toprule
\tabhdrH
\midrule
\endfirsthead
\multicolumn{3}{r}{\footnotesize（续上表）} \\
\midrule
\tabhdrH
\midrule
\endhead
\bottomrule
\endlastfoot
高一06 第 9 题、★高一15 附加 17 题 & \textbf{2009 北京卷（文）填空 6}「孤立元」 & \textbf{原题照录}（高一15 版最长公共子串 76 字，整题一字不差；两版互为逐字同题） \\ \addlinespace[1pt]
高一09 第 5 题 ★（旧卷漏检） & \textbf{2008 浙江卷（理）选择 1} & \textbf{原题照录}（\(U=\allowbreak \mathbb{R}\)、\(A=\allowbreak \{x\mid\allowbreak  x>0\}\)、\(B=\allowbreak \{x\mid\allowbreak  x\leqslant\allowbreak -1\}\) 与所求式 \((A\cap\allowbreak \overline{B})\cup\allowbreak (B\cap\allowbreak \overline{A})\) 一字不差，只去掉四个选项改成填空） \\ \addlinespace[1pt]
★高一17 第 15 题 & \textbf{2012 湖北卷（文）} & \textbf{原题照录}（「存在一个无理数，它的平方是有理数」的否定；题干与四个选项逐字相同，\(r=\allowbreak 0.977\)） \\ \addlinespace[1pt]
高一b 第 9 题 & \textbf{2014 福建卷（文）填空 4} & \textbf{原题照录}（\(\{a,b,c\}=\allowbreak \{0,\allowbreak 1,\allowbreak 2\}\) 且三条关系只有一条成立，求 \(100a+\allowbreak 10b+c\)） \\ \addlinespace[1pt]
高一b 第 13 题 & \textbf{2018 全国 II 卷（理）选择 2} & \textbf{原题照录}（\(A=\allowbreak \{(x,y)\mid\allowbreak  x^2+y^2\leqslant\allowbreak  3,\allowbreak \ x,y\in\mathbb{Z}\}\) 求元素个数） \\ \addlinespace[1pt]
高一b 第 15 题 & \textbf{2009 北京卷（文）选择 8} & \textbf{原题照录}（正三角形内 \(\lvert PO\rvert\leqslant\allowbreak \lvert PO_i\rvert\) 的区域，只把 T/D、O/P$_{0}$ 换了记号） \\ \addlinespace[1pt]
高一08 第 15 题、高一32 第 13 题 & \textbf{1999 广东卷 / 全国卷 选择 1} & \textbf{原题照录}（Venn 图阴影区域；两版的题干、四个选项\textbf{与图形}均逐字相同，答案同为 C------见同题表那一行） \\ \addlinespace[1pt]
高一c 第 13 题 & \textbf{2011 湖南卷（理）选择 2} & \textbf{原题照录}（\(M=\allowbreak \{1,\allowbreak 2\},\allowbreak \ N=\allowbreak \{a^2\}\)，问「\(a=\allowbreak 1\)」是「\(N\subseteq\allowbreak  M\)」的什么条件） \\ \addlinespace[1pt]
高一a 第 11 题 & \textbf{2013 山东卷（理）选择 2} & \textbf{原题照录}（\(A=\allowbreak \{0,\allowbreak 1,\allowbreak 2\}\)，\(B=\allowbreak \{x-y\mid\allowbreak  x,y\in A\}\)；题干一字不差，只把四个选项的次序倒了过来） \\ \addlinespace[1pt]
高一03 第 5 题、★高一09 第 15 题 & 2017 江苏卷 填空 1 & 改编·\textbf{两种改法}：真题是 \(A=\allowbreak \{1,\allowbreak 2\}\)、\(B=\allowbreak \{a,a^2+\allowbreak 3\}\)、\(A\cap\allowbreak  B=\allowbreak \{1\}\) 求 \(a\)（答 1）；高一03 版把 \(a\) 换成 \(-a\) 并改问 \(A\cup\allowbreak  B=\allowbreak \{1,\allowbreak 2,\allowbreak 4\}\)，高一09 版把 \(a^2+\allowbreak 3\) 换成 \(a^2+\allowbreak 1\)、改成选择（答案随之由 1 变 0） \\ \addlinespace[1pt]
高一08 第 13 题 & 2004 全国 IV 卷（理）选择 1 & 改编（\(M\)、\(N\) 一字未动，把 \(M\cap\allowbreak  N\) 改成 \(M\cup\allowbreak  N\) 并换掉四个选项） \\ \addlinespace[1pt]
高一09 第 7 题 & 2012 新课标卷（理）选择 1 & 改编（\(A\) 的元素由 \(\{1,\allowbreak 2,\allowbreak 3,\allowbreak 4,\allowbreak 5\}\) 改为 \(\{1,\allowbreak 2,\allowbreak 3,\allowbreak 4,\allowbreak 5,\allowbreak 6\}\)） \\ \addlinespace[1pt]
高一03 第 12 题 & 2012 新课标卷（理）选择 1 & 改编（与上一条同源、改法不同：把条件 \(x-y\in A\) 换成 \(\tfrac xy\in\mathbb N\)，\(A\) 换成 \(\{-1,\allowbreak 0,\allowbreak 1\}\)） \\ \addlinespace[1pt]
高一08 第 11 题、★高一a 附加题第 2 题、高一39 附加 21 题 & 2010 湖南卷（文）填空 7 & 改编·\textbf{两种改法}：\textbf{高一a 版的定义与真题一字不差}（\(k=\allowbreak 2^{i_1-1}+\allowbreak 2^{i_2-1}+\allowbreak \cdots\)、集合 \(M=\allowbreak \{a_1,\allowbreak \dots,a_n\}\)），只把问法换成「第 25 个子集」；高一08 版把指数改成元素值（集合取 \(\{0,\allowbreak 1,\allowbreak \dots,n\}\)、\(k=\allowbreak 2^{a_1}+\allowbreak \cdots\)；原卷此题上标用 \(n\)、子集下标却用 \(m\)，本身不统一，本稿照原卷），保留「第 211 个子集」这一问、删去原题第一问 \\ \addlinespace[1pt]
高一a 第 4 题、高一b 第 7 题 & 2008 山东卷（\textbf{文/理}）选择 1 & 改编（原题 \(M\subseteq\allowbreak \{a_1,\allowbreak \dots,a_4\}\)，本稿扩成 5 个元素，答案由 2 变 4）。\textbf{该卷文、理两卷都印了这一道}，两边相似度相同，故出处写作「文/理」 \\ \addlinespace[1pt]
★高一17 第 16 题 & 2015 上海卷（春） & 改编（\(P_1,P_2,Q_1,Q_2\) 四集合版；\textbf{C、D 两项真题作「存在 \(a\)」，原卷印成「对任意 \(a\)」}------原卷笔误，答案 B 不受影响） \\ \addlinespace[1pt]
★高一13 第 14 题 & 2015 上海卷（春） & 改编（与上一行同源：只留 \(P_1,P_2\) 半题，四个选项改写） \\ \addlinespace[1pt]
★高一17 第 7 题、\textbf{▲高一20 第 15 题} & 2003 上海卷（秋理） & \textbf{同一道真题的两种改法}：17 版把真题的「两个一元二次\textbf{不等式}解集相同」改成「两个一元二次\textbf{方程}解集相同」的填空题（\(r=\allowbreak 0.882\)）；20 版\textbf{保留不等式}（\(r=\allowbreak 0.935\)、LCS 55 字），只补了 \(\varnothing\subsetneq\allowbreak  M\subsetneq\allowbreak \mathbb R\)、\(\varnothing\subsetneq\allowbreak  N\subsetneq\allowbreak \mathbb R\) 两个条件------\textbf{接近于原题照录}。两版答案不同（充分非必要／必要非充分） \\ \addlinespace[1pt]
★高一15 第 9 题 & 2006 山东卷（理/文） & 改编（\(A\odot B=\allowbreak \{z\mid\allowbreak  z=xy(x+y),\allowbreak \ x\in A,\allowbreak \ y\in B\}\) 定义一字不差；\(A,B\) 由 \(\{0,\allowbreak 1\},\allowbreak \{2,\allowbreak 3\}\) 换成 \(\{1,\allowbreak 2\},\allowbreak \{3,\allowbreak 4\}\)，选择改填空） \\ \addlinespace[1pt]
★高一17 第 9 题 & 2021 全国乙卷（理）选择 2 & 借同一对集合另问（\(S=\allowbreak \{s\mid\allowbreak  s=\allowbreak 2n+\allowbreak 1,n\in\mathbb{Z}\}\)、\(T=\allowbreak \{t\mid\allowbreak  t=\allowbreak 4n+\allowbreak 1,n\in\mathbb{Z}\}\) \textbf{定义逐字相同}；真题问 \(S\cap\allowbreak  T\)，本卷改问 \(S\) 与 \(T\) 的包含关系） \\ \addlinespace[1pt]
★高一12 第 21 题、★高一05 第 19 题 & 2009 北京卷（理）「性质 \(P\)」 & 同构改编／改编：真题定义为「\(a_ia_j\) 与 \(a_j/a_i\) 至少有一个属于 \(A\)」------高一05 第 19 题是它的\textbf{乘除版}（换名「性质 \(M\)」，改小问），高一12 第 21 题「单封集」是它的\textbf{加减版}（\(a_i+a_j\) 与 \(a_j-a_i\)） \\ \addlinespace[1pt]
高一a 第 2 题 & 2012 辽宁卷（理/文）选择 & 改编（\(U\) 由 \(\{0,\allowbreak \dots,\allowbreak 9\}\) 改为 \(\{0,\allowbreak \dots,\allowbreak 8\}\)；问法未变，答案随之由 \(\{7,\allowbreak 9\}\) 变为 \(\{7\}\)） \\ \addlinespace[1pt]
高一b 第 11 题 & 2010 上海卷（秋理）填空 14 & 改编（4 个子集 $\to$ 2 个子集，\(\varnothing\)、\(U\) $\to$ \(a\)、\(b\)） \\ \addlinespace[1pt]
高一c 第 10 题 & 2015 浙江卷（理）选择 6 & 借 \(d(A,B)=\allowbreak \operatorname{card}(A\cup\allowbreak  B)-\operatorname{card}(A\cap\allowbreak  B)\) 的定义另出题 \\ \addlinespace[1pt]
★高一02 第 15 题 & 2012 江苏卷（\(f(n)\) 计数题） & 改编（\textbf{同构换参数}：真题是 \(P_n=\allowbreak \{1,\allowbreak \dots,n\}\)、\(A\subseteq\allowbreak  P_n\)、\(x\in A\Rightarrow2x\notin A\) 且 \(x\in P_n\setminus\allowbreak  A\Rightarrow2x\notin P_n\setminus\allowbreak  A\)，求 \(f(4)\) 与 \(f(n)\)；本卷改成 \(U=\allowbreak \{1,\allowbreak \dots,\allowbreak 6\}\)、\(2x\to\allowbreak 3x\)、改问选择的元素个数------答案由 \(f(4)=\allowbreak 4\) 变为 \(16\)） \\ \addlinespace[1pt]
★高一17 第 11 题 & 2008 福建卷（文/理）「数域」 & 改编（\textbf{同一定义、四个命题全换}：真题文理两卷各给一组「数域必含 \(0,\allowbreak 1\)／整数集是数域／\(\mathbb{Q}\subseteq\allowbreak  M\) 则 \(M\) 必为数域／数域必为无限集」式命题，本卷换成「\(0\) 是任何数域的元素／有非零元素则 \(2026\in F\)／\(\{3k\}\) 是否数域／\(\mathbb{Q}\) 是否数域」） \\ \addlinespace[1pt]
高一42 第 20 题 & 2012 江苏卷（\(f(n)\) 计数题） & 改编：\(P_n\) 与「禁选相邻倍数」的计数结构、求 \(f(4)\) 与 \(f(n)\) 均同源，但条件③把真题的相对补集写成「若 \(x\notin A\) 则 \(2x\in A\)」。按此字面读法，\(2x>n\) 时与 \(f(4)=\allowbreak 4\) 的印刷结论矛盾；本卷解析按链内关系计数，详见本卷专记。 \\ \addlinespace[1pt]
高一43 第 18 题 & \textbf{1997 年全国卷（文/理）解答题} & \textbf{改编}：同一高速运输成本模型，固定距离、速度区间与费用参数后求最小全程成本，另加「每小时成本不高于指定值」的小问；真题文、理版本是同一道题，故合并记一条对应。 \\ \addlinespace[1pt]
●高一22 第 11 题 & \textbf{2011 年大纲卷（文/理）选择} & \textbf{原题照录}（题干「下面四个条件中，使 \(a>b\) 成立的充分而不必要的条件是」与四个选项 \(a>b+\allowbreak 1\)／\(a>b-1\)／\(a^2>b^2\)／\(a^3>b^3\) \textbf{逐字相同}，\(r=\allowbreak 1.000\)；答案同为 A。\textbf{该卷文、理两卷都印了这一道}，两边相似度相同，故出处写作「文/理」） \\ \addlinespace[1pt]
●高一22 第 13 题 & \textbf{2010 年辽宁卷（理）选择} & \textbf{原题照录}（题干「已知 \(a>0\)，则 \(x_0\) 满足关于 \(x\) 的方程 \(ax=b\) 的充要条件是」与四个选项------\(\exists x\in\mathbb R\)／\(\forall x\in\mathbb R\) 配 \(\geqslant\allowbreak \)／\(\leqslant\allowbreak \) 两种，见本卷源码------\textbf{逐字相同}，\(r=\allowbreak 1.000\)；答案同为 C） \\ \addlinespace[1pt]
●高一22 附加 19 题 & \textbf{2006 年湖北卷（理）选择} & \textbf{原题照录}（四条命题与真题相同、答案同为①②；差别只有两处表述------第 ③ 条真题印 \(\nsubseteq\) 的符号写法、本卷写成文字「\(A\) 不是 \(B\) 的子集」（同义），以及由选择改成了填空。\(r=\allowbreak 0.952\)、最长公共子串 66 字） \\ \addlinespace[1pt]
◆高一23 第 12 题 & \textbf{2004 年北京卷（秋文/秋理）选择} & \textbf{改编}（借真题题干「已知 \(a,b,c\) 满足 \(c<b<a\) 且 \(ac<0\)」与 A、B 两个选项，把问法由「一定成立的是」\textbf{反转}为「不成立的是」，C 由 \(cb^2<ab^2\) 改为 \(c(b^2+\allowbreak 1)<a(b^2+\allowbreak 1)\)、D 由 \(ac(a-c)>0\) 改为 \(ac(a-c)<0\)，答案由 A 变 B。真题的 B 项 \(c(b-a)<0\) 恰是「不成立」的那一个，故本卷答案正是真题里不成立的选项。\textbf{该卷文、理两卷都印了这一道}，故出处写作「文/理」；机判 \(r=\allowbreak 0.742\)、LCS 13，落变体档，回原文复核确认同源） \\ \addlinespace[1pt]
◇高一28 第 16 题 & \textbf{2013 年浙江卷（文）选择} & \textbf{原题照录}（两个自定义运算 \(a\wedge b=\allowbreak \min\{a,b\}\)、\(a\vee b=\allowbreak \max\{a,b\}\) 的定义式，四个选项 \(a\wedge b\geqslant\allowbreak 2\)／\(c\wedge d\leqslant\allowbreak 2\) 等，与答案 \textbf{C} 逐字相同；\(r=\allowbreak 0.966\)、最长公共子串 19 字） \\ \addlinespace[1pt]
◇高一25 第 16 题 & \textbf{2015 年上海卷（春）选择} & \textbf{原题照录}（\(P_1\)／\(P_2\)／\(Q_1\)／\(Q_2\) 四集合版，题干、四个选项与答案 \textbf{B} 逐字相同；\(r=\allowbreak 0.916\)）。与下面 高一17 第 16 题那一行是\textbf{同一道真题的第三种改法}------17 那版的 C、D 两项把真题的「存在 \(a\)」印成「对任意 \(a\)」（原卷笔误），25 这版与真题一致 \\ \addlinespace[1pt]
◇高一30 第 4 题 & \textbf{2010 年安徽卷（秋文）填空} & \textbf{原题照录}（命题「存在 \(x\in\R\)，使得 \(x^2+\allowbreak 2x+\allowbreak 5=\allowbreak 0\)」的否定；\(r=\allowbreak 0.909\)、最长公共子串 18 字，答案同为「对任意 \(x\in\R\)，都有 \(x^2+\allowbreak 2x+\allowbreak 5\neq0\)」） \\ \addlinespace[1pt]
◇高一30 第 12 题 & \textbf{2019 年上海卷（春）填空} & \textbf{原题照录}（\(A=\allowbreak [t,t+\allowbreak 1]\cup\allowbreak [t+\allowbreak 4,t+\allowbreak 9]\)、\(0\notin A\)、存在正数 \(\lambda\) 使 \(\dfrac\lambda a\in A\)；\textbf{整题一字不差}，\(r=\allowbreak 1.000\)、最长公共子串 47 字，答案同为 \(t=\allowbreak 1\) 或 \(-3\)） \\ \addlinespace[1pt]
高一45 第 10 题 & \textbf{2019 年上海卷（春）填空} & \textbf{改编}：同为正区间集合在倒数映射 \(a\mapsto\lambda/a\) 下不变的问题；两段区间端点及条件由 \([t,t+\allowbreak 1]\cup\allowbreak [t+\allowbreak 4,t+\allowbreak 9]\)、\(0\notin A\) 改为 \([t,t+\allowbreak 1]\cup\allowbreak [t+\allowbreak 3,t+\allowbreak 6]\)、\(t>0\)（机判 LCS 22，\(r=\allowbreak 0.832\)）。与上行高一30 第12题是同一道真题的第二种改法 \\ \addlinespace[1pt]
◇高一31 第 14 题 & \textbf{2009 年天津卷（理）选择} & \textbf{原题照录}（\(0<b<a+\allowbreak 1\)、关于 \(x\) 的不等式 \((x-b)^2>(ax)^2\) 的解集中整数恰有 \(3\) 个；\(r=\allowbreak 0.736\)、最长公共子串 22 字，答案同为 \(1<a<3\)） \\ \addlinespace[1pt]
◇高一34 第 8 题 & \textbf{2008 年江苏卷 填空} & \textbf{原题照录}（\(x\)、\(y\)、\(z\) 为正实数且 \(x-2y+\allowbreak 3z=\allowbreak 0\)，求 \(\dfrac{y^2}{xz}\) 的最小值；\(r=\allowbreak 0.967\)、最长公共子串 28 字，答案同为 \(3\)） \\ \addlinespace[1pt]
◇高一35 第 12 题 & \textbf{2012 年浙江卷（理）填空} & \textbf{原题照录}（\(x>0\) 时均有 \([(a-1)x-1](x^2-ax-1)\geqslant\allowbreak 0\)，求 \(a\)；\(r=\allowbreak 0.915\)，答案同为 \(\dfrac32\)） \\ \addlinespace[1pt]
◇高一35 附加 17 题 & \textbf{2009 年天津卷（理）选择} & \textbf{改编}：与高一31 第 14 题是\textbf{同一道真题的两种题型}------本卷把选择改成了填空「求实数 \(a\) 的范围」，条件与解集都相同，答案同为 \(1<a<3\) \\ \addlinespace[1pt]
◇高一36 第 3 题 & \textbf{2004 年湖北卷（文）选择} & \textbf{改编}：真题的 \(B\) 作 \(x\leqslant\allowbreak 6\)，本卷作 \(x\leqslant\allowbreak 5\)，答案由 \(\{1,\allowbreak 4,\allowbreak 6\}\) 变为 \(\{1,\allowbreak 4\}\)；且由选择改成填空 \\ \addlinespace[1pt]
◇高一36 第 13 题 & \textbf{2003 年上海卷（秋理）选择} & \textbf{原题照录}（``\(\dfrac{a_1}{a_2}=\allowbreak \dfrac{b_1}{b_2}=\allowbreak \dfrac{c_1}{c_2}\)''是``\(M=N\)''的什么条件；\(r=\allowbreak 0.994\)、最长公共子串 55 字，答案同为 D）。与 高一17 第 7 题、高一20 第 15 题是\textbf{同一道真题的第三种} \\ \addlinespace[1pt]
◇高一36 第 14 题 & \textbf{2010 年福建卷（文）选择} & \textbf{原题照录}（非空集合 \(S=\allowbreak \{x\mid\allowbreak  m\leqslant\allowbreak  x\leqslant\allowbreak  l\}\) 满足 \(x\in S\Rightarrow x^2\in S\) 的三个命题；\(r=\allowbreak 0.963\)、最长公共子串 38 字，答案同为 D）。\textbf{真题作「如下三个命题」}，可确证本卷题干的「四个命题」是笔误（见「高一36 专记」） \\ \addlinespace[1pt]
◇高一37 第 8 题 & \textbf{2009 年重庆卷（理）选择} & \textbf{改编}：左侧 \(\lvert x+\allowbreak 3\rvert-\lvert x-1\rvert\) 与真题\textbf{一字未动}（那正是本题的核心------左式最大值 \(4\) 的来源），右侧由 \(a^2-3a\) 改为 \(\lvert a-3\rvert\)、由选择改成填空，答案随之由「\(a\leqslant\allowbreak -1\) 或 \(a\geqslant\allowbreak 4\)」变为「\(a\leqslant\allowbreak -1\) 或 \(a\geqslant\allowbreak 7\)」。机判只报了「对任意实数 \(x\) 恒成立，则实数 \(a\) 的取值范围为」这类套话串（\(r=\allowbreak 0.892\)），但同源证据不止于此 \\ \addlinespace[1pt]
◇高一37 第 11 题 & \textbf{2010 年四川卷（文）选择} & \textbf{原题照录}（设 \(a>b>0\)，则 \(a^2+\allowbreak \dfrac1{ab}+\allowbreak \dfrac1{a(a-b)}\) 的最小值是------\textbf{题干一字不差}、答案同为 \(4\)，只去掉四个选项改成填空）。\textbf{同卷（理）版是另一道题}（条件作 \(a>b>c>0\)、式子里多出 \(-10ac+\allowbreak 25c^2\)、答案 B），机判把它也报了一条，复核后判\textbf{假命中}、不计 \\ \addlinespace[1pt]
■高一20 第 8 题 & \textbf{2023 年全国乙卷（理）选择} & \textbf{借同一对集合另问}（真题 \(U=\allowbreak \R\)、\(M=\allowbreak \{x\mid\allowbreak  x<1\}\)、\(N=\allowbreak \{x\mid\allowbreak -1<x<2\}\) 与问法「用交、并、补表示某个集合」都与本卷相同，只把所求集合换成 \(\{x\mid\allowbreak  x\geqslant\allowbreak 2\}\)；\(r=\allowbreak 0.727\)、最长公共子串 22 字） \\ \addlinespace[1pt]
■高一21 第 4 题 & \textbf{2009 年上海卷（秋文）填空} & \textbf{改编}（真题 \(A=\allowbreak \{x\mid\allowbreak  x\leqslant\allowbreak 1\}\)、\(B=\allowbreak \{x\mid\allowbreak  x\geqslant\allowbreak  a\}\)、\(A\cup\allowbreak  B=\allowbreak \R\) 求 \(a\)；本卷把端点与字母对调成 \(A=\allowbreak \{x\mid\allowbreak  x<a\}\)、\(B=\allowbreak \{x\mid\allowbreak  x\geqslant\allowbreak 1\}\)，答案随之由 \(a\leqslant\allowbreak 1\) 变 \(a\geqslant\allowbreak 1\)。\(r=\allowbreak 0.912\)、最长公共子串 16 字。同题位另有一条 \textbf{2007 福建卷（理）} 的机判候选，只共用 \(A=\allowbreak \{x\mid\allowbreak  x<a\}\) 与 \(A\cup\allowbreak  B=\allowbreak \R\) 的框架，见下面「① 机判报了、但未计入」表） \\ \addlinespace[1pt]
◇高一37 第 19 题 & \textbf{2005 年北京卷（春季·文/理）解答} & \textbf{改编}：分子 \(920\)、分母 \(v^2+\allowbreak 3v+\allowbreak 1600\)、第 (2) 问的设问与答案「\(25<v<64\)」\textbf{全同}；唯一出入是第 (1) 问的作答要求由「精确到 \(0.1\) 千辆/小时」改为「保留分数形式」（答案随之由 \(\approx11.1\) 写成 \(\dfrac{920}{83}\)）。真题文、理两卷印的是同一道，故出处写作「文/理」 \\ \addlinespace[1pt]
高一47 第 6 题 & \textbf{2012 年江苏卷填空} & \textbf{改编}：真题已知二次函数值域为 \([0,\allowbreak +\allowbreak \infty)\)，由 \(f(x)<c\) 的解集 \((m,m+\allowbreak 6)\) 求 \(c=\allowbreak 9\)；本卷保留这一独特设定、把区间长度改为 \(8\)，题目写成 \(y=x^2+bx+\allowbreak 3\)、区间 \(m-8<x<m\)，答案相应为 \(16\)（机判 LCS 15、\(r=\allowbreak 0.714\)） \\ \addlinespace[1pt]
\end{xltabular}

另有 1 道\textbf{同定义但换了问法}、未计入上表：高一02 第 3 题用了 2009 北京卷（文）「孤立元」的
定义（改称「孤立元素」，并把问题简化成「给定 \(A=\{1,2,3,5\}\) 求其孤立元素」）。

\textbf{存疑一条（未计入上表）}：高一11 第 3 题「已知集合 \(A=\{2,a\}\)，且 \(1\in A\)，则 \(a=\)」与
\textbf{2026 上海卷（秋）填空 1}「\(A=\{2,1+a\}\)，且 \(-1\in A\)」同构而数字全换。两条证据方向相反：
① 这个句式在 1952--2026 的全库里\textbf{只出现这一次}（正是 2026 秋），而本卷发布于高考后三个月，
不能排除照当年高考题改写；② 但「元素属于含参集合 $\Rightarrow$ 求参数」本身是教材同步练习的标准题型
（人教 A 版《1.1 集合的概念》就有一道 \(A=\{a-2,2a^2+5a,12\}\)、\(-3\in A\) 的经典题），而这两道
题都极简------都只含 2 个元素、都问 \(a\)------更像是各自从同一常见题型简化而来。证据不足以断，
故单列于此，\textbf{不计入上面的 43 道}。

其中 \textbf{2010 上海秋理第 14 题}这次以原卷原文核实过：题面是「从集合 \(U=\{a,b,c,d\}\) 的
子集中选出 \textbf{4 个}不同的子集，需同时满足 ① \(\varnothing\)、\(U\) 都要选出 ② 任意两个子集
\(A\)、\(B\) 必有 \(A\subseteq B\) 或 \(A\supseteq B\)」，官方答案 36------与 高一b 第 11 题
源码注记的记载一致（网上流传的「选出 2 个」版本是讹传）。

\textbf{① 机判报了、但未计入本表的题位}（2026-09-28 建表：把 题目溯源.py ①「与高考库比对」报过的
\textbf{每一个本库题位}都摊开判一遍并留档；判据是「共享的是那道真题的\textbf{核心设定}，还是通用句式」------
前者登记进上表，后者记在这里。这张表已由 \texttt{工具/\allowbreak{}计数核对.py} 按 \texttt{工具/\allowbreak{}真题命中清单.txt} 守着，
\textbf{以后每轮跑完 ① 若多出一个报过的题位，闸就会报错、逼人在这里或上表里 disposition}）：

\newcommand{\tabhdrI}{\rowcolor{white} 机判报过的题位 & 机判出处 & 为什么不计入 \\}
\rowcolors{1}{white}{rowgray}
\begin{xltabular}{\textwidth}{@{}>{\raggedright\arraybackslash}p{2.82cm}>{\raggedright\arraybackslash}p{4.23cm}>{\raggedright\arraybackslash}p{7.69cm}@{}}
\toprule
\tabhdrI
\midrule
\endfirsthead
\multicolumn{3}{r}{\footnotesize（续上表）} \\
\midrule
\tabhdrI
\midrule
\endhead
\bottomrule
\endlastfoot
高一02 第 18 题 & 2018 上海卷（春）填空 & 只共用 \(A=\allowbreak \{x\mid\allowbreak 0<x<2\}\) 这个区间；\(B\)（真题是固定区间、本卷是含参）与问法都不同 \\ \addlinespace[1pt]
高一04 第 15 题 & 1999 上海卷（文）、（理）解答 & 只共用「\(A\subseteq\allowbreak  B\)，求实数 \(a\) 的取值范围」这句套话 \\ \addlinespace[1pt]
高一05 第 7 题 & 2007 北京卷（理）填空 & 只共用「若 \(A\cap\allowbreak  B=\allowbreak \varnothing\)，则实数 \(a\) 的取值范围是」这句套话 \\ \addlinespace[1pt]
高一06 第 7 题 & 2021 全国乙卷（理）选择 & 借了真题第一集合 \(S=\allowbreak \{s\mid\allowbreak  s=\allowbreak 2n+\allowbreak 1\}\) 的取法，但 \(T\) 换了（\(2n-1\) vs \(4n+\allowbreak 1\)）、问法也换（判 \(S\)、\(T\) 的包含关系 vs 求 \(S\cap\allowbreak  T\)）；与已登记的 高一17 第 9 题不是同一版 \\ \addlinespace[1pt]
高一11 第 4 题 & 2001 新课标卷（文）、2007 安徽卷（文）选择 & 各只共用其中一个二次方程（\(x^2+x=\allowbreak 0\)／\(x^2=\allowbreak 1\)） \\ \addlinespace[1pt]
高一11 第 5 题 & 2009 新课标卷（文）、2009 海南宁夏新课标卷（理）、2018 全国 II 卷（文）选择 & 只共用「已知集合 \(A=\allowbreak \{1,\allowbreak 3,\allowbreak 5,\allowbreak \cdots\}\)」这个开头 \\ \addlinespace[1pt]
高一11 第 20 题 & 2017 北京卷（文）选择 & 只共用「集合 \(A=\allowbreak \{x\mid\allowbreak  x<-2\) 或 \(x\cdots\}\)」 \\ \addlinespace[1pt]
高一15 第 4 题 & 2009 上海卷（秋文）填空 & 只共用「且 \(A\cup\allowbreak  B=\allowbreak \R\)，则实数 \(a\) 的取值范围是」这句套话 \\ \addlinespace[1pt]
高一19 第 3 题 & 2013 湖北卷（文）、2016 浙江卷（文）、2018 浙江卷、2021 全国乙卷（文）选择 & 只共用「已知全集 \(U=\allowbreak \{1,\allowbreak 2,\allowbreak 3,\allowbreak 4,\allowbreak 5\}\)」 \\ \addlinespace[1pt]
高一20 第 3 题 & 2007 北京卷（理）填空 & 只共用「\(\varnothing\)，则实数 \(a\) 的取值范围是」这句套话 \\ \addlinespace[1pt]
高一26 第 7 题 & 2019 江苏卷 填空 & 只共用「\(B=\allowbreak \{x\mid\allowbreak  x>0,\allowbreak \ x\in\R\}\)，则 \(A\cap\allowbreak  B=\allowbreak \)」 \\ \addlinespace[1pt]
高一26 第 8 题 & 1999 上海卷（文）解答 & 只共用「若 \(A\subseteq\allowbreak  B\)，求实数 \(a\) 的取值范围」这句套话 \\ \addlinespace[1pt]
高一28 第 5 题 & 2007 北京卷（理）填空 & 只共用 \(B=\allowbreak \{x\mid\allowbreak  x^2-5x+\allowbreak 4\geqslant\allowbreak 0\}\) 这个常见集合 \\ \addlinespace[1pt]
高一28 第 18 题 & 2007 上海卷（春）填空 & 只共用「\(>0\) 的解集为 \((-\infty,\allowbreak \cdots\)」 \\ \addlinespace[1pt]
高一30 第 19 题 & 2017 北京卷（文）选择 & 只共用「集合 \(A=\allowbreak \{x\mid\allowbreak  x<-2\) 或 \(x\cdots\}\)」（2026-09-25 那一轮已判过） \\ \addlinespace[1pt]
高一36 第 9 题 & 2014 浙江卷（文）填空 & 只共用「已知实数 \(a\)、\(b\)、\(c\) 满足 \(a+b+c=\allowbreak 0\)」；真题问的是平方和条件下的最值，与本卷 \(\dfrac ca\) 的范围无关（2026-09-26 那一轮已判过） \\ \addlinespace[1pt]
高一40 第 9 题 & 2012 福建卷（文）填空 & 真题是 \(x^2-ax+\allowbreak 2a>0\) 在 \(\R\) 上恒成立、本卷是 \(x^2-ax-2<0\) 在 \((1,\allowbreak 2)\) 内恒成立，除套话外全不同（2026-09-27 那一轮已判过） \\ \addlinespace[1pt]
高一47 第 11 题 & 2012 浙江卷（理）单选；2007 安徽卷（文）单选 & 浙江题只与本卷共用 \(B=\allowbreak \{x\mid\allowbreak  x^2-2x-3\leqslant\allowbreak 0\}\)，但 \(A\) 的定义与所求交集不同；安徽题用同一二次式的零点集（等式 \(=\allowbreak 0\)）而非本卷的区间（不等式 \(\leqslant\allowbreak 0\)），且 \(A\) 也不同。两条都只有一个集合片段相同，非同题。 \\ \addlinespace[1pt]
\end{xltabular}

\clearpage
\section{考据不出来的：绝大多数}

把 \textbf{30 份}试卷里的\textbf{自定义名词}逐个在全库检索：

\begin{itemize}
\tightlist
\item
  \textbf{全库皆无}：闭集合、和谐集、好集、划分集、互斥子集、可积数集、二分划、三数循环、
  全食／偏食、高斯函数、单封集、\textbf{萌数}、\textbf{可分集合}、\textbf{破晓集}（高一30 第 10 题；
  对加、减、乘都封闭的数集）、\textbf{变项和}（高一30 第 11 题）、\textbf{交替和／交替积}（高一30 第 21 题）。（这些名词在 1952--2026 的高考真题库里确实一次都没出现；
  但「检索不到」只等于「不是高考真题」------例如「和谐集」就已考据出非高考的出处，即上海市
  上海中学东校 2023--2024 学年高一上期中卷，见「附加题专查」；「萌数」也在\textbf{百度题库}里
  检到同题面，见「非高考但能考据到来源的题」。）
\item
  \textbf{同名不同题}：「特征数」见于 2013 湖南卷（文）与 1977 江苏常州卷（前者的指子集的 0/1
  指示数列，后者的含义与之也不同），与 高一07／高一a 的「最大元素与最小元素之和」
  \textbf{同名不同义}。
\item
  \textbf{「封闭集」那条要订正}（2026-09-15 重跑）：此前写「封闭集只与 2010 四川卷（文/理；
  定义是 \(x+y,\ x-y,\ xy\in S\)）\textbf{借名不借题}」------这条对 2010 四川卷而言仍是「借名」，
  但\textbf{高一05 那道题的真正源头不是它}：高一05 第 19 题（「性质 \(M\)」）与高一05 第 20 题
  （「封闭集」）、高一12 第 21 题（「单封集」）同出一源，即 \textbf{2009 北京卷（理）的
  「性质 \(P\)」}（\(a_ia_j\) 与 \(a_j/a_i\) 至少有一个属于 \(A\)）：05 第 19 题是它的乘除版，
  05 第 20 题与 12 第 21 题是它的加减版（\(a_i+a_j\) 与 \(a_j-a_i\)）。三者已改列入上面的
  出处表，此处不再算「考据不出来」。
\end{itemize}

也就是说，这些「新定义压轴题」\textbf{都不是高考真题}，应属校本自编，或取材自模考／竞赛／自招
与教辅。

\clearpage
\section{非高考但能考据到来源的题}

高考真题库里查不到的题，有一部分能在\textbf{教辅/题库/校际试卷}里找到同题或同定义版本
（2026-09-15 第二轮复核时逐题做了网络检索，取「同题面」或「同定义＋同小问」为准）：

\newcommand{\tabhdrJ}{\rowcolor{white} 题 & 考据结果 \\}
\rowcolors{1}{white}{rowgray}
\begin{xltabular}{\textwidth}{@{}>{\raggedright\arraybackslash}p{6.97cm}>{\raggedright\arraybackslash}p{8.18cm}@{}}
\toprule
\tabhdrJ
\midrule
\endfirsthead
\multicolumn{2}{r}{\footnotesize（续上表）} \\
\midrule
\tabhdrJ
\midrule
\endhead
\bottomrule
\endlastfoot
高一11 第 21 题「闭集合」 & \textbf{教辅/题库常见题}------「给定数集 \(M\)，若对任意 \(a\)、\(b\in M\) 有 \(a+b\in M\) 且 \(a-b\in M\)，则称 \(M\) 为闭集合」这一题面（连同「\(C\cup\allowbreak  D\) 是否一定为闭集合」这一小问）在百度文库、作业帮、青夏教育、百度题库都能检索到 \\ \addlinespace[1pt]
高一a 第 15 题「性质 \(P\)」（\(\operatorname{card}(A+A)=\allowbreak \frac{n(n+\allowbreak 1)}{2}\)） & \textbf{菁优网收录同题}：该站版本是 \(A=\allowbreak \{1,\allowbreak 3,\allowbreak 5\}\)、\(B=\allowbreak \{2,\allowbreak 4,\allowbreak 8\}\) 与 \(2021\)，本卷换成 \(\{1,\allowbreak 4,\allowbreak 7\}\) 与 \(2022\) \\ \addlinespace[1pt]
高一06 第 14 题「好集合」 & \textbf{百度题库收录同题}，且题库里有小问更多的完整版（如「若 \(\lvert S\rvert=\allowbreak 2019\)，求证 \(S\) 中存在元素 \(m\)，使 \(S\) 中所有元素均为 \(m\) 的整数倍」），本卷只取前两问 \\ \addlinespace[1pt]
高一13 第 10 题「好子集」 & \textbf{菁优网收录同定义题}：该站版本是 \(I=\allowbreak \{1,\allowbreak 3,\allowbreak 5,\allowbreak 7\}\)，本卷换成 \(\{1,\allowbreak 2,\allowbreak 3,\allowbreak 4,\allowbreak 5\}\) \\ \addlinespace[1pt]
高一02 第 6 题「伙伴关系集合」 & \textbf{教辅/题库常见题}：作业帮、青夏教育、菁优网各有不同 \(M\) 的版本（\(\{-1,\allowbreak 0,\allowbreak \frac12,\allowbreak 1,\allowbreak 2\}\)／\(\{-1,\allowbreak 0,\allowbreak \frac13,\allowbreak \frac12,\allowbreak 1,\allowbreak 2\}\) 等），本卷用的是 \(\{-1,\allowbreak 0,\allowbreak \frac12,\allowbreak 1,\allowbreak 2,\allowbreak 3,\allowbreak 4\}\) \\ \addlinespace[1pt]
\textbf{高一22 第 16 题「萌数」} & \textbf{百度题库收录同题}：该站的题面与三个条件（\(A\)、\(B\)、\(C\) 为 \(M=\allowbreak \{1,\allowbreak 2,\allowbreak \dots,n\}\) 的三个非空子集且两两不交、并集为 \(M\)；\(A\) 全奇、\(B\) 全偶、\(3\) 的倍数都在 \(C\)；\(S_1=S_2=S_3\)）及三小问（\(n=\allowbreak 10\)、\(n=\allowbreak 8\)、证明 \(n=\allowbreak 6k+\allowbreak 2\)）与本卷\textbf{逐字相同}------2026-09-16 复核本卷题面时检得（该站详情页需验证码，故只核到检索结果的题面片段，未取到原始出处与年份） \\ \addlinespace[1pt]
\end{xltabular}

其余「真题库与网络都检索不到」的，是各卷的自编压轴与自编小题（如高一02 第 16 题只借用了
「指示函数」这个集合论标准名词；高一15 附加 18、高一09 第 12 题「划分集」、\textbf{高一22 附加 17／18}
等整题无同题）。另有一处卷面小疵顺带记下：\textbf{高一02 第 15 题}的两个条件原卷编号印作「①」「④」（跳过②③），
本稿照原卷录入、未擅改。

\clearpage
\section{附加题专查（最早带附加题的五份 14 题）}

本节逐题考据的是\textbf{最早带附加题的五份、合计 14 道}：\texttt{高一09}（进才）、\texttt{高一a}（大同中学）、
\texttt{高一b}（大同）、\texttt{高一15}（七宝）、\texttt{高一22}（七宝）。\textbf{全稿现共 8 份卷带附加题、合计 26 道}
------另三份是 \texttt{高一35}（七宝 · 抱孩子练习9.22，附加 17---20）、\texttt{高一39}（大同 · 周末作业4.1，
附加 21---24）与 \texttt{高一47}（七宝 · 抱孩子练习9.26，附加 1---4），各 4 道；它们的考据见「高一31---高一36 专记」、「高一37---高一40 专记」与「高一47 专记」
（其中高一35 附加 17 题＝2009 天津卷（理）已入「真题出处」表；高一39 附加 21 题与 高一a 附加 2 题
是同一道，故也是 \textbf{2010 湖南卷（文）填空 7} 那道真题的一个题位------2026-09-28 补登进表，
两处并见同题表与出处表）。\textbf{「8 份／26 道」与本节这「五份／14 道」都已由 \texttt{工具/\allowbreak{}计数核对.py}
按源码逐卷复算守住}（原先这里写的是「全稿共 5 份／14 道」，那个「全稿」随加卷过期了）。
其中 \textbf{高一09 的附加题卷面上没有编号}（只印「附加题：」
直接接题干，本稿按通例排成「1.」------已在下表注明），其余四份都\textbf{另起编号}或接续正文
（高一a 另起「1.」「2.」；高一b 接续作 19---21；高一15 与 高一22 都接续作 17---20）。十四题都在原图
（高一a／高一b 见 \texttt{04.jpg}，高一09 见 \texttt{07.png}，高一15 见第 6---7 页，高一22 见第 5---7 页）上逐字核对过
（高一a 附加题第 2 题的下标 \(i\) 一处笔误已订正），2026-09-15 逐题考据如下
（\textbf{高一22 的四道是 2026-09-16 入库时补入的}）：

\newcommand{\tabhdrK}{\rowcolor{white} 卷·题 & 内容 & 考据结果 \\}
\rowcolors{1}{white}{rowgray}
\begin{xltabular}{\textwidth}{@{}>{\raggedright\arraybackslash}p{2.04cm}>{\raggedright\arraybackslash}p{4.91cm}>{\raggedright\arraybackslash}p{7.78cm}@{}}
\toprule
\tabhdrK
\midrule
\endfirsthead
\multicolumn{3}{r}{\footnotesize（续上表）} \\
\midrule
\tabhdrK
\midrule
\endhead
\bottomrule
\endlastfoot
09 附加题（卷面无编号） & 「和谐集」三小问（判断 \(\{1,\allowbreak 2,\allowbreak 3,\allowbreak 4,\allowbreak 5\}\)、求证 \(\{1,\allowbreak 3,\allowbreak 5,\allowbreak 7,\allowbreak 9,\allowbreak 11,\allowbreak 13\}\)、证明元素个数为奇数） & \textbf{可确指校际试卷}：同题见于 \textbf{《上海市上海中学东校 2023--2024 学年高一上学期数学期中试卷》}------组卷网（21 世纪教育网）标注其来源为「(2023 高一上·上海市期中)」并直接引用该卷，同题另被《备考 2024 年高考数学优生冲刺专题特训》等收录。本卷三小问与该卷\textbf{完全一致}；\textbf{高一02 第 19 题是它的改编}（改问「元素个数的最小值」）。真题库无 \\ \addlinespace[1pt]
a 附加 1 & 「互斥子集」\(f(6)\) & 真题库无（参考库里的「互斥」7 条全是概率的「互斥事件」）。教辅/题库常见题（青夏教育、作业帮、百度题库、洛谷 T421094 均可检索到），\textbf{但网上有两个版本}：主流版作「\textbf{视 \((A,B)\) 与 \((B,A)\) 为同一组}」$\to$ \(f(n)=\allowbreak \frac{3^n-2^{n+\allowbreak 1}+\allowbreak 1}{2}\)、\(f(6)=\allowbreak 301\)；\textbf{本卷原卷作「当且仅当 \(A=B\) 时\ldots 为同一组」}（与洛谷 T421094 同），因 \(A\)、\(B\) 非空且不交、\(A=B\) 不可能成立，等价于「两者总是不同组」$\to$ 有序对 $\to$ \(f(6)=\allowbreak 602\)。\textbf{卷面蓝笔订正写的正是 602}，本稿即按原卷口径给出 \\ \addlinespace[1pt]
a 附加 2 & 「第 \(k\) 个子集」（2 的幂编码），问第 25 个子集 & \textbf{2010 湖南卷（文）填空 7}（见出处表；本条是 2026-09-15 新补出的漏检，与高一08 第 11 题同源、改法不同） \\ \addlinespace[1pt]
b 附加 19 & \(A\subseteq\allowbreak  B\subseteq\allowbreak  C\subseteq\allowbreak \{1,\allowbreak \dots,\allowbreak 2024\}\)，有序集合组的个数 \(4^{2024}\) & 真题库\textbf{零命中}（16,273 道里连「有序集合组」或 \(A\subseteq\allowbreak  B\subseteq\allowbreak  C\) 的结构都没有）。同类题的近亲是「\(A,B\subseteq\allowbreak  C\subseteq\allowbreak \{1,\allowbreak \dots,\allowbreak 2020\}\) $\to$ \(5^{2020}\)」，见于\textbf{强基计划}题解（博客园《{[}强基{]}3，有趣的集合个数算法》）$\to$ 判为\textbf{强基／自招风格的计数题＋校本改编}（上界改成 2024、条件收成链式包含） \\ \addlinespace[1pt]
b 附加 20 & \(M=\allowbreak \{1,\allowbreak \dots,\allowbreak 10\}\) 的非空子集按 \(\sum(-1)^kk\) 求和，总和 \(2560\) & 真题库无；\textbf{与高一09 第 11 题是同一道}（两版只换举例：09 用 \(\{1,\allowbreak 4,\allowbreak 7\}\to\allowbreak -4\)，大同用 \(\{2,\allowbreak 3,\allowbreak 6\}\to\allowbreak 5\)）。网络上是\textbf{教辅/题库常见题}：青夏教育上至少三个版本，分别叫「无名称」「\textbf{非常元素和}」「\textbf{独立和}」，答案都是 2560 \\ \addlinespace[1pt]
b 附加 21 & 四个有理数、两两乘积集合给定，求四数之和 \(\pm\frac94\) & 真题库无；\textbf{同一道题 2015-09-13 见于 Math StackExchange}（编号 1432967，英文原题、未标出处，两个解答给出的答案 \(\pm\frac94\) 与本稿一致），百度知道亦有中文版 $\to$ \textbf{2015 年前已在网上流传的竞赛／自招风格组合题}，确切出处考据不到 \\ \addlinespace[1pt]
15 附加 17 & 「孤立元」（\(S=\allowbreak \{1,\allowbreak \dots,\allowbreak 8\}\) 的 3 元子集不含孤立元） & \textbf{2009 北京卷（文）填空 6 的原题照录}（见出处表；与高一06 第 9 题互为逐字同题） \\ \addlinespace[1pt]
15 附加 18 & \(A\cap\allowbreak  B=\allowbreak \varnothing\)，\(M=\allowbreak \{T\mid\allowbreak  T\subseteq\allowbreak  A\}\)，\(P=\allowbreak \{S\mid\allowbreak  S\subseteq\allowbreak  B\}\)，求 \(M\cap\allowbreak  P\)（答 \(\{\varnothing\}\)） & 真题库\textbf{零命中}（LCS \(=\allowbreak 0\)），网上亦检索不到 $\to$ 附加题里检索不到同题的一道（幂集交为 \(\{\varnothing\}\) 的辨析小题）；22 附加 17／18 同样检不到\textbf{原式}，但那两道的\textbf{题型}在教辅里常见，本题连题型都无同题 \\ \addlinespace[1pt]
15 附加 19 & \(A=\allowbreak \{x\mid\allowbreak  x=m^2-n^2,\allowbreak \ m,n\in\mathbb{Z}\}\)，写出所有满足集合 \(A\) 的偶数（答 \(4k\)） & 真题库无；\textbf{教材/教辅常见题}（百度题库、作业帮、青夏教育都有「\(A=\allowbreak \{x\mid\allowbreak  x=m^2-n^2\}\)」的同型题，其中一份标为「高一上学期数学期末考试试卷第 73 套真题」）。\textbf{与高一13 第 18 题是同一个 \(A\)}------那版是三小问（判断 \(8,\allowbreak 9,\allowbreak 10\)、证必要非充分、求所有偶数），本卷只留最后一问的结论 \\ \addlinespace[1pt]
15 附加 20 & \(n\) 元「调和子集」（\(a_1+\allowbreak \cdots+a_n=a_1\cdots a_n\)），求自然数集 \(\mathbb{N}\) 的 3 元调和子集（答 \(\{1,\allowbreak 2,\allowbreak 3\}\)） & 真题库无；\textbf{是教辅/题库里「好集」那道题的改名版}------青夏教育、作业帮上的同题作「若 \(a_1+\allowbreak \cdots+a_n=a_1a_2\cdots a_n\)，则称集合 \(A\) 是集合 \(M\) 的 \(n\) 元`好集'」。注意与高一08／高一14 那道同名的「好集」，其定义（\(0,\allowbreak 1\in A\)；\(x-y\in A\)；\(x\neq0\) 时 \(1/x\in A\)）\textbf{完全不同}：同名两义，本卷改称「调和子集」反倒避开了撞名 \\ \addlinespace[1pt]
22 附加 17 & 方程 \(x^2-(2m-2)x+\allowbreak 5m+\allowbreak 8=\allowbreak 0\) 的两个解是直角三角形两直角边的长、斜边长为 10，求 \(m\)（答 \(8\)） & 真题库\textbf{零命中}。命制思路是两步套：\textbf{用韦达定理得 \(x_1+x_2=\allowbreak 2m-2\)、\(x_1x_2=\allowbreak 5m+\allowbreak 8\)，再由 \(x_1^2+x_2^2=\allowbreak 100\) 联立}------解得 \(m=\allowbreak 8\) 或 \(m=\allowbreak -\frac72\)，还须用「两解皆正」舍去后者，这是本题的实际难点（原卷只给答案 \(8\)、未展开）。网络检索\textbf{未找到含这组系数的原式}（「二次方程两根作直角三角形两直角边」这一题型本身常见，但换一组系数即无从比对），判为\textbf{校本编写或改编} \\ \addlinespace[1pt]
22 附加 18 & \(M=\allowbreak \{(x,y)\mid\allowbreak \frac{y-3}{x-2}=a+\allowbreak 1\}\)、\(T=\allowbreak \{(x,y)\mid\allowbreak (a^2-1)x+\allowbreak (a-1)y=\allowbreak 15\}\)，\(M\cap\allowbreak  T=\allowbreak \varnothing\) 求 \(a\)（答 \(\pm1\)、\(-4\)、\(\frac52\)） & 真题库\textbf{零命中}。\textbf{设计点在「\(M\) 是挖掉 \((2,\allowbreak 3)\) 的直线」}------四组解各有来历：\(a=\allowbreak 1\) 时 \(T=\allowbreak \varnothing\)，\(a=\allowbreak -1\) 时 \(M\) 与 \(T\) 是两条平行线，\(a=\allowbreak \frac52\) 与 \(a=\allowbreak -4\) 两组的交点恰好落在被挖掉的 \((2,\allowbreak 3)\) 上。网络检索\textbf{未找到同题}（「含参直线与挖点直线求交」是教辅里成体系的一类，但本组系数与问法无对应版本），判为\textbf{校本编写} \\ \addlinespace[1pt]
22 附加 19 & 有限集合的基数 \(\operatorname{card}\) 四条命题，问真命题序号（答①②） & \textbf{2006 年湖北卷（理）选择 原题照录}（见「真题出处」表）：四条命题与真题相同、答案同为①②；差别只有两处表述------第 ③ 条真题印 \(\nsubseteq\) 的符号写法、本卷写成文字「\(A\) 不是 \(B\) 的子集」（同义），以及由选择改成了填空 \\ \addlinespace[1pt]
22 附加 20 & \(M=\allowbreak \{930,\allowbreak -5,\allowbreak -1,\allowbreak 0,\allowbreak \pi\}\) 的 31 个非空子集各自的元素乘积之和（答 \(-1\)） & 真题库\textbf{零命中}。\textbf{是「含 \(0\) 与 \(-1\) 的集合配成相反数对」这一经典构造}：含 \(0\) 的 16 个子集贡献全为 \(0\)，剩下的按「含不含 \(-1\)」两两配对、乘积互为相反数，最后只剩 \(\{-1\}\) 一个。网上同型题很多------\textbf{洛谷 U657612 是它的程序化推广}（题面自述「做到一道题\ldots\ldots 把题面推而广之」，样例元素即取 \(0\) 与 \(-1\)），故判为\textbf{教辅/题库常见题 + 校本换元素}（把常见的 \(1\sim n\) 换成 \(930,\allowbreak -5,\allowbreak \pi\) 这类彼此无关的数） \\ \addlinespace[1pt]
\end{xltabular}

小结：\textbf{十四道附加题里能直接对上高考真题的有三道}（\textbf{22 附加 19 ＝ 2006 湖北卷（理）原题照录}、
15 附加 17 ＝ 2009 北京卷（文）原题照录、a 附加 2 ＝ 2010 湖南卷（文）改编），一道可确指校际试卷
（09 附加题 ＝ 上海中学东校
2023--2024 学年高一上期中卷，大同那份是它的改编），五道在教辅/题库里有迹可循（a 附加 1、
b 附加 20、15 附加 19、15 附加 20、22 附加 20），b 附加 19 与 b 附加 21 的最近亲分别是强基计划题解
与 2015 年就在网上流传的竞赛／自招组合题，22 附加 17／18 与 15 附加 18 则连网络检索都找不到同题
（前两道检不到的是\textbf{原式}、题型在教辅里常见，15 附加 18 连题型都无同题）。这与
「附加题＝压轴拔高」的定位吻合：各校多从题库、强基／自招题、校际试卷与真题改编里挑。

\begin{quote}
\textbf{2026-09-16 订正}：本小结原写「十道附加题没有一道是高考真题」，与其后紧跟的「两道能直接
对上真题」自相矛盾（15 附加 17 正是 2009 北京卷（文）的\textbf{原题照录}）。并入 高一22 时一并
按实改写：上表里能对上真题的是\textbf{三道}（两道原题照录 + 一道改编），并非「零道」。
\end{quote}

\clearpage
\section{方法与局限}

\begin{itemize}
\tightlist
\item
  方法：把题目脱去 LaTeX 与标点做「骨架」文本 $\to$ 6-gram 倒排索引全召回 $\to$ 最长公共子串 +
  整段相似度判定 $\to$ 人工复核高置信项。\textbf{参考实现已随仓库提交：\texttt{papers/\allowbreak{}工具/\allowbreak{}题目溯源.py}}
  （\texttt{python3\ 题目溯源.py\ 全跑} 即抽取两库并比对；判据、阈值与踩过的坑都写在文件头注释里，
  中间结果落在系统临时目录，不污染仓库）。\textbf{两个判据要一起看}：长题靠「最长公共子串 $\geqslant$ 12」，
  短题（如一句式的选择题）最长串天然长不了，靠「整段相似度 $\geqslant$ 0.71」补上。上表 50 条对应里，
  高一08 第 11 题、高一a 第 4 题与 高一b 第 7 题、高一03 第 12 题 都是只靠后者才捞出来的。
\item
  \textbf{两处实现细节}（2026-09-15 重跑时踩到）：① 本库那侧要先剥 \texttt{\%} 注释再按 \texttt{\textbackslash{}item} 切分------
  源码注释里会出现字面 \texttt{\textbackslash{}item}（如「在其 \texttt{\textbackslash{}item} 前先量一下版面」），留着会被当成一道题，
  实测让高一12 多出一道假题、题号整体后移一位；② 归一化时 \texttt{\textbackslash{}cap} \texttt{\textbackslash{}cup} \texttt{\textbackslash{}subseteq}
  \texttt{\textbackslash{}varnothing} \texttt{\textbackslash{}mid} 必须先换成字符再剥命令（否则剥成空串，见下条），且\textbf{选项文字也会
  制造假阳性}------\texttt{高一c\ 第\ 13\ 题\ ×\ 高一13\ 第\ 13\ 题} 的 LCS 达 27，公共串却是四个选项的
  套话「充分不必要／必要不充分／充要／既不充分也不必要」，并非同题；③ \textbf{记号的字母不同也会
  把公共串切断}------\texttt{高一a\ 附加题第\ 2\ 题} 与 2010 湖南卷（文）填空 7 的定义本该逐字相同，只因
  一个写 \(\{a_{i_1},\dots\}\)、另一个写 \(\{a_{b_1},\dots\}\)（外加集合名 \(M\) 与 \(E\) 不同），
  最长公共子串被切到只有 10 字、\(r=0.596\)，两项判据双双落在线下，是靠「同一道真题的两种改编」
  这条内容线索才捞回来的（已补进出处表与手工补记）。\textbf{字母级别不同的同题，一定要人工过一遍。}
\item
  \textbf{两个曾导致漏检的口径陷阱}（已修，记此备查）：① 两边必须用\textbf{同一套}归一化------参考库
  那侧若只做骨架化、不先剥 LaTeX 命令，\texttt{\textbackslash{}operatorname\{card\}} 会留下 \texttt{operatornamecard}
  这类残渣，与已经剥干净的高一文本对不上，实测把同一道题的相似度从 38 压到 12；
  ② 预筛不能用「共享片段数 top-K」截断，候选一多真匹配就会被挤出榜外。
\item
  \textbf{题号口径}：一律按\textbf{卷面实际编号}（源码里靠 \texttt{\textbackslash{}setcounter\{enumi\}} 还原，与 \texttt{\textbackslash{}item} 的
  先后顺序不同）；\textbf{附加题也照卷面}------高一b 接续正文作 19---21、高一15 与 高一22 都接续作 17---20、
  高一a 另起编号作「1.」「2.」，而高一09 卷面\textbf{根本没有编号}（只印「附加题：」，本稿按
  通例排成「1.」），故本稿一律称它「09 附加题」，不写作「09 第 22 题」。引用出处时同理：
  不能用 \texttt{\textbackslash{}item} 的出现次序编号。比对前会把参考库里个别「答案/解析被误排进 \texttt{problem} 块」的条目先截断，
  否则会压低整段的相似度而漏检。
\item
  \textbf{原图核对}：跨卷同题的旧一轮 \textbf{46 道成员逐题对照 \texttt{原图/\allowbreak{}} 里的原卷图}核过一遍（不只核
  题干，连集合的端点开闭、条件用的是 \(\cap\) 还是 \(\cup\)、答案的方括号都逐字比对）。这轮核出
  \textbf{5 处转录错误}（见前面订正清单第 7---11 条），全部是「解析对、题干或 \texttt{\textbackslash{}ans} 错」这一类------
  因为解析是从原卷整段抄的，抄错的多在需要转写的地方。\textbf{骨架文本比对查不出这类错}
  （\texttt{\textbackslash{}cap}、\texttt{\textbackslash{}cup}、\texttt{\textbackslash{}subseteq} 脱去 LaTeX 后都变成空串），必须回原图。
  新增的 6 道成员（高一14 第 20 题，高一15 第 14／16／17 题，高一12 第 3／4 题）属于
  高一12／13／14／15／17 五份，已在 2026-09-15 那轮\textbf{整卷回原图}核对里逐字比过（见专记）。
  \textbf{2026-09-15 末并入高一19---21 后，两档新增或改判的每一对又都回原图核了一遍}（逐字档 10 个
  题位、变体档 8 个题位），本轮\textbf{没有查出转录错误}，但查出\textbf{一对假同题}------见下条。
\item
  \textbf{骨架里没有图形，图形题必须单独看图（本轮新增的教训）}：\texttt{高一08\ 第\ 15\ 题\ ×\ 高一19\ 第\ 7\ 题}
  两边的题干几乎\textbf{一字不差}（都是「\(M\)、\(P\)、\(S\) 是 \(X\) 的三个子集，则阴影部分所表示的集合
  是」，\(r=0.857\)），骨架判据因此把它判成逐字同题。可回原图一看------\textbf{图形配置不同}：08 版
  \(S\) 内含于 \(P\)、阴影是 \((M\cap P)\setminus S\)（选 C），19 版 \textbf{\(P\) 内含于 \(S\)}、阴影是
  \((M\cap S)\setminus P\)------\textbf{不是同题}，已从两档剔除。凡题干高度相同而含「如图／图中」的
  配对，都要把两边的 \texttt{figures/\allowbreak{}venn\_*.png} 并排看一遍再定档。
\item
  \textbf{选择题的选项不计入骨架}（2026-09-15 补跑高一10 时落实）：四个选项
  （「A. 充分不必要条件 B. 必要不充分条件 C. 充要条件 D. 既不充分又不必要条件」）是\textbf{跨题复用的
  套话}，留在骨架里会让两道毫不相干的选择题凭选项文字拿到 LCS 27--30 的假命中------下一条「制造
  假阳性」记的就是这个坑。补跑时先在骨架里按行首「A.」把选项切掉，逐字档的分组才稳定（本次
  据此剔除了 1 组纯选项命中）；上表 24 组 = 上一轮人工核定的 15 组 + 补跑高一10 新增的 2 组

  \begin{itemize}
  \tightlist
  \item
    2026-09-15 末并入高一19---21 后新增或改判的 7 组。
  \end{itemize}
\item
  \textbf{纯数字的公共串不算文本证据}：候选取线（LCS $\geqslant$ 12 字）本身很松，元素表一类的数字串
  （\texttt{=1,2,3,4,5,6,7}）就能凑够 12 字。故判「逐字同题」时还要求\textbf{公共串里至少 8 个汉字}，
  变体档与符号档则要求\textbf{公共串里含「内容字符」}------汉字、拉丁字母或关系符号三者有一个即可
  （2026-09-15 末把「至少 1 个汉字」放宽成这个口径，纯数字串仍被挡掉，但
  \texttt{B=\{x\textbar{}x²+2(a+1)x+a²-\allowbreak{}5=0\}} 这类全符号的公共串不再误伤），
  这一条也正是「高一10 全部 19 题无真题」这个结论的判据（它最好的几个匹配都栽在这一条）。
\item
  \textbf{通用句尾也会凑出够长的公共串，机械结果必须人工过一遍}：像「则实数 \(m\) 的取值范围是」
  「求实数 \(a\) 的取值范围」这类\textbf{题干套话}，LCS 轻松到 15--16 字、汉字数也够------它只是「含参
  求范围」这一大类的共同句尾，并非同题。2026-09-15 补跑高一10 时机械跑出 4 对这样的「变体」
  （与 11\#11、12\#4、13\#3、13\#5，LCS 15--16、\(r\le0.59\)），逐条看过\textbf{都不算}，故未收进变体表；
  下表 20 组都是「换了数字或设定但题干骨架高度相同」的那一类。
  \textbf{2026-09-15 末这条判据本身也修了}（见上面「下一条」与本节末）：原先变体档写的是
  「子串 $\geqslant$ 12 且含 $\geqslant$ 8 个汉字，\textbf{或} \(r\ge0.72\) 且含 $\geqslant$ 5 个汉字」------\textbf{前半句没有 \(r\) 门}，
  于是「则实数 \(a\) 的取值范围是」这类套话只要凑够 12 字、8 个汉字就被判成变体（\(r\) 低到 0.380）；
  同时它又只有正文档、没有符号档，把 02\#17 $\times$ 05\#17（\(r=0.933\)）、03\#12 $\times$ 09\#7（\(r=0.800\)）
  这类\textbf{公共串里一个汉字都没有}的真变体挡在门外。现改成 README 本来写的那条合取式
  「\(r\ge0.71\) $\wedge$ 子串 $\geqslant$ 12 $\wedge$ 公共串含内容字符」，两头都顺了：套话被 \(r\) 门挡掉，
  符号题被内容字符闸放行。
\item
  \textbf{两路补扫（2026-09-15 第二轮）}：判据（LCS ＋ 整段相似度）对\textbf{换了数字或字母的同题}天生
  不敏感，故补了两路------① \textbf{数字脱敏}：把两边骨架里的阿拉伯数字全部抽掉再比一次，专找
  「同结构、换数字」的改编，揪出了高一02 第 15 题（2012 江苏卷、\(2x\to3x\)）与高一09 第 15 题
  （2017 江苏卷、\(a^2+3\to a^2+1\)）；② \textbf{自定义名词检索}：把每题「称\ldots 为「X」」里的名词抽出来
  逐个在参考库正文里做字符串检索，揪出了高一17 第 11 题（「数域」＝2008 福建卷（文/理））------
  该题定义与真题实质相同、四个命题全换，LCS 只有 11、\(r\) 也低，判据完全看不见。\textbf{启示：同题
  检测不能只看一处最长连续子串，还得看「数字/字母之外的骨架」与「概念名」这两路。}
\item
  \textbf{局限 1}：参考库只含\textbf{高考真题}，不含模考、竞赛与教材习题，所以「检索不到」只能说明
  「不是高考真题」，\textbf{不能据此推断为自编}（见「非高考但能考据到来源的题」一节：闭集合、
  性质 \(P\)、好集合、好子集、伙伴关系集合等都能在教辅/题库里找到同题）。
\item
  \textbf{局限 2}：同模板换数字的改编题，机器判定边界模糊；出处表的 \textbf{50 条对应（43 道真题）} 均已逐条核对原文（含数学
  符号的具体取值与选项文字），但仍有判断余地。「跨卷同题」同理------组数随判据松紧而变，
  本稿把两档分开列（逐字 48 组 / 同模板变体 20 组），只取可确证的；分数低于变体线、但读过
  原文可确证的 8 对，另列「判据之外手工补记」。
  「存疑」单列的那一条也属此类------判据给不出结论，留给读者自己判断。
\item
  \textbf{短题要补看整段相似度}：一条式选择题的骨架往往不足 25 字，只用「最长公共子串 $\geqslant$ 25」
  会把\textbf{整组}漏掉------\texttt{高一a\ 第\ 5\ 题\ ×\ 高一b\ 第\ 8\ 题}（改写自同一道，相似度 1.000）
  当初就是这么漏的。故判据里并上「整段相似度 $\geqslant$ 0.95」，变体档则用「$\geqslant$ 0.71 且最长串 $\geqslant$ 12
  且公共串含内容字符」。\textbf{每一条路径都写死「最长串 $\geqslant$ 12」}------比对前的 12-gram 预筛丢掉的
  正是「最长串 \textless{} 12」的对，两者口径必须一致，否则会出现「\(r\) 够但这道题根本没进候选池」的漏检
  （2026-09-15 末修的就是这一处：原先「\(r\ge0.95\)」那条没有长度下限，与预筛互相矛盾）。
  0.71 不是随手取的：表内最低那对是「全食／偏食」\(r=0.7129\)，而它往下一直要到 \(r=0.6957\)
  才有下一个候选（一对纯套话），\textbf{0.696--0.713 这一段是空的}，线落在这一档里哪个位置都一样。
\item
  \textbf{阈值压到底是什么样}（记下边界在哪；下面这组数字是 \textbf{12 份}那一轮跑的，17 份重跑后的
  对应数字附在本条末尾）：把判据放宽到「相似度 $\geqslant$ 0.50」，跨卷配对由本稿
  收的二十几对涨到 \textbf{646 组}，但\textbf{最长公共子串}把它切成两半------LCS $\geqslant$ 11 的只有 \textbf{56 对}，
  LCS 9--10 一档就骤增到 \textbf{123 对}，\textbf{这个跳变就是边界}：

  \begin{itemize}
  \tightlist
  \item
    LCS $\geqslant$ 11（56 对）：本稿列的 \textbf{28 对}（逐字 18 ＋ 变体 10）全在里面，已占一半；
    剩下 28 对逐条看过，\textbf{8 对是判据漏掉的真同题、分属 3 个模板}------5 对是上面那个
    「含参二次方程」大族的边（\texttt{高一02\ 第\ 17\ 题} 配 07／08／09／11 那四份，再加
    \texttt{高一05\ 第\ 17\ 题\ ×\ 高一09\ 第\ 19\ 题}；\(r=0.660\sim0.815\)，LCS 全卡在 11），
    另 3 对见上文「判据之外手工补记」（该节共 4 对，第 4 对连 \(r\ge 0.50\) 都没到，不在本档）；其余 \textbf{20 对}的公共串清一色是不含题意的套话
    （「则实数 \(a\) 的取值范围是」之类）；
  \item
    LCS 9--10（123 对）里 \(r\ge 0.80\) 的只剩 \textbf{2 对}：\texttt{高一05\ 第\ 5\ 题\ ×\ 高一08\ 第\ 7\ 题}
    （含参区间包含关系的繁简两版）和 \texttt{高一04\ 第\ 15\ 题\ ×\ 高一09\ 第\ 19\ 题}（即上面那族的
    反向形式）；其余都在 0.80 以下；
  \item
    LCS $\leqslant$ 8：只是共享套话（「已知集合\ldots，则实数 \(a\) 的取值范围是」）。注意\textbf{整段相似度在
    短题上会骗人}------\texttt{高一07\ 第\ 4\ 题\ ×\ 高一b\ 第\ 8\ 题} 的 \(r\) 高达 0.783，LCS 却只有 4，
    是两句都短的缘故，属同题型巧合。
    所以本稿取 LCS $\geqslant$ 12 为线、并给变体档留了整段相似度这一条补充，不再往下压。
    \textbf{17 份重跑后的同一档数字}（2026-09-15，跨卷 45,150 对里取 LCS $\geqslant$ 8 的 936 对）：
    LCS $\geqslant$ 12 共 \textbf{71 对}、LCS $\geqslant$ 11 共 \textbf{91 对}（其中 \(r\ge0.80\) 的 22 对），LCS 9--10 骤增到
    \textbf{346 对}、\(r\ge0.80\) 的\textbf{一对也没有}------\textbf{跳变仍在同一处}，判据不用改。
  \end{itemize}
\end{itemize}


\end{document}
